\documentclass[11pt]{article}

\usepackage{amsmath,amssymb,amsthm}
\usepackage{booktabs}
\usepackage{geometry}
\usepackage{hyperref}
\usepackage{xcolor}

\geometry{margin=2.7cm}
\hypersetup{
  colorlinks=true,
  linkcolor=blue!60!black,
  citecolor=green!45!black,
  urlcolor=blue!60!black,
  pdftitle={Erdos Problem 848: Exact Hall Closure by Square-Sieve Descent and Prefix Colouring},
  pdfauthor={Alex Chengyu Li}
}

\newtheorem{theorem}{Theorem}[section]
\newtheorem{proposition}[theorem]{Proposition}
\newtheorem{lemma}[theorem]{Lemma}
\newtheorem{corollary}[theorem]{Corollary}
\theoremstyle{definition}
\newtheorem{definition}[theorem]{Definition}

\newcommand{\NN}{\mathbb{N}}
\newcommand{\CN}{C_N}
\newcommand{\TN}{T_N}
\newcommand{\calP}{\mathcal{P}}

\title{\textbf{Erd\H{o}s Problem 848:\\
Exact Hall Closure by Square-Sieve Descent and Prefix Colouring}}
\author{Alex Chengyu Li\\[3pt]
\small Independent Researcher, Tokyo, Japan}
\date{July 2026}

\begin{document}
\maketitle

\begin{abstract}
For $N\geq1$, let $A\subseteq[1,N]$ have the property that $ab+1$ is
nonsquarefree for every $a,b\in A$.  We prove that
\[
 |A|\leq \bigl|\{n\leq N:n\equiv7\pmod {25}\}\bigr|,
\]
with equality attained by the progression $7\pmod {25}$.  The proof first
rewrites the assertion as an exact Hall inequality.  An exact prefix
colouring proves the inequality through $1.5\cdot10^6$; two finite
square-sieve arguments cover the interval up to $5\cdot10^6$; and a
valuation--cell--fibre descent covers the four intervals
$[5,10)$, $[10,20)$, $[20,40)$, and $[40,\infty)$ million.  All finite
enumerations used in the proof employ integer or rational arithmetic and are
reproduced by the accompanying source files.
\end{abstract}

\section{Statement and Hall formulation}

Write
\[
  \CN=\{n\in[1,N]:n\equiv 7\pmod {25}\}.
\]
We prove
\begin{equation}\label{eq:target}
  |A|\leq |\CN|
\end{equation}
for every $N$ and every $A\subseteq[1,N]$ such that $ab+1$ is
nonsquarefree for all $a,b\in A$.  Sawhney proved the assertion for all
sufficiently large $N$ by a density sieve and structural reduction
\cite{Sawhney848}; the argument below supplies an explicit all-$N$ proof.

\begin{proposition}[Sharpness]\label{prop:sharpness}
The set $\CN$ is admissible: if $a,b\in\CN$, then $ab+1$ is not squarefree.
\end{proposition}

\begin{proof}
Since $a\equiv b\equiv7\pmod {25}$,
\[
  ab+1\equiv 7^2+1=50\equiv0\pmod {25}.
\]
Thus $5^2\mid ab+1$.
\end{proof}

\section{The direct Hall cut}

Join $b\in[1,N]\setminus\CN$ to $a\in\CN$ when $ab+1$ is squarefree.
For $B\subseteq[1,N]\setminus\CN$, write
\[
  \Gamma_N(B)=\{a\in\CN:\text{$ab+1$ is squarefree for some $b\in B$}\}
\]
and
\[
  \TN(B)=\CN\setminus\Gamma_N(B).
\]
Thus $a\in\TN(B)$ precisely when $ab+1$ is non-squarefree for every
$b\in B$.

For every set $B\subseteq[1,N]\setminus\CN$ such that
$bc+1$ is nonsquarefree for all $b,c\in B$, the Hall inequality to be proved is
\begin{equation}\label{eq:hall-defect}
  |B|+|\TN(B)|\leq |\CN|,
\end{equation}
equivalently $|\Gamma_N(B)|\geq |B|$.

\begin{theorem}[Hall reduction]\label{thm:hall-reduction}
If the Hall inequality \eqref{eq:hall-defect} holds, then the exact bound
\eqref{eq:target} holds for every $N$.
\end{theorem}

\begin{proof}
Let $A$ be admissible and put $B=A\setminus\CN$.  Every subset
$D\subseteq B$ is again a compatible outside set, so
$|\Gamma_N(D)|\geq|D|$ by the Hall inequality \eqref{eq:hall-defect}.  Hall's marriage theorem
gives an injection from $B$ into $\Gamma_N(B)$.  Admissibility gives
\[
  \Gamma_N(B)\cap(A\cap\CN)=\varnothing,
\]
because every vertex of $\Gamma_N(B)$ has a squarefree edge to some element
of $B\subseteq A$.  Hence
\[
  |B|\leq|\Gamma_N(B)|\leq|\CN\setminus A|,
\]
and therefore
\[
  |A|=|A\cap\CN|+|B|\leq|A\cap\CN|+|\CN\setminus A|=|\CN|.
\]
\end{proof}

\begin{theorem}[Hall equivalence]\label{thm:hall-equivalence}
the Hall inequality \eqref{eq:hall-defect} is equivalent to the exact Erd\H{o}s 848 upper bound
\eqref{eq:target}.
\end{theorem}

\begin{proof}
The forward implication is Theorem~\ref{thm:hall-reduction}.  Conversely,
assume \eqref{eq:target} and let $B$ satisfy the hypotheses of
the Hall inequality \eqref{eq:hall-defect}.  Put
\[
  A=B\cup\TN(B).
\]
Products internal to $B$ are non-squarefree by hypothesis.  Products
internal to $\TN(B)$ are divisible by $25$, because all its elements are
$7\pmod {25}$.  Finally, every cross-product between $B$ and $\TN(B)$ is
non-squarefree by the definition of $\TN(B)$.  Hence $A$ is admissible.
Moreover $B$ and $\TN(B)$ are disjoint, so \eqref{eq:target} gives
\[
  |B|+|\TN(B)|=|A|\leq|\CN|,
\]
which is \eqref{eq:hall-defect}.
\end{proof}

Thus the Hall inequality is the original theorem written at the exact
exchange boundary.  All subsequent estimates prove \eqref{eq:hall-defect}.

\section{An exact finite prefix close}

The whole finite remainder admits a single exact colouring certificate.

\begin{theorem}[Two-matching prefix colouring]
\label{thm:finite-prefix-colouring-close}
For every $1\leq N\leq1.5\cdot10^6$, let
\[
 \mathcal D_N=\{x\leq N:x^2+1\text{ is nonsquarefree}\},
\]
and join two distinct members $x,y\in\mathcal D_N$ when $xy+1$ is
nonsquarefree.  This graph has a proper colouring with exactly $|\CN|$
colours.  Consequently \eqref{eq:target} and
the Hall inequality \eqref{eq:hall-defect} hold throughout this interval.
\end{theorem}

\begin{proof}
Every member of an admissible set belongs to $\mathcal D_N$, and every two
distinct members are joined.  Thus it is enough to construct the asserted
proper colouring.

Put $L=1.5\cdot10^6$.  The diagonal set $\mathcal D_L$ is enumerated without
factoring $L$ separate integers.  If $p^2\mid x^2+1$, then $p$ is odd,
$p\equiv1\pmod4$, and $p\leq x\leq L$.  Conversely every solution of
$z^2\equiv-1\pmod {p^2}$ marks an arithmetic progression of diagonal
vertices.  For each prime $p\leq L$, sieve primality, find the two roots
modulo $p$, lift them uniquely modulo $p^2$, verify the lifted congruence,
and mark both progressions.  This is an exact equivalence.

Process the marked vertices in increasing order.  At a given prefix write
\[
 U=\{u:u\equiv7\pmod {25}\},\qquad
 V=\{v:v\equiv18\pmod {25}\},\qquad
 E=\mathcal D_N\setminus(U\cup V).
\]
Maintain two injections into the colour set $U$.  The first is
$\mu:V\to U$ with
\[
 v\mu(v)+1\quad\hbox{squarefree}.
\]
The second is $\nu:E\to U$, and it satisfies
\[
 e\nu(e)+1\quad\hbox{squarefree},\qquad
 e\mu^{-1}(\nu(e))+1\quad\hbox{squarefree}
\]
whenever the second factor exists.  Then for every $u\in U$ the class
\[
 \{u\}\cup\mu^{-1}(u)\cup\nu^{-1}(u)
\]
has at most three vertices, and every product plus one inside it is
squarefree.  These are therefore $|U|=|\CN|$ independent colour classes.

The two injections are maintained prefix by prefix, using no future vertex.
When a new $7$-class vertex appears, add its singleton colour.  When a new
$18$-class vertex appears, either put it in the unique colour not yet used
by $\mu$, or exchange it with one resident $18$-class vertex; both new
products in the exchange are checked.  On the complete interval there are
$39290$ direct insertions and $20710$ one-exchange insertions, and no longer
augmenting path is used.  An exceptional vertex is placed in a colour not
yet used by $\nu$ after checking its one or two required products.  If an
exchange in $\mu$ invalidates a resident exceptional vertex, that vertex is
immediately reinserted.  There are $40692$ such insertions or repairs:
$40688$ use a free colour directly, while only the initial vertices
\[
 41,\quad70,\quad99,\quad117
\]
use augmenting paths, of maximum depth four.  Hence the two injections, and
therefore the proper colouring, exist after every marked vertex is processed.
Between marked vertices the graph is unchanged, while a new colour appears
exactly at a new member of $\CN$.  The certificate consequently covers
every prefix, not merely the upper endpoint.

The terminal exact counts are
\[
\begin{array}{c|r@{\qquad}c|r}
\text{quantity}&\text{value}&\text{quantity}&\text{value}\\ \hline
\text{primes through }L&114155&
\text{lifted roots inspected}&114044\\
|\mathcal D_L|&157733&|U|=|V|&60000\\
|E|&37733&\text{pure edge tests}&125789\\
\text{exceptional edge tests}&88407&
\text{final verification pairs}&135466.
\end{array}
\]
 All $270784$ construction-oracle product queries are factored exactly.
 The terminal verifier then passes all $135466$ class pairs through a fresh
 factorization oracle; in this run their products are distinct.  Deterministic
 $64$-bit Miller--Rabin verifies every factor as prime, their product is
 reconstructed, and repeated factors are rejected; $128$-bit modular
 multiplication prevents overflow.  A final pass checks one diagonal witness
 for every marked vertex, checks that no unmarked vertex was coloured, and
 refactors every pair within every terminal colour class.  Compilation fails
 if assertions are disabled.  The complete producer and verifier are
\path{scripts/systematic_finite_prefix_coloring.cpp}.

Thus every admissible set is a clique in a graph properly coloured with
$|\CN|$ colours, and so has at most $|\CN|$ members.  Hall equivalence,
Theorem~\ref{thm:hall-equivalence}, gives the last assertion.
\end{proof}

\section{An aggregate explicit large-N close}

We now make the known sufficiently-large-$N$ part numerically explicit and
sharpen its threshold.  We retain Sothanaphan's explicit progression estimate
\cite{Sothanaphan848}, but aggregate the $23$ non-base residue classes before
estimating the diagonal Pell tail.  This makes the large-square tail a single
set of integers rather than $23$ separately charged copies.

Let
\[
 C_{2,5}=1-\frac{25}{3\pi^2}<0.155656803,
 \qquad
 C_5=1-\frac{25}{4\pi^2}<0.366742603,
\]
and
\[
 C_{\rm quad}
 =1-\prod_{\substack{p\equiv1\ (4)\\p\geq13}}
       \left(1-\frac2{p^2}\right)
 <0.027346508.
\]
Exact rational multiplication through $10^7$ gives
\[
 1-\prod_{\substack{13\leq p\leq10^7\\p\equiv1\ (4)}}
 \left(1-\frac2{p^2}\right)<0.027346490.
\]
For the omitted factors,
\[
 \begin{aligned}
 1-\prod_{p>10^7}\left(1-\frac2{p^2}\right)
 &\leq2\sum_{p>10^7}\frac1{p^2}\\
 &<0.000000017856.
 \end{aligned}
\]
The last inequality retains the negative boundary term in partial summation.
The calculation is given in Lemma~\ref{lem:boundary-corrected-square-tail}.
Adding the two strict bounds proves the displayed constant.
Fix $q_0\in\{25,100\}$ and $b\leq N$.  For every residue class $t$ such
that $p^2\nmid bt+1$ for every prime $p\mid q_0$, the explicit
off-diagonal progression estimate gives
\begin{equation}\label{eq:explicit-off-diagonal}
 \#\{a\leq N:a\equiv t\pmod {q_0},\ \mu(ab+1)=0\}
 \leq
 \frac N{q_0}\left(1-
   \prod_{(p,q_0b)=1}\left(1-\frac1{p^2}\right)\right)
 +\Phi_{q_0}N^{3/4},
\end{equation}
where $\Phi_{25}=4.7$ and $\Phi_{100}=5.3$.  For completeness, the source's
normalized error function is
\begin{align*}
 E_q(N)={}&\frac{2.367}{N^{3/4}}+\alpha_q+\frac1{N^{3/4}}
 +\frac{N/q+1}{\alpha_qN^{3/2}}
 +\frac1{\alpha_q\sqrt q\,N^{3/2}}\\
 &+\frac{2\lambda q^{1/4}(N/q+1)}{\alpha_qN^{5/4}}
 +\frac{2\lambda}{\alpha_q q^{3/8}N^{3/2}}
 +\frac{\lambda q^{5/4}(N/q+1)}{\alpha_q^2N}
 +\frac{\lambda q^{1/4}}{\alpha_q^2N^2},
\end{align*}
where
\[
 \alpha_{25}=\frac{25}{8},\qquad \alpha_{100}=\frac72,\qquad
 \lambda=\frac{128}{(8\cdot3\cdot5\cdot7\cdot11\cdot13)^{1/4}}<6.876.
\]
It bounds the coefficient of $N^{3/4}$ in
\eqref{eq:explicit-off-diagonal}.  Every nonconstant term decreases with $N$,
and direct substitution at $N=1.6\cdot10^{15}$ gives
\[
 E_{25}(1.6\cdot10^{15})<4.69949,
 \qquad E_{100}(1.6\cdot10^{15})<5.27503.
\]
Thus \eqref{eq:explicit-off-diagonal} with the displayed rounded constants is
valid throughout the range used below, not only at the older published
threshold.

The root-count factor in this estimate can be made much smaller in the present
application.  Let $\rho(m)$ be the maximum number of roots of
$y^2\equiv u\pmod m$ over units $u$, and let $\ell_j$ be the $j$-th odd prime,
$P_k=\prod_{j\leq k}\ell_j$.  Put
\begin{align*}
 Q_*&=8P_{12}=1217001054108840,\\
 N_{100}&=Q_*/20=60850052705442,
 \qquad N_{25}=Q_*/5=243400210821768.
\end{align*}
For $7.67769\cdot10^9\leq N\leq1.6\cdot10^{15}$ define
\[
 H_{25}(N)=\begin{cases}8192,&N<N_{25},\\16384,&N\geq N_{25},\end{cases}
\qquad
 H_{100}(N)=\begin{cases}8192,&N<N_{100},\\16384,&N\geq N_{100}.\end{cases}
\]
For $q_0\in\{25,100\}$, $H=H_{q_0}(N)$, set
\[
 d_H=(2H)^{1/3}N^{2/3}
\]
and define the normalized error envelope
\begin{align}
 \mathcal E_{q_0,H}(N)=\frac1N\bigg(&2.367+d_H+1
 +\frac{(1+2H)(N/q_0+1)}{d_H}
 +\frac{1+2H}{\sqrt{q_0}d_H}\notag\\
 &+\frac{Hq_0N(N/q_0+1)}{d_H^2}+\frac{H}{d_H^2}\bigg).
 \label{eq:primorial-ap-error}
\end{align}

\begin{lemma}[Primorial progression error]
\label{lem:primorial-progression}
In \eqref{eq:explicit-off-diagonal}, throughout
$7.67769\cdot10^9\leq N\leq1.6\cdot10^{15}$, the term
$\Phi_{q_0}N^{3/4}$ may be replaced by
\[
 N\mathcal E_{q_0,H_{q_0}(N)}(N).
\]
On every interval on which the corresponding $H_{q_0}$ is constant, the error
envelope is decreasing.
\end{lemma}

\begin{proof}
In the source reduction the relevant modulus is $d_1=q_0b/g$, where $g$ divides
the radical of $q_0$.  Removing one redundant factor
$5$ when its exponent is at least two shows that there is an $m\leq q_0N/5$
with $\rho(m)=\rho(d_1)$.  If $m=2^\nu m_0$ with $m_0$ odd, then
\[
 \rho(m)\leq c_\nu2^{\omega(m_0)},\qquad
 c_0=c_1=1,\quad c_2=2,\quad c_\nu=4\ (\nu\geq3).
\]
The least modulus with root count greater than $8192$ is $8P_{12}=Q_*$.
The least one with root count greater than $16384$ is
$8P_{13}=52331045326680120$, which exceeds $20(1.6\cdot10^{15})$.
This proves the two step-function bounds for $\rho(d_1)$.

Repeat the M\"obius split underlying \eqref{eq:explicit-off-diagonal}, but keep
$\rho(d_1)\leq H$ instead of replacing it by $6.876d_1^{1/4}$.  With
$X\leq N/q_0+1$, $\sqrt{q_0}\leq d_1\leq q_0N$, and
$D=\lceil d_H\rceil$, the total additive error is at most
\[
 2.367+D+\frac{(1+2H)X}{D}
 +\frac{1+2H}{\sqrt{q_0}D}
 +\frac{Hq_0NX}{D^2}+\frac H{D^2}.
\]
Use $d_H\leq D\leq d_H+1$, increasing the positive $D$ term and decreasing
all denominators, to obtain \eqref{eq:primorial-ap-error}.  Every normalized
term is decreasing while $H$ is fixed.
\end{proof}

For the diagonal part put
\[
 A_7=\{a\in A:a\equiv7\pmod {25}\},\quad
 A_{18}=\{a\in A:a\equiv18\pmod {25}\},\quad
 A^*=A\setminus(A_7\cup A_{18}),
\]
and define
\[
 U_{25}=\{x\leq N:x\not\equiv7,18\pmod {25}\},\qquad
 U_{50}=\{x\in U_{25}:x\text{ is odd}\}.
\]
For $r\in\{0,1,2,3\}$ also put
\[
 U_{100,r}=\{x\in U_{25}:x\equiv r\pmod4\};
\]
this is a union of $23$ classes modulo $100$.
Two coprimality moduli will be used.  Put
\[
\begin{array}{c|c|c|c}
i&M_i&\eta_i=\varphi(M_i)/M_i&
h_i>\displaystyle\sum_{\substack{1\leq a\leq M_i\\(a,M_i)=1}}1/a\\ \hline
0&2310&16/77&2.189\\
1&43890&288/1463&2.685
\end{array}
\]
and $e_0=64$, $e_1=128$.  For $R\geq M_i^2$, put
\[
 K_i(R)=2\eta_i R\left(h_i+\eta_i\left(1+\log\frac{\sqrt R}{M_i}\right)\right)
      +e_i\sqrt R.
\]
For $q\in\{25,50,100\}$ write $\theta_q=23/q$, put
\[
 r_N=\frac{12}{5}N^{2/3},
\]
and set, whenever $r_N\geq M_i^2$,
\begin{align}
 \Delta_q^{(i)}(N)={}&\frac{23}{N}+\frac{3K_i(r_N+1)}N
 +\frac{46(\log(r_N+1)+2)}{q r_N}\notag\\
 &+\frac{N^2+1}{N r_N^2}
       \left(1+\log_{2+\sqrt3}N\right).
 \label{eq:aggregate-diagonal-error}
\end{align}

\begin{lemma}[Aggregate diagonal estimate]
\label{lem:aggregate-diagonal}
For $q=25$ and $U_q=U_{25}$, for $q=50$ and $U_q=U_{50}$, or for
$q=100$ and $U_q=U_{100,r}$ with any fixed $r\pmod4$,
\begin{equation}\label{eq:aggregate-diagonal}
 \#\{x\in U_q:\mu(x^2+1)=0\}
 \leq N\bigl(\theta_q C_{\rm quad}+\Delta_q^{(i)}(N)\bigr)
 \qquad(i=0,1),
\end{equation}
provided $r_N\geq M_i^2$.  The $i=0$ envelopes decrease for
  $N\geq7.8\cdot10^{10}$, and at that endpoint
 \[
  \Delta_{25}^{(0)}(N)<0.002663263,
  \quad \Delta_{50}^{(0)}(N)<0.002662851,
  \quad \Delta_{100}^{(0)}(N)<0.002662646.
\]
The $i=1$ envelopes decrease for $N\geq3\cdot10^{13}$, and at that endpoint
\[
 \Delta_{25}^{(1)}(N)<0.000402450,
 \qquad \Delta_{50}^{(1)}(N)<0.000402441.
\]
\end{lemma}

\begin{proof}
Let $R=\lceil r_N\rceil$.  In the M\"obius expansion of the indicator of
$x^2+1$ being squarefree, only squarefree $s$ composed of primes
$p\equiv1\pmod4$, $p\geq13$, contribute.  For such an $s$, the congruence
$x^2\equiv-1\pmod {s^2}$ has exactly $2^{\omega(s)}$ roots.  Counting complete
blocks in the $23$ allowed residue classes gives
\[
 \left|A_s-\theta_qN\frac{2^{\omega(s)}}{s^2}\right|
 \leq3\,2^{\omega(s)},
\]
where $A_s$ counts $x\in U_q$ with $s^2\mid x^2+1$.

Every such $s$ is coprime to both $M_i$.  If
$C_{M_i}(y)=\#\{n\leq y:(n,M_i)=1\}$, inclusion--exclusion gives
$C_{M_i}(y)\leq\eta_i y+2^{\omega(M_i)}$.  The hyperbola decomposition
therefore gives
\begin{align*}
 \sum_{\substack{s\leq R\\(s,M_i)=1}}d(s)
 &\leq2\sum_{\substack{a\leq\sqrt R\\(a,M_i)=1}}C_{M_i}(R/a)\\
 &\leq2\eta_i R\sum_{\substack{a\leq\sqrt R\\(a,M_i)=1}}\frac1a
      +e_i\sqrt R
 \leq K_i(R).
\end{align*}
The two first-block harmonic bounds are finite rational computations.  The
last inequality follows by separating the first block modulo $M_i$ and
bounding each later block by $\eta_i/k$.  Hence the small-$s$ rounding loss is
at most $3K_i(R)$, while
\[
 \sum_{s>R}\frac{d(s)}{s^2}\leq\frac{2(\log R+2)}R.
\]
The remaining large-square tail is charged only once.  Writing
$x^2+1=ds^2$, the negative Pell equation has at most
$1+\log_{2+\sqrt3}N$ positive solutions for each fixed $d$, so
\[
 \sum_{s>R}A_s\leq\frac{N^2+1}{R^2}
 \left(1+\log_{2+\sqrt3}N\right).
\]
Finally $|U_q|\leq\theta_qN+23$.  Combining the bounds and using
$r_N\leq R\leq r_N+1$ proves \eqref{eq:aggregate-diagonal}.  Notice that
 $r_N\geq M_0^2$ already for $N\geq7.8\cdot10^{10}$, whereas
$r_N\geq M_1^2$ for $N\geq2.274\cdot10^{13}$.  Elementary differentiation
proves the asserted monotonicity on the displayed ranges; direct endpoint
evaluation gives the constants.
\end{proof}

For the last refinement define
\[
 V_8=\{x\leq N:x\not\equiv7,18\pmod {25},\quad
                  (x\text{ odd or }8\mid x)\}
\]
and, whenever $r_N\geq M_i^2$,
\begin{align}
 \Delta_8^{(i)}(N)={}&\frac{115}{N}+\frac{6K_i(r_N+1)}N
 +\frac{23(\log(r_N+1)+2)}{20r_N}\notag\\
 &+\frac{N^2+1}{Nr_N^2}\left(1+\log_{2+\sqrt3}N\right).
 \label{eq:eight-divisible-diagonal-error}
\end{align}

\begin{lemma}[Two-adic split of the even non-base branch]
\label{lem:two-adic-even-branch}
Let $3\cdot10^{13}\leq N\leq6\cdot10^{13}$ and let $b$ be even.  In the
two $q_0=25$ progression estimates obtained from $b$, one may use
\[
 H=2048\quad\text{if }\nu_2(b)=1,\qquad
 H=4096\quad\text{if }\nu_2(b)=2,\qquad
 H=8192\quad\text{if }\nu_2(b)\geq3.
\]
Moreover
\[
 \#\{x\in V_8:\mu(x^2+1)=0\}
 \leq N\left(\frac{23}{40}C_{\rm quad}+\Delta_8^{(1)}(N)\right),
\]
and at the left endpoint $\Delta_8^{(1)}(N)<0.000667629$.
\end{lemma}

\begin{proof}
After removal of the redundant factor $5$, the effective root-count modulus is
$5b$.  On the stated interval it is at most $3\cdot10^{14}$.  If
$\nu_2(b)=1$, a root count greater than $2048$ would require a modulus at least
$2P_{12}=304250263527210$; if $\nu_2(b)=2$, a count greater than $4096$ would
require at least $4P_{12}=608500527054420$.  If $8\mid b$, the existing
$8192$ bound applies.  This proves the progression assertions.

For the diagonal assertion, the odd non-base classes have density $23/50$ and
the classes divisible by $8$ have density $23/200$.  For each root modulo
$s^2$, counting the two parts successively modulo $2,8$, and $25$ gives total
discrepancy at most $6$.  Their union has $115$ classes modulo $200$, so the
proof of Lemma~\ref{lem:aggregate-diagonal}, with small-modulus loss $6K_1(R)$
and the Pell tail still paid once, gives
 \eqref{eq:eight-divisible-diagonal-error} with $i=1$.  Its terms decrease, and direct
substitution gives the stated endpoint value.
\end{proof}

\begin{corollary}[Low-range two-adic diagonal]
\label{cor:low-two-adic-diagonal}
For $N\geq7.8\cdot10^{10}$,
\[
 \#\{x\in V_8:\mu(x^2+1)=0\}
 \leq N\left(\frac{23}{40}C_{\rm quad}+\Delta_8^{(0)}(N)\right).
\]
The envelope decreases, and at the left endpoint
\[
 \Delta_8^{(0)}(N)<0.004510951.
\]
\end{corollary}

\begin{proof}
Repeat the diagonal part of Lemma~\ref{lem:two-adic-even-branch} with $M_0$
in place of $M_1$.  Here $r_N\geq M_0^2$ throughout the stated range, so the
same $115$ residue-class endpoint loss, $6K_0(R)$ small-square loss, and single
Pell tail give \eqref{eq:eight-divisible-diagonal-error} with $i=0$.
Elementary differentiation shows that every term decreases; direct
substitution at $N=7.8\cdot10^{10}$ gives the displayed value.
\end{proof}

For $\varepsilon\in\{0,1,2,3\}$ put
\[
 W_\varepsilon=\{x\in U_{25}:x\not\equiv\varepsilon\pmod4\}
\]
and define
\begin{align}
 \Delta_W(N)={}&\frac{69}{N}+\frac{6K_0(r_N+1)}N
 +\frac{69(\log(r_N+1)+2)}{50r_N}\notag\\
 &+\frac{N^2+1}{Nr_N^2}
       \left(1+\log_{2+\sqrt3}N\right).
 \label{eq:three-mod-four-diagonal-error}
\end{align}

\begin{lemma}[Three-mod-$4$-class diagonal]
\label{lem:three-mod-four-diagonal}
For $N\geq7.8\cdot10^{10}$ and every choice of $\varepsilon$,
\[
 \#\{x\in W_\varepsilon:\mu(x^2+1)=0\}
 \leq N\left(\frac{69}{100}C_{\rm quad}+\Delta_W(N)\right).
\]
The envelope decreases, and at the left endpoint
\[
 \Delta_W(N)<0.004511053.
\]
\end{lemma}

\begin{proof}
For every squarefree small-square modulus $s$ in the proof of
Lemma~\ref{lem:aggregate-diagonal}, its root count on $W_\varepsilon$ is the
difference of its root counts on $U_{25}$ and $U_{100,\varepsilon}$.
Subtracting their main terms gives density $69/100$, while the two
discrepancies give at most $6\,2^{\omega(s)}$.  Also $W_\varepsilon$ is a union
of $69$ classes modulo $100$.  Thus the same hyperbola calculation has endpoint
loss $69$, small-square loss $6K_0(R)$, harmonic-tail coefficient $69/50$, and
one common Pell tail.  These are exactly the four terms in
\eqref{eq:three-mod-four-diagonal-error}.  Their monotonicity and the displayed
endpoint value follow by differentiation and direct substitution.
\end{proof}

The final low-range refinement uses the stronger coprimality modulus before a
complete residue block is available.  Put
\[
 M_*=43890,\qquad \eta_*=\frac{288}{1463},\qquad
 h_*:=\sum_{\substack{1\leq a\leq7194\\(a,M_*)=1}}\frac1a<2.329039
\]
and
\[
 K_*(R)=2\eta_*Rh_*+128\sqrt R.
\]
The displayed bound for $h_*$ is an exact finite rational summation.  On
$6.7\cdot10^{10}\leq N\leq1.001\cdot10^{11}$ define
\begin{align*}
 \widetilde\Delta_q(N)&=\Delta_q^{(0)}(N)
       -\frac{3K_0(r_N+1)}N+\frac{3K_*(r_N+1)}N
       \qquad(q=25,50,100),\\
 \widetilde\Delta_W(N)&=\Delta_W(N)
       -\frac{6K_0(r_N+1)}N+\frac{6K_*(r_N+1)}N,\\
 \widetilde\Delta_8(N)&=\Delta_8^{(0)}(N)
       -\frac{6K_0(r_N+1)}N+\frac{6K_*(r_N+1)}N.
\end{align*}

\begin{lemma}[Short-block diagonal refinement]
\label{lem:short-block-diagonal}
Throughout this interval the conclusions of Lemma~\ref{lem:aggregate-diagonal},
Corollary~\ref{cor:low-two-adic-diagonal}, and
Lemma~\ref{lem:three-mod-four-diagonal} hold with their respective error terms
replaced by $\widetilde\Delta_q$, $\widetilde\Delta_8$, and
$\widetilde\Delta_W$.  These envelopes decrease, and at
$N=6.7\cdot10^{10}$,
\[
 \begin{array}{lll}
 \widetilde\Delta_{25}<0.002514345,&
 \widetilde\Delta_{50}<0.002513892,&
 \widetilde\Delta_{100}<0.002513665,\\
 \widetilde\Delta_W<0.004175714,&
 \widetilde\Delta_8<0.004175602.
 \end{array}
\]
\end{lemma}

\begin{proof}
For $R=r_N+1$ on the stated interval, direct endpoint evaluation gives
$\sqrt R<7194$.  Inclusion--exclusion gives
\[
 C_{M_*}(y)=\#\{n\leq y:(n,M_*)=1\}
 \leq\eta_*y+64.
\]
The hyperbola step in Lemma~\ref{lem:aggregate-diagonal} therefore yields,
without requiring $R\geq M_*^2$,
\[
 \sum_{\substack{s\leq R\\(s,M_*)=1}}d(s)
 \leq2\eta_*R
   \sum_{\substack{a\leq\sqrt R\\(a,M_*)=1}}\frac1a+128\sqrt R
 \leq K_*(R).
\]
Every squarefree modulus $s$ occurring in the three diagonal arguments is
coprime to $M_*$, because its prime factors are $1\pmod4$, whereas the added
prime is $19$.  Thus only the displayed small-square rounding term changes;
all main terms, harmonic tails, and the single Pell tail remain unchanged.
This gives the three tilded formulas.  Each normalized term decreases by
elementary differentiation, and direct substitution gives the endpoint
bounds.
\end{proof}

\begin{theorem}[Explicit large range]\label{thm:explicit-large-close}
The exact bound \eqref{eq:target}, and equivalently
the Hall inequality \eqref{eq:hall-defect}, holds for every $N\geq3\cdot10^{13}$.
\end{theorem}

\begin{proof}
If $A$ is contained in $A_7$ or $A_{18}$, the result is immediate.  Otherwise
the exhaustive parity decomposition gives the following four normalized upper
bounds, in order: $A^*$ has an even element; $A^*$ is odd and the base part is
odd; $A^*$ is odd and a base class has an even element; and $A^*=\varnothing$
with both base classes nonempty:
\begin{align*}
 \frac{23}{25}C_{\rm quad}+\frac2{25}C_{2,5}&<0.0376115,\\
 \frac{23}{50}C_{\rm quad}+\frac1{50}+\frac1{50}C_{2,5}&<0.0356927,\\
 \frac{23}{50}C_{\rm quad}+\frac1{25}C_5+\frac1{25}C_{2,5}&<0.0334755,\\
 \frac2{25}C_5&<0.0293395.
\end{align*}
First let $3\cdot10^{13}\leq N<6\cdot10^{13}$.  In the last three cases,
using $H=8192$ in the progression envelope gives at the left endpoint the
respective bounds
\[
 0.038548,\qquad0.036331,\qquad0.031792,
\]
and all three decrease thereafter.  In the first case choose an even
$b\in A^*$.  If $\nu_2(b)=1$, Lemma~\ref{lem:two-adic-even-branch} with the
unrestricted diagonal estimate gives an endpoint bound below $0.039559$; if
$\nu_2(b)=2$, it gives a bound below $0.039961$.  If neither type of element
exists, every even element of $A^*$ is divisible by $8$, so $A^*\subseteq V_8$.
The restricted diagonal estimate and $H=8192$ then give
\[
 \frac{23}{40}C_{\rm quad}+\frac2{25}C_{2,5}
  +\Delta_8^{(1)}(N)+2\mathcal E_{25,8192}(N)<0.031297
\]
at the left endpoint, again with decreasing right side.  Thus the entire
interval $3\cdot10^{13}\leq N<6\cdot10^{13}$ is closed.

It remains to prove the result for $N\geq6\cdot10^{13}$.
Indeed, in the first case an even member of $A^*$ sieves both base classes.
In the second, the four odd base progressions modulo $100$ split into two on
which $4\mid ab+1$ automatically and two sieved progressions.  In the third,
an even base member sieves the opposite class and an odd member of $A^*$
sieves its own base class.  In the fourth, one member of each base class
sieves the other.  Lemmas~\ref{lem:primorial-progression} and
\ref{lem:aggregate-diagonal} therefore give
\begin{align*}
 U_1(N)&=0.0376115+\Delta_{25}^{(1)}(N)+2\mathcal E_{25,H_{25}(N)}(N),\\
 U_2(N)&=0.0356927+\Delta_{50}^{(1)}(N)+2\mathcal E_{100,H_{100}(N)}(N)+2/N,\\
 U_3(N)&=0.0334755+\Delta_{50}^{(1)}(N)+2\mathcal E_{25,H_{25}(N)}(N),\\
 U_4(N)&=0.0293395+2\mathcal E_{25,H_{25}(N)}(N).
\end{align*}
The only increases can occur at $N_{100}$ and $N_{25}$.  Evaluating the
right-hand envelopes at the initial point and immediately after each jump
gives
\[
\begin{array}{c|cccc}
N & U_1&U_2&U_3&U_4\\ \hline
6\cdot10^{13}&0.0399044&0.0379855&0.0357684&0.0312858\\
N_{100}&0.0398938&0.0384785&0.0357578&0.0312767\\
N_{25}&0.0393740&0.0374551&0.0352380&0.0308770
\end{array}
\]
All entries are strictly below $1/25$.  Between rows they decrease.
For $N\geq1.6\cdot10^{15}$, use the original rounded progression errors
$4.7N^{-1/4}$ and $5.3N^{-1/4}$ together with the decreasing aggregate
diagonal bound; at the endpoint the four bounds are respectively below
\[
 0.039224,\qquad0.037495,\qquad0.035088,\qquad0.030827.
\]
Thus every case is closed.  The Hall statement follows from
Theorem~\ref{thm:hall-equivalence}.
\end{proof}

\begin{corollary}[Parity pruning below the large close]
\label{cor:even-nonbase-residual}
Let $8.858\cdot10^{12}\leq N<3\cdot10^{13}$.  If a compatible outside set $B$ violates
the Hall inequality, then $B\setminus C_N^{18}$ contains an even element.
\end{corollary}

\begin{proof}
Put $A=B\cup T_N(B)$ as in Theorem~\ref{thm:hall-equivalence}; then
$A^*=B\setminus C_N^{18}$.  If $A^*=\varnothing$, the pure opposite-base
matching, Corollary~\ref{cor:opposite-base-matching}, closes the Hall
inequality.  Suppose $A^*$ is nonempty and odd.  If the base part
$A_7\cup A_{18}$ is odd, use the low-range diagonal envelope
\[
 0.0356927+\Delta_{50}^{(0)}(N)
 +2\mathcal E_{100,8192}(N)+2/N<0.039999984
\]
at $N=8.858\cdot10^{12}$; it remains decreasing throughout the stated interval.  If
the base part has an even element, the corresponding bound
\[
 0.0334755+\Delta_{50}^{(0)}(N)
 +2\mathcal E_{25,8192}(N)<0.037783
\]
holds at $N=8.858\cdot10^{12}$ and also decreases.  Either alternative
would give $|A|<N/25$, contradicting the assumed Hall defect.  Hence $A^*$,
and therefore $B\setminus C_N^{18}$, contains an even element.
\end{proof}

\begin{corollary}[Quantitative even-pivot pruning]
\label{cor:dense-even-nonbase-residual}
Let $8.858\cdot10^{12}\leq N<3\cdot10^{13}$, and suppose that a compatible outside set
$B$ violates the Hall inequality.  If
\[
 E=\{b\in B\setminus C_N^{18}:2\mid b\},\qquad e=|E|,
\]
Define
\[
 U_{\rm odd}(N)=0.0356927+\Delta_{50}^{(0)}(N)
 +2\mathcal E_{100,8192}(N)+2/N.
\]
Then
\[
 e>\left(\frac1{25}-U_{\rm odd}(N)\right)N-\frac7{25}>2.
\]
Thus there are always two distinct even elements of
$B\setminus C_N^{18}$.  If $N\geq10^{13}$, the stronger bound
\begin{equation}\label{eq:even-pivot-density}
 e>\frac{21N}{125000}-\frac7{25}.
\end{equation}
holds, and there are distinct even $b,c\in B\setminus C_N^{18}$ with
\[
 1\leq |b-c|\leq5952.
\]
For this pair the prime $2$ cannot occur in the same-prime part and
\begin{equation}\label{eq:even-pair-same-prime}
 |S_N(b,c)|\leq\frac{58N}{11025}+2.
\end{equation}
The original Hall inequality follows if
\begin{equation}\label{eq:even-pair-distinct-debit}
 |T_N(B)\cap D_N(b,c)|
 \leq\frac{383N}{11025}-|B|-\frac{57}{25}.
\end{equation}
If moreover $2.25\cdot10^{13}\leq N<3\cdot10^{13}$, the pair may be chosen
with $|b-c|\leq880$.  Then at most one same-prime witness occurs,
\begin{equation}\label{eq:upper-even-pair-same-prime}
 |S_N(b,c)|\leq\frac N{225}+1,
\end{equation}
and it is sufficient that
\begin{equation}\label{eq:upper-even-pair-distinct-debit}
 |T_N(B)\cap D_N(b,c)|
 \leq\frac{8N}{225}-|B|-\frac{32}{25}.
\end{equation}
Thus every residual defect in this range is unconditionally excluded from the
one-even-pivot and $2$-adic same-prime branches.
\end{corollary}

\begin{proof}
Put $A=B\cup T_N(B)$ and $A^*=B\setminus C_N^{18}$, and remove the set $E$
from $A$.  The remaining non-base part is odd.  If it is nonempty, the larger
of the two low-range odd envelopes in the preceding proof is at most
$U_{\rm odd}(N)N$.  If it is empty, retain any $b\in E$, whose existence follows
from Corollary~\ref{cor:even-nonbase-residual}.  The two off-diagonal
progression estimates from $b$ give instead
\[
 |A\setminus E|\leq N\left(
 \frac2{25}C_{2,5}+2\mathcal E_{25,8192}(N)\right)
 <0.016136N
\]
at $N=8.858\cdot10^{12}$, with a decreasing right side.  Hence in all cases
$|A\setminus E|\leq U_{\rm odd}(N)N$.  A Hall defect has
$|A|>|C_N|\geq N/25-7/25$, so subtraction proves the first displayed bound.
At the left endpoint $U_{\rm odd}(N)<0.039999984$, so its right side exceeds
$1.6\cdot10^{-8}N-7/25>2$.  At $N=10^{13}$ one has
$U_{\rm odd}(N)<0.039832$, and monotonicity proves
\eqref{eq:even-pivot-density} on the remaining range.

List the $e$ even elements increasingly.  Since
\[
 \frac{21}{125000}-\frac1{5953}=\frac{13}{744125000}
\]
and $N\geq10^{13}$, equation~\eqref{eq:even-pivot-density} gives
$e-1>(N-1)/5953$.  A least consecutive gap is therefore less than $5953$;
being even, it is at most $5952$.

For an even pivot, $ab+1$ is odd, so $2$ cannot witness a square divisor.
Every same-prime witness also differs from $5$ and has $p^2\mid|b-c|$.
There are at most two such primes, because
$(3\cdot7\cdot11)^2=53361>5952$.  The two largest residue-class budgets are
therefore those for $p=3,7$, and
\[
 |S_N(b,c)|\leq\frac N{25}\left(\frac19+\frac1{49}\right)+2
 =\frac{58N}{11025}+2.
\]
Combining this with $|C_N|\geq N/25-7/25$ proves that
\eqref{eq:even-pair-distinct-debit} is sufficient for the original Hall
inequality, exactly as in Proposition~\ref{prop:bounded-gap-residual}.

At $N=2.25\cdot10^{13}$, direct substitution in the decreasing low-range odd
envelope gives
\[
 0.04-\left(0.0356927+\Delta_{50}^{(0)}(N)
 +2\mathcal E_{100,8192}(N)+2/N\right)
 >\frac1{881}+4\cdot10^{-6}.
\]
The preceding argument therefore gives $e-1>(N-1)/881$ throughout the final
subinterval.  The least even gap is less than $881$, hence at most $880$.
Two distinct same-prime witnesses would force
$(3\cdot7)^2=441$ to divide this even gap, and consequently would force the
gap to be at least $882$, a contradiction.  Thus only $p=3$ need be charged,
which proves \eqref{eq:upper-even-pair-same-prime}.  Subtracting its bound from
$|C_N|\geq N/25-7/25$ gives
\eqref{eq:upper-even-pair-distinct-debit}.
\end{proof}

\begin{lemma}[Finite two-pivot sieve on the base classes]
\label{lem:finite-two-pivot-base-sieve}
Let $8.858\cdot10^{12}\leq N<3\cdot10^{13}$, and let $b,c\leq N$ be
distinct even integers, neither congruent to $7$ or $18\pmod {25}$.  Then
\[
 \#\{a\leq N:a\equiv7\text{ or }18\pmod {25},\
       \ \mu(ab+1)=\mu(ac+1)=0\}<0.013992N.
\]
If in addition $N\geq10^{13}$ and $|b-c|\leq5952$, the right side improves
to $0.012132N$.
\end{lemma}

\begin{proof}
Let
\[
 \mathcal P=\{3,7,11,13,17,19,23,29,31,37,41,43,47\}.
\]
For $p\in\mathcal P$, each of $p^2\mid ab+1$ and $p^2\mid ac+1$
specifies at most one residue of $a$ modulo $p^2$.  When both residues exist,
they coincide exactly when $p^2\mid b-c$.

Fix one of the two classes $a\equiv7,18\pmod {25}$.  Inclusion--exclusion for
the event that some $p\in\mathcal P$ divides $ab+1$ to second order, the
analogous event for $ac+1$, and their union gives the local intersection
density
\[
 I_J=1-2\prod_{p\in\mathcal P}\left(1-\frac1{p^2}\right)
 +\prod_{p\in J}\left(1-\frac1{p^2}\right)
  \prod_{p\in\mathcal P\setminus J}\left(1-\frac2{p^2}\right),
\]
where $J=\{p\in\mathcal P:p^2\mid b-c\}$.  In general it is largest when
$J=\mathcal P$, and direct rational multiplication gives
\begin{equation}\label{eq:small-two-pivot-density}
 I_J\leq I_{\mathcal P}
 =1-\prod_{p\in\mathcal P}\left(1-\frac1{p^2}\right)<0.152377.
\end{equation}
If $|b-c|\leq5952$, at most two primes lie in $J$, because
$(3\cdot7\cdot11)^2>5952$.  In that case the worst choice is
$J=\{3,7\}$ and $I_J<0.129757$.
Every nonempty inclusion--exclusion intersection is one CRT class after
$a\pmod {25}$ is fixed.  The total rounding discrepancy is at most
\[
 \epsilon=2(2^{13}-1)+(3^{13}-1)
 =1610704.
\]
Thus the small-prime common failures in both base classes contribute at most
\begin{equation}\label{eq:small-two-pivot-count}
 \frac{2I_J}{25}N+2\epsilon.
\end{equation}

It remains to charge a witness prime larger than $47$ in at least one of the
two linear forms.  Put $Y=\lfloor N/333\rfloor$.  The finite computation
\[
 \sum_{47<p\leq1000}\frac1{p^2}
 +\sum_{n>1000}\frac1{n^2}<0.004751=:\tau
\]
gives $\sum_{p>47}p^{-2}<\tau$.  For $47<p\leq Y$, a fixed prime contributes
at most $N/(25p^2)+1$ in a fixed base class.  Hence one form in one base
class contributes at most
\[
 \frac{\tau N}{25}+\pi(Y)
\]
in this intermediate range.

For $p>Y$, use the tail calculation in Proposition~1 of
Sothanaphan~\cite{Sothanaphan848}.  After the same removal of one redundant
factor $5$ as in Lemma~\ref{lem:primorial-progression}, the root count is at
most $8192$ because $5N<Q_*$.  If $X\leq N/25+1$ is the length of the base
progression and $d_1$ is its transformed modulus, then
\[
 \sum_{r>Y}A_r
 \leq8192\left(\frac{2x}{d_1Y}+\frac{x}{Y^2}\right)
 \leq L(N),
\]
where
\[
 L(N)=8192\left(
 \frac{2(N/25+2)}Y+\frac{N^2+25N+1}{Y^2}\right).
\]
This bounds, in particular, the prime-square tail.  There are two forms and
two base classes, so \eqref{eq:small-two-pivot-count} and the union bound give
\begin{equation}\label{eq:finite-two-pivot-envelope}
 \frac{|G_N(b,c)|}{N}\leq
 \frac{2I_J}{25}+\frac{2\epsilon}{N}
 +\frac{4\tau}{25}+\frac{4\pi(Y)}N+\frac{4L(N)}N,
\end{equation}
where $G_N(b,c)$ denotes the set in the statement.

Rosser and Schoenfeld's explicit estimate
$\pi(y)<1.25506y/\log y$ for $y>1$~\cite{RosserSchoenfeld} applies.  Since
$N/334\leq Y\leq N/333$, all nonconstant terms in the resulting upper
envelope decrease for $N\geq8.858\cdot10^{12}$.  At the left endpoint its
five terms, using $I_J<0.152377$, are respectively less than
\[
 0.012190160,\quad0.000000364,\quad0.000760160,\quad
 0.000628127,\quad0.000412773.
\]
Their sum is less than $0.013992$.  Under the additional gap hypothesis,
evaluate instead at $N=10^{13}$ with $I_J<0.129757$; the five terms
are respectively less than
\[
 0.01038056,\quad0.000000323,\quad0.000760160,\quad
 0.000624970,\quad0.000365635.
\]
and their sum is less than $0.012132$.  This proves both assertions.
\end{proof}

\begin{theorem}[Improved explicit large range]
\label{thm:8858-large-close}
The exact bound \eqref{eq:target}, and equivalently the Hall inequality \eqref{eq:hall-defect},
holds for every $N\geq8.858\cdot10^{12}$.
\end{theorem}

\begin{proof}
Theorem~\ref{thm:explicit-large-close} handles $N\geq3\cdot10^{13}$.
Suppose instead that $8.858\cdot10^{12}\leq N<3\cdot10^{13}$ and that $B$ is a Hall
defect.  Put $A=B\cup T_N(B)$.  Corollary~
\ref{cor:dense-even-nonbase-residual} supplies distinct even
$b,c\in B\setminus C_N^{18}$.

Every $x\in A^*$ has $\mu(x^2+1)=0$, so the repaired low-range aggregate
diagonal estimate gives
\[
 |A^*|\leq N\left(\frac{23}{25}C_{\rm quad}
                         +\Delta_{25}^{(0)}(N)\right)<0.025784N.
\]
Every member of $A_7\cup A_{18}$ is a common non-neighbour of $b,c$.
Lemma~\ref{lem:finite-two-pivot-base-sieve} therefore gives
\[
 |A|< (0.025784+0.013992)N=0.039776N
 <\frac N{25}-\frac7{25}\leq|C_N|,
\]
contradicting the Hall defect.  Thus the middle interval is closed, and the
Hall equivalence finishes the proof.
\end{proof}

\section{A direct singleton Hall close}

The full large-$N$ argument spends density on all parts of an admissible set.
For a singleton Hall defect we only need one squarefree neighbour, and the
same explicit progression estimate becomes effective much earlier.

We use the following all-parameter form of the M\"obius splitting estimate
proved in Proposition~1 of Sothanaphan~\cite{Sothanaphan848}.  If
$x\geq q\geq1$, $(u,q)=1$, and $D\geq1$ is an integer, then
\begin{equation}\label{eq:explicit-squarefree-lower}
 S_2(x;q,u)\geq
 \frac{x}{q}\prod_{p\nmid q}\left(1-\frac1{p^2}\right)
 -D-\frac{x}{qD}
 -\lambda q^{1/4}\left(\frac{2x}{qD}+\frac{x}{D^2}\right),
\end{equation}
where $S_2(x;q,u)$ counts squarefree positive integers at most $x$ in the
class $u\pmod q$ and
\[
 \lambda=\frac{128}{(8\cdot3\cdot5\cdot7\cdot11\cdot13)^{1/4}}<6.876.
\]
For $N\geq2\cdot10^9$ define
\begin{align}
 \mathcal L(N)={}&0.6079\left(\frac N{25}-2\right)
 -(3.2N^{3/4}+1)
 -\frac{N/25+1}{3.2N^{3/4}}\notag\\
 &-6.876\sqrt5N^{1/4}
 \left(
   \frac{2(N/25+1)}{3.2N^{3/4}}
   +\frac{N^2+25N}{(3.2N^{3/4})^2}
 \right).
 \label{eq:uniform-pivot-expansion}
\end{align}

\begin{theorem}[Uniform one-pivot Hall expansion]\label{thm:singleton-large-close}
Let $N\geq2\cdot10^9$, and let $b\in[1,N]\setminus\CN$.  Then there is an
$a\in\CN$ for which $ab+1$ is squarefree.  More precisely,
\[
 |\Gamma_N(\{b\})|\geq\mathcal L(N)>0.00207N,
\]
and $\mathcal L(2\cdot10^9)>4.14\cdot10^6$.  Consequently every nonempty
compatible outside set $B$ with $|B|\leq\mathcal L(N)$ satisfies the Hall
inequality~\eqref{eq:hall-defect}.  This includes all singleton defects.
\end{theorem}

\begin{proof}
Write $a=7+25n$ and put
\[
 X=|\CN|,\qquad g=(7b+1,25b),\qquad
 q=\frac{25b}{g},\qquad c=\frac{7b+1}{g}.
\]
Since $(b,7b+1)=1$, one has $g\mid25$.  The exclusion
$b\not\equiv7\pmod {25}$ gives $g\in\{1,5\}$, $(c,q)=1$, and
$g\mid q$ but $g\nmid c$.  Therefore $g$ is squarefree and
\[
 (7+25n)b+1=g(qn+c)
\]
is squarefree exactly when $qn+c$ is squarefree.

The $X$ values $qn+c$, $0\leq n<X$, are precisely the positive integers in
the class $c\pmod q$ up to
\[
 x=q(X-1)+c.
\]
Here $c\leq q$, $q\leq25N$, $x/q\leq X$, and
\[
 \frac{x}{q}\geq X-1,\qquad
 \frac N{25}-1\leq X\leq\frac N{25}+1,
 \qquad x\leq N^2+25N.
\]
Take $D=\lceil3.2N^{3/4}\rceil$ in
\eqref{eq:explicit-squarefree-lower}.  Since
\[
 \prod_{p\nmid q}\left(1-\frac1{p^2}\right)
 \geq\prod_p\left(1-\frac1{p^2}\right)
 =\frac6{\pi^2}>0.6079
\]
and $q^{1/4}\leq\sqrt5N^{1/4}$, the number of squarefree values is at least
$\mathcal L(N)$.
After division by $N$, every subtracted term is decreasing for
$N\geq2\cdot10^9$, while the positive main term is increasing.  Using
$\sqrt5<2.237$, at that endpoint this lower bound exceeds
$4.14\cdot10^6$ and its quotient by $N$ exceeds $0.00207$.  Thus a
squarefree neighbour exists, and the normalized lower bound remains above
$0.00207$ thereafter.
For a nonempty $B$, choose $b\in B$.  Then
$\Gamma_N(\{b\})\subseteq\Gamma_N(B)$, so
$|\Gamma_N(B)|\geq\mathcal L(N)\geq|B|$ whenever the displayed size
hypothesis holds.  This is precisely the Hall inequality.
\end{proof}

For the matching range we can reuse the primorial progression error instead
of the coarser all-parameter singleton bound.  Put
\[
 N_{4096}=\frac{8P_{11}}5=5936590507848,
 \qquad
 H_{\rm deg}(N)=\begin{cases}4096,&N<N_{4096},\\8192,&N\geq N_{4096},\end{cases}
\]
and set
\[
 \kappa_5=\frac{12769}{20164}.
\]
For $1.8263\cdot10^{11}\leq N<N_{25}$, put
\[
 \mathcal M(N)=\frac{\kappa_5N}{25}-\frac7{25}
                -N\mathcal E_{25,H_{\rm deg}(N)}(N).
\]

\begin{proposition}[Primorial one-pivot degree]
\label{prop:primorial-one-pivot-degree}
On this interval every outside vertex has at least $\mathcal M(N)$
squarefree neighbours in $C_N$, and every vertex of $C_N$ has at least the
 same number in $C_N^{18}$.  Moreover $\mathcal M(N)/N$ is increasing on each
 of the two intervals separated by $N_{4096}$, and
\[
 \mathcal M(1.8263\cdot10^{11})>\frac N{50}+2.8\cdot10^4.
\]
\end{proposition}

\begin{proof}
For an outside pivot, Corollary~1 of Sothanaphan's progression estimate gives
the bad-value upper bound used in \eqref{eq:explicit-off-diagonal}.  The
effective root modulus is at most $5N$.  The least moduli with root count
greater than $4096$ and $8192$ are $8P_{11}$ and $Q_*$, respectively, so
Lemma~\ref{lem:primorial-progression} allows the displayed $H_{\rm deg}$.  Also
\[
 \prod_{(p,25b)=1}\left(1-\frac1{p^2}\right)
 \geq\prod_{p\neq5}\left(1-\frac1{p^2}\right)
 =\frac{25}{4\pi^2}>\kappa_5.
\]
Indeed the classical bound $\pi<355/113$ gives
\[
 \frac{25}{4\pi^2}>\frac{25}{4}\left(\frac{113}{355}\right)^2
 =\frac{12769}{20164}=\kappa_5.
\]
Since $|C_N|\geq N/25-7/25$, subtraction gives the claimed lower degree.
 The reverse $7$--$18$ calculation is identical.  Every nonconstant term of
 $\mathcal E_{25,H_{\rm deg}}$ decreases while $H_{\rm deg}$ is fixed.  Direct
 evaluation gives the displayed lower-end margin and
 $\mathcal M(N_{4096})/N_{4096}>0.0222$ immediately after the sole jump.
 This proves both the asserted piecewise monotonicity and the strict
 half-class lower bound throughout the interval.
\end{proof}

For $4.5654\cdot10^{10}\leq N<3.4375\cdot10^{11}$ define the correlated
degree envelope
\[
 \mathcal M_1(N)=\frac{\kappa_5N}{25}-\frac7{25}
                 -N\mathcal E_{25,1024}(N).
\]

\begin{proposition}[Root-product correlated degree]
\label{prop:correlated-one-pivot-degree}
Throughout this interval every outside vertex has at least $\mathcal M_1(N)$
squarefree neighbours in $C_N$, and the same reverse degree holds from $C_N$
to $C_N^{18}$.  Moreover $\mathcal M_1(N)/N$ is increasing and
\[
 \mathcal M_1(4.5654\cdot10^{10})>\frac N{50}+800.
\]
\end{proposition}

\begin{proof}
Let $m=2^\nu m_0\leq5N$ be the effective root modulus.  If
$\rho(m)\leq1024$, the proof of Proposition~
\ref{prop:primorial-one-pivot-degree} applies directly with $H=1024$.
Suppose instead that $\rho(m)>1024$.  Since $5N<P_{11}$ throughout the
present interval, the cases $\nu\leq1$ are impossible: they require at least
eleven odd prime factors.  Removing a factor $5$ did not change the
$2$-adic valuation, so $\nu\geq2$ implies $4\mid b$.  The Euler product then omits
$p=2$ as well as $p=5$, and is at least $4\kappa_5/3$; meanwhile
$\rho(m)\leq4096$, since the least modulus with more than $4096$ roots is
$8P_{11}>5N$.  Thus the $H=4096$ lower degree is at least the asserted
$H=1024$ lower degree provided
\[
 \mathcal E_{25,4096}(N)-\mathcal E_{25,1024}(N)
 \leq\frac{\kappa_5}{75}.
\]
The left side decreases; at $N=4.5654\cdot10^{10}$ it is below $0.003132$,
whereas $\kappa_5/75>0.008443$.  This proves the correlated bound in both
directions.  The monotonicity and endpoint margin follow directly from the
formula for $\mathcal E_{25,1024}$.
\end{proof}

For $2.2827\cdot10^{10}\leq N<2.47\cdot10^{11}$ put
\[
 \mathcal M_{1/2}(N)=\frac{\kappa_5N}{25}-\frac7{25}
                 -N\mathcal E_{25,512}(N).
\]

\begin{proposition}[Second root--Euler correlated degree]
\label{prop:half-root-correlated-degree}
Throughout this interval every outside vertex has at least
$\mathcal M_{1/2}(N)$ squarefree neighbours in $C_N$, and the same reverse
degree holds from $C_N$ to $C_N^{18}$.  The quotient
$\mathcal M_{1/2}(N)/N$ increases, and
\[
 \mathcal M_{1/2}(2.2827\cdot10^{10})>\frac N{50}+411.
\]
\end{proposition}

\begin{proof}
Let $m=2^\nu m_0\leq5N$ be the effective root modulus.  If
$\rho(m)\leq512$, use the progression estimate directly with $H=512$.
Suppose $\rho(m)>512$ and $\nu\leq1$.  Then $m_0$ has at least ten distinct
odd prime factors.  At most one is the modulus prime $5$, so the pivot is
divisible by at least nine distinct odd primes other than $5$.  If $3$ did
not divide the pivot, then
\[
 b\geq7\cdot11\cdot13\cdot17\cdot19\cdot23\cdot29\cdot31\cdot37
 =247357937827>N,
\]
a contradiction.  Hence the Euler product omits $3$ and is at least
$9\kappa_5/8$.  Also $\rho(m)\leq1024$, because $5N<P_{11}$, and at the lower
endpoint
\[
 \mathcal E_{25,1024}-\mathcal E_{25,512}<0.001386
 <\frac{\kappa_5}{200}.
\]
Thus the missing-$3$ main-term gain pays the larger root error.

If $\nu\geq2$, then $4\mid b$, the Euler product omits $2$, and
$\rho(m)\leq4096$.  At the same endpoint
\[
 \mathcal E_{25,4096}-\mathcal E_{25,512}<0.005332
 <\frac{\kappa_5}{75},
\]
so the missing-$2$ gain pays this error.  Both differences decrease.  The
reverse progression has the identical alternatives.  This proves the degree
bounds; monotonicity and the endpoint margin follow by direct evaluation.
\end{proof}

Put
\[
 \beta_8:=0.00097181149.
\]
For $1.382\cdot10^{10}\leq N\leq2.2827\cdot10^{10}$ define
\[
 \mathcal M_{1/4}(N)=\frac{\kappa_5N}{25}-\frac7{25}
                 -N\mathcal E_{25,512}(N)+\beta_8N.
\]

\begin{proposition}[Eight-prime root--Euler correlated degree]
\label{prop:eight-prime-correlated-degree}
Throughout this interval every outside vertex has at least
$\mathcal M_{1/4}(N)$ squarefree neighbours in $C_N$, and the same reverse
degree holds from $C_N$ to $C_N^{18}$.  The quotient
$\mathcal M_{1/4}(N)/N$ increases, and
\[
 \mathcal M_{1/4}(1.382\cdot10^{10})>\frac N{50}+16356.
\]
\end{proposition}

\begin{proof}
If the effective root count is at most $256$, use $H=256$.  Direct endpoint
evaluation gives
\[
 \mathcal E_{25,512}-\mathcal E_{25,256}>0.001099>\beta_8
\]
throughout the interval.

Suppose next that the root count exceeds $256$ and the effective modulus has
$2$-adic valuation at most one.  It has at least nine distinct odd prime
factors.  At most one is the modulus prime $5$, so the pivot is divisible by
at least eight distinct non-$5$ odd primes.  Exact ordered-prime enumeration
under product at most $2.2827\cdot10^{10}$ gives the unique least Euler gain
at
\[
 \{7,13,17,19,23,29,31,37\},
 \qquad \prod p=22487085257,
\]
and
\[
 \prod_p\left(1-\frac1{p^2}\right)^{-1}-1
 >0.038365586453.
\]
Consequently its normalized main-term gain is greater than
\[
 \frac{\kappa_5}{25}(0.038365586453)
 >0.000971811492>\beta_8.
\]
If the $2$-adic valuation is at least two, the missing-$2$ gain
$\kappa_5/75$ is larger still.  Thus every root count at most $512$ pays the
claimed improved degree.

For the remaining higher-root strata, at the lower endpoint
\[
 \mathcal E_{25,1024}-\mathcal E_{25,512}+\beta_8<0.002610
 <\frac{\kappa_5}{200},
\]
and
\[
 \mathcal E_{25,4096}-\mathcal E_{25,512}+\beta_8<0.007275
 <\frac{\kappa_5}{75}.
\]
The forced missing-$3$ and missing-$2$ gains from the proof of
Proposition~\ref{prop:half-root-correlated-degree} therefore continue to pay
those strata.  The reverse progression is identical.  Monotonicity and the
endpoint margin follow from the displayed formula.
\end{proof}

Put
\[
 \beta'_8:=0.00111752951,
 \qquad
 \mathcal M_{1/8}(N)=\frac{\kappa_5N}{25}-\frac7{25}
                 -N\mathcal E_{25,512}(N)+\beta'_8N.
\]

\begin{proposition}[Second eight-prime root--Euler degree]
\label{prop:second-eight-prime-degree}
For $1.29\cdot10^{10}\leq N\leq1.382\cdot10^{10}$, every outside vertex
has at least $\mathcal M_{1/8}(N)$ squarefree neighbours in $C_N$, with the
same reverse degree.  Moreover
\[
 \mathcal M_{1/8}(1.29\cdot10^{10})>\frac N{50}+6728.
\]
\end{proposition}

\begin{proof}
Repeat Proposition~\ref{prop:eight-prime-correlated-degree}, but impose the
current product bound $1.382\cdot10^{10}$.  Complete ordered-prime
enumeration now visits $828$ nodes and is uniquely minimized by
\[
 \{7,11,13,17,23,29,31,37\},
 \qquad \prod p=13018838833.
\]
Exact rational multiplication gives
\[
 \prod_p\left(1-\frac1{p^2}\right)^{-1}-1
 =\frac{13286935886873651}{301166059585536000}
 >0.044118304383,
\]
and hence normalized gain greater than
$0.001117529515>\beta'_8$.  The low-root saving
$\mathcal E_{25,512}-\mathcal E_{25,256}$ exceeds $0.001299$, while at the
lower endpoint the two remaining comparisons are
\[
 \mathcal E_{25,1024}-\mathcal E_{25,512}+\beta'_8<0.002794
 <\frac{\kappa_5}{200}
\]
and
\[
 \mathcal E_{25,4096}-\mathcal E_{25,512}+\beta'_8<0.007568
 <\frac{\kappa_5}{75}.
\]
Thus every root stratum pays the asserted degree.  Direct endpoint evaluation
gives the margin.
\end{proof}

Put
\[
 \beta''_8:=0.00113583855,
 \qquad
 \mathcal M_{1/16}(N)=\frac{\kappa_5N}{25}-\frac7{25}
                 -N\mathcal E_{25,512}(N)+\beta''_8N.
\]

\begin{proposition}[Third eight-prime root--Euler degree]
\label{prop:third-eight-prime-degree}
For $1.279\cdot10^{10}\leq N\leq1.29\cdot10^{10}$, every outside vertex
has at least $\mathcal M_{1/16}(N)$ squarefree neighbours in $C_N$, with the
same reverse degree.  Moreover
\[
 \mathcal M_{1/16}(1.279\cdot10^{10})>\frac N{50}+5090.
\]
\end{proposition}

\begin{proof}
At product bound $1.29\cdot10^{10}$, complete enumeration is uniquely
minimized by
\[
 \{7,11,13,17,19,29,31,43\},
 \qquad \prod p=12498697211.
\]
Its exact relative Euler gain is
\[
 \frac{12438467965272539}{277389791723520000}>0.044841116495,
\]
which gives normalized gain $>0.001135838556>\beta''_8$.  At the lower
endpoint the low-root saving and the two higher-root comparisons remain
strict:
\[
 \mathcal E_{25,512}-\mathcal E_{25,256}>\beta''_8,
\]
\[
 \mathcal E_{25,1024}-\mathcal E_{25,512}+\beta''_8
 <0.002799<\frac{\kappa_5}{200},
\]
\[
 \mathcal E_{25,4096}-\mathcal E_{25,512}+\beta''_8
 <0.007586<\frac{\kappa_5}{75}.
\]
Thus the same root-stratified degree proof applies, and endpoint evaluation
gives the margin.
\end{proof}

Put
\[
 \mathcal M_{\rm QR2}(N)=\frac{\kappa_5N}{25}-\frac7{25}
                 -N\mathcal E_{25,128}(N).
\]

\begin{proposition}[Quadratic-residue root hierarchy]
\label{prop:qr-root-hierarchy}
For $7.67769\cdot10^9\leq N\leq1.15\cdot10^{10}$, every outside vertex has
at least $\mathcal M_{\rm QR2}(N)$ squarefree neighbours in $C_N$, with the
same reverse degree.  Moreover
\[
 \mathcal M_{\rm QR2}(7.67769\cdot10^9)>\frac N{50}+3853771.
\]
\end{proposition}

\begin{proof}
The root count is at most $2048$: more roots would require at least nine
non-$5$ odd prime divisors, whose least product is
\[
 3\cdot7\cdot11\cdot13\cdot17\cdot19\cdot23\cdot29\cdot31
 =20056049013>N.
\]
If the root count is at most $128$, use
Lemma~\ref{lem:primorial-progression} directly.  For the four remaining
root strata, let $k$ be the number of non-$5$ odd prime divisors.  Root count
in $(128,256]$, $(256,512]$, $(512,1024]$, or $(1024,2048]$ forces
$k\geq5,6,7,8$, respectively.

Each stratum also forces enough small prime divisors.  If the asserted number
were absent, the least possible radical would already exceed the upper
endpoint:
\[
\begin{array}{c|c|c|c}
H&\text{forced small primes}&\text{excluded least product}&\text{value}\\ \hline
256&3\text{ primes }\leq1000&3\cdot7\cdot1009\cdot1013\cdot1019&21872281683\\
512&4\text{ primes }\leq360&3\cdot7\cdot11\cdot367\cdot373\cdot379&11984670159\\
1024&5\text{ primes }\leq160&3\cdot7\cdot11\cdot13\cdot163\cdot167\cdot173&14141826699\\
2048&6\text{ primes }\leq60&3\cdot7\cdot11\cdot13\cdot17\cdot61\cdot67\cdot71&14813826027.
\end{array}
\]

Choose $D=20000000$.  Then the complementary variable satisfies
\[
 m\leq\frac{N^2+25N+1}{D^2}\leq330626.
\]
Intersect every possible collection of the forced quadratic-residue masks.
The exact maxima are
\[
\begin{array}{c|c|c|c}
H&\text{mask cutoff/size}&\max\#m&H\max\#m\\ \hline
256&1000/3&42697&10930432\\
512&360/4&21570&11043840\\
1024&160/5&10663&10918912\\
2048&60/6&4798&9826304.
\end{array}
\]
Thus every higher-root tail is at most $11043840$.  The calculation is
reproduced by
\path{scripts/quadratic_residue_hierarchy_scan.py}.

Discarding every Euler gain, the progression proof gives at the new lower
endpoint more than
\[
 \frac{\kappa_5(N-43)}{25}-D-\frac{N/25+1}{D}
 -\frac1{5D}-11043840-1
 >\frac N{50}+9933718.
\]
This lower bound increases with $N$, whereas
$\mathcal M_{\rm QR2}(N)-N/50<1.277\cdot10^7$ throughout the interval.
Hence every higher-root stratum also supplies the asserted degree.  Direct
evaluation gives the displayed endpoint margin.
\end{proof}

Put
\[
 \mathcal M_{\rm QR}(N)=\frac{\kappa_5N}{25}-\frac7{25}
                 -N\mathcal E_{25,256}(N).
\]

\begin{proposition}[Quadratic-residue screened degree]
\label{prop:qr-screened-degree}
For $1.15\cdot10^{10}\leq N\leq1.2643\cdot10^{10}$, every outside vertex
has at least $\mathcal M_{\rm QR}(N)$ squarefree neighbours in $C_N$, with
the same reverse degree.  Moreover
\[
 \mathcal M_{\rm QR}(1.15\cdot10^{10})>\frac N{50}+154290.
\]
\end{proposition}

\begin{proof}
If the transformed root count is at most $256$, this is
Lemma~\ref{lem:primorial-progression} with $H=256$.  We replace the crude
root payment only in the stratum $256<\rho(d_1)\leq512$ whose $2$-adic
exponent is at most one.  There the effective modulus contains exactly eight
non-$5$ odd prime divisors $S$.  In the
large-square tail of Proposition~1 of Sothanaphan, the complementary variable
$m$ and a root $d$ satisfy
\[
 d^2m\equiv c_1\pmod {d_1}.
\]
For every $p\in S$ one has $c_1\equiv1\pmod p$; hence any contributing $m$
must be a nonzero quadratic residue modulo every $p\in S$.

Choose $D=8000000$.  On the whole interval,
\[
 m\leq\frac{N^2+25N+1}{D^2}\leq2497585.
\]
Complete ordered-prime enumeration at the upper endpoint gives $5700$
feasible supports, involving $101$ possible primes.  Intersecting the exact
quadratic-residue masks through $2497585$ leaves at most $6024$ values of
$m$ for every support.  Since each has at most $512$ roots, the complete tail
is at most
\[
 512\cdot6024=3084288.
\]
The scan is reproduced by
\texttt{scripts/quadratic\_residue\_support\_scan.py}.

We may ignore the Euler gain from all eight primes.  If $X$ is either base
progression length, then $X-1\geq(N-43)/25$, $X\leq N/25+1$, and the proof of
the progression estimate now gives at least
\[
 \frac{\kappa_5(N-43)}{25}-D-\frac{N/25+1}{D}
 -\frac1{5D}-3084288-1
 >\frac N{50}+50214005
\]
squarefree values at the lower endpoint.  The left side minus $N/50$
increases, while $\mathcal M_{\rm QR}(N)-N/50$ is below $2.26\cdot10^6$ on
the interval.  Thus this root stratum more than pays the asserted degree.

For root count above $512$ with $2$-adic exponent at most one, the forced
missing-$3$ gain applies and, at the lower endpoint,
\[
 \mathcal E_{25,1024}-\mathcal E_{25,256}
 <0.003124<\frac{\kappa_5}{200}.
\]
For $2$-adic exponent at least two, the missing-$2$ gain applies and
\[
 \mathcal E_{25,4096}-\mathcal E_{25,256}
 <0.008083<\frac{\kappa_5}{75}.
\]
Both differences decrease with $N$.  This covers every root stratum in both
directions.  Direct endpoint substitution gives the stated margin.
\end{proof}

Put
\[
 \beta'''_8:=0.00116052,
 \qquad
 \mathcal M_{1/32}(N)=\frac{\kappa_5N}{25}-\frac7{25}
                 -N\mathcal E_{25,512}(N)+\beta'''_8N.
\]

\begin{proposition}[Modulus-correlated eight-prime degree]
\label{prop:modulus-correlated-eight-prime-degree}
For $1.2643\cdot10^{10}\leq N\leq1.279\cdot10^{10}$, every outside vertex
has at least $\mathcal M_{1/32}(N)$ squarefree neighbours in $C_N$, with the
same reverse degree.  Throughout this interval the degree is strictly larger
than half a base class, and at the lower endpoint
\[
 \mathcal M_{1/32}(1.2643\cdot10^{10})>\frac N{50}+1442.
\]
\end{proposition}

\begin{proof}
We retain the transformed modulus in the proof of
Lemma~\ref{lem:primorial-progression}.  For $q_0=25$, root bound $H$, pivot
$b$, and $X\leq N/25+1$, the last term before the uniform enlargement
$d_1\leq25N$ is
\[
 \frac{Hd_1X}{D^2}\leq\frac{25HbX}{D^2}.
\]
Consequently, with $H=512$ and
$d_{512}=1024^{1/3}N^{2/3}$, the actual error is at most
\begin{equation}\label{eq:retained-pivot-saving}
 N\mathcal E_{25,512}(N)
 -\frac{512\cdot25(N-b)(N/25+1)}{d_{512}^2}.
\end{equation}
The reciprocal $d_1$ term is still bounded with $d_1\geq5$, so this is a
one-sided strengthening of the already proved envelope.

Only the root stratum $256<\rho(d_1)\leq512$ with odd $b$ needs the new
correlation.  In this stratum $b$ has exactly eight distinct non-$5$ odd
prime divisors.  Write their set as $S$ and their product as $r(S)$.  Every
prime divisor of $b$ lies in $S\cup\{5\}$; hence
\[
 b=r(S)k,\qquad k\leq N/r(S),\qquad
 p\mid k\Longrightarrow p\in S\cup\{5\}.
\]
For fixed $S$, the Euler gain is
\[
 G(S)=\frac{\kappa_5}{25}
 \left(\prod_{p\in S}(1-p^{-2})^{-1}-1\right),
\]
whereas \eqref{eq:retained-pivot-saving} is worst at the largest admissible
$k$.  Complete ordered-prime enumeration through the upper endpoint visits
$5831$ supports and $157$ support/multiplier breakpoints.  Direct
differentiation shows that the degree for each fixed admissible pair $(S,b)$
increases between consecutive breakpoints.  The unique global minimum is at
$N=12643000000$, with
\[
 S=\{7,11,13,17,19,29,31,43\},\qquad
 r(S)=b=12498697211,
\]
and exact Euler contribution
\[
 G(S)=\frac{158826797448565050491}
 {139832194007826432000000}.
\]
Adding the retained-modulus saving gives
\[
 G(S)+\frac{512\cdot25(N-b)(N/25+1)}{d_{512}^2N}
 >0.001160529889>\beta'''_8.
\]
This finite calculation is reproduced by
\texttt{scripts/modulus\_correlated\_eight\_prime\_scan.py}.

If the root count is at most $256$, then at the lower endpoint
\[
 \mathcal E_{25,512}-\mathcal E_{25,256}>0.00133914>\beta'''_8.
\]
If $b$ is even, the missing-$2$ gain is already larger than the asserted
bonus.  The remaining higher-root strata satisfy
\[
 \mathcal E_{25,1024}-\mathcal E_{25,512}+\beta'''_8
 <0.002848<\frac{\kappa_5}{200},
\]
\[
 \mathcal E_{25,4096}-\mathcal E_{25,512}+\beta'''_8
 <0.007653<\frac{\kappa_5}{75},
\]
so the forced missing-$3$ and missing-$2$ gains pay them exactly as before.
The reverse progression has the same alternatives.  Endpoint substitution
gives the displayed margin; the breakpoint enumeration and piecewise monotonicity
preserve it on the full interval.
\end{proof}

Define the degree function
\[
 \mathcal D(N)=
 \begin{cases}
  \mathcal M_{\rm QR2}(N),&7.67769\cdot10^9\leq N<1.15\cdot10^{10},\\
  \mathcal M_{\rm QR}(N),&1.15\cdot10^{10}\leq N<1.2643\cdot10^{10},\\
  \mathcal M_{1/32}(N),&1.2643\cdot10^{10}\leq N<1.279\cdot10^{10},\\
  \mathcal M_{1/16}(N),&1.279\cdot10^{10}\leq N<1.29\cdot10^{10},\\
  \mathcal M_{1/8}(N),&1.29\cdot10^{10}\leq N<1.382\cdot10^{10},\\
  \mathcal M_{1/4}(N),&1.382\cdot10^{10}\leq N<2.2827\cdot10^{10},\\
  \mathcal M_{1/2}(N),&2.2827\cdot10^{10}\leq N<2.47\cdot10^{11},\\
  \mathcal M_1(N),&2.47\cdot10^{11}\leq N<3.4375\cdot10^{11},\\
  \mathcal M(N),&3.4375\cdot10^{11}\leq N<N_{25},\\
  \mathcal L(N),&N\geq N_{25}.
 \end{cases}
\]

\begin{corollary}[Global matching of the opposite base class]
\label{cor:opposite-base-matching}
Let
\[
 C_N^{18}=\{b\in[1,N]:b\equiv18\pmod {25}\}.
\]
For every $N\geq7.67769\cdot10^9$, the bipartite squarefree-edge graph between
$C_N^{18}$ and $\CN$ has a matching that saturates $C_N^{18}$.  Consequently
every outside set $B\subseteq C_N^{18}$ satisfies the Hall inequality
\eqref{eq:hall-defect}.
\end{corollary}

\begin{proof}
Below $1.15\cdot10^{10}$ use Proposition~\ref{prop:qr-root-hierarchy},
from there to $1.2643\cdot10^{10}$ use Proposition~\ref{prop:qr-screened-degree},
from there to $1.279\cdot10^{10}$ use Proposition~
\ref{prop:modulus-correlated-eight-prime-degree}, from there to
$1.29\cdot10^{10}$ use Proposition~\ref{prop:third-eight-prime-degree},
from there to $1.382\cdot10^{10}$ use
Proposition~\ref{prop:second-eight-prime-degree},
from there to $2.2827\cdot10^{10}$ use
Proposition~\ref{prop:eight-prime-correlated-degree}; from there to
$2.47\cdot10^{11}$, Proposition~
\ref{prop:half-root-correlated-degree} gives degree at least $\mathcal D(N)$
in both directions; from there to $3.4375\cdot10^{11}$ use Proposition~
\ref{prop:correlated-one-pivot-degree}, and from there to $N_{25}$ use Proposition~
\ref{prop:primorial-one-pivot-degree}.  Above that point, Theorem~
\ref{thm:singleton-large-close} gives degree at least $\mathcal D(N)$
from every vertex of $C_N^{18}$ into $\CN$.  Applying the same proof with the
roles reversed and the progression $b\equiv18\pmod {25}$ gives the same lower
bound from every vertex of $\CN$ into $C_N^{18}$: in that calculation
$(18a+1,25a)=1$ for $a\equiv7\pmod {25}$, and all the inequalities used in
Theorem~\ref{thm:singleton-large-close} remain valid.

The quotient $\mathcal L(N)/N$ is increasing.  Direct substitution at
$N=1.409\cdot10^{12}$ gives
\[
 \mathcal L(N)>\frac N{50}+5.8\cdot10^5>
 \frac12\max\{|C_N^{18}|,|\CN|\}.
\]
For the lower intervals, Propositions~\ref{prop:qr-root-hierarchy},
\ref{prop:qr-screened-degree},
\ref{prop:modulus-correlated-eight-prime-degree},
\ref{prop:third-eight-prime-degree},
\ref{prop:second-eight-prime-degree},
\ref{prop:eight-prime-correlated-degree},
\ref{prop:half-root-correlated-degree},
\ref{prop:correlated-one-pivot-degree}, and
\ref{prop:primorial-one-pivot-degree}, together with their endpoint and switch
evaluations, give the same strict half-size inequality with $\mathcal D(N)$.
Write
$L=C_N^{18}$ and $R=\CN$.  If $S\subseteq L$ is nonempty and
$|S|\leq\mathcal D(N)$, the neighbourhood of any member of $S$ already has
size at least $|S|$.  If $|S|>\mathcal D(N)$ and $|\Gamma(S)|<|S|$, then,
because $|L|\leq|R|$, there is a vertex $y\in R\setminus\Gamma(S)$.  All
neighbours of $y$ lie in $L\setminus S$, so
\[
 \deg(y)\leq|L|-|S|<|L|-\mathcal D(N)<\mathcal D(N),
\]
contradicting the reverse minimum-degree bound.  Hall's theorem supplies the
matching.  Restricting it to any $B\subseteq C_N^{18}$ proves the last claim.
\end{proof}

The half-minimum-degree argument is not needed below the preceding endpoint.
Splitting both base classes by parity makes the square witnesses $2$ and $5$
identically inactive, while a Hall defect itself supplies a short triple.
This gives a systemic matching down to the one-pivot scale.

Put
\[
 L_0=\{x\leq N:x\equiv18\pmod {50}\},\qquad
 R_0=\{y\leq N:y\equiv7\pmod {50}\},
\]
and
\[
 L_1=\{x\leq N:x\equiv43\pmod {50}\},\qquad
 R_1=\{y\leq N:y\equiv32\pmod {50}\}.
\]
Thus $C_N^{18}=L_0\sqcup L_1$, $C_N=R_0\sqcup R_1$, and
$|L_i|\leq|R_i|$.  Join $L_i$ to $R_i$ when the product plus one is
squarefree.

\begin{lemma}[Zero-coincidence three-form block]
\label{lem:zero-coincidence-triple-block}
Let $x_1,x_2,x_3$ be distinct elements of one $L_i$ with
$\max x_j-\min x_j\leq200$.  The number of $y\in R_i$ for which each of
$x_1y+1,x_2y+1,x_3y+1$ has a prime-square divisor at most $353$ is less than
\begin{equation}\label{eq:zero-coincidence-triple-block}
 0.003707\left(\frac N{50}+1\right)+1628952.
\end{equation}
\end{lemma}

\begin{proof}
Let $\mathcal P=\{p:3\leq p\leq353,\ p\neq5\}$ and put $z_p=p^{-2}$.
There are $69$ such primes.  No residue classes belonging to two different
forms can coincide at a prime in $\mathcal P$: coincidence would give
$p^2\mid x_j-x_k$, whereas $50\mid x_j-x_k$, so a nonzero difference would
have absolute value at least $50p^2\geq450$.

For a list $w=(w_p)_{p\in\mathcal P}$ write
\[
 L_3(w)=1-e_1(w)+e_2(w)-e_3(w),\qquad
 U_2(w)=1-e_1(w)+e_2(w),
\]
where the $e_j$ are elementary symmetric functions.  The degree-$3$ lower
and degree-$2$ upper Bonferroni inequalities applied to the survivor events
give the common-bad density at most
\[
 I_3:=1-3L_3(z)+3U_2(2z)-L_3(3z).
\]
Exact rational evaluation gives
\begin{equation}\label{eq:systemic-triple-constants}
 \sum_{p\in\mathcal P}\frac1{p^2}<0.161841,\qquad
 I_3<0.003707.
\end{equation}
Every nonconstant CRT term has one terminal incomplete class.  The total
number of such terms in the displayed Bonferroni combination is
\[
 3\sum_{j=1}^3e_j(1,\ldots,1)
 +3\sum_{j=1}^2e_j(2,\ldots,2)
 +\sum_{j=1}^3e_j(3,\ldots,3)
 =1628952.
\]
This proves \eqref{eq:zero-coincidence-triple-block}.
\end{proof}

For $N\geq2\cdot10^9$ define
\[
 P(N)=\frac{N}{19(\log(N/20)-1.1)},\qquad
 \tau=0.000407315,
\]
and let
\[
 H(N)=
 \begin{cases}
 1024,&2\cdot10^9\leq N<2587877292,\\
 2048,&2587877292\leq N<7.67769\cdot10^9.
 \end{cases}
\]
Put
\begin{align}
 E_H(N)&=\left(\frac N{50}+1\right)\tau+P(N)+400.801H(N),
 \label{eq:systemic-pure-tail}\\
 D_H(N)&=\frac N{50}-1
 -0.161841\left(\frac N{50}+1\right)-69-E_H(N),
 \label{eq:systemic-pure-degree}\\
 U_H(N)&=0.003707\left(\frac N{50}+1\right)
 +1628952+3E_H(N).
 \label{eq:systemic-pure-codegree}
\end{align}

\begin{lemma}[Uniform medium-prime payment in the parity blocks]
\label{lem:systemic-pure-medium-tail}
On $2\cdot10^9\leq N<7.67769\cdot10^9$, every vertex of each graph
$L_i$--$R_i$ has degree at least $D_H(N)$, in both directions.  Moreover,
every triple in Lemma~\ref{lem:zero-coincidence-triple-block} has common
non-neighbourhood in $R_i$ of size at most $U_H(N)$, and
\begin{equation}\label{eq:systemic-pure-strict-margin}
 D_H(N)>\frac N{100},\qquad U_H(N)<D_H(N).
\end{equation}
\end{lemma}

\begin{proof}
For a pivot $t$ in either block and a variable in the opposite block,
$2$ and $5$ cannot be square witnesses.  For primes through $353$, the
union bound and \eqref{eq:systemic-triple-constants} give the first two
subtractions in \eqref{eq:systemic-pure-degree}.

For the remaining primes take $Y=\lfloor N/19\rfloor$.  The density part is
at most $(N/50+1)\tau$, since the rational prime-square tail calculation gives
\[
 \sum_{p>353}\frac1{p^2}<0.000407315.
\]
For $353<p\leq Y$, pay one terminal class per prime.  Dusart's estimate
\eqref{eq:dusart-pi}, together with $N/20\leq Y\leq N/19$, gives
\[
 \pi(Y)\leq P(N).
\]

For $p>Y$, write the variable as $r+50n$.  The transformed progression has
modulus $50t$ and coprime initial term $rt+1$.  Removing one redundant factor
$5$ leaves a root modulus at most $10N$.  The least modulus with more than
$1024$ unit square roots is
\[
 8P_9=25878772920,
\]
and the least with more than $2048$ is
\[
 8P_{10}=802241960520.
\]
This proves the displayed step bound $H(N)$.  The standard complementary
variable split used in Lemma~\ref{lem:exact-split-prime-root} bounds the
large-prime terminal classes by
\[
 H(N)\left(
 \frac{2(N/50+2)}Y+
 \frac{N^2+50N+1}{Y^2}
 \right)<400.801H(N).
\]
Thus the full tail is bounded by $E_H(N)$, proving the degree assertion.
For a common non-neighbour of three pivots, either all witnesses are at most
$353$, when Lemma~\ref{lem:zero-coincidence-triple-block} applies, or at
least one form uses a larger witness; the union bound pays $3E_H(N)$.  This
gives \eqref{eq:systemic-pure-codegree}.

With $H=1024$, direct substitution at $N=2\cdot10^9$ gives
\[
 \frac{D_H(N)}N>0.01351113,\qquad
 \frac{U_H(N)}N<0.01064464.
\]
At the sole root jump, $N=2587877292$ and $H=2048$, it gives
\[
 \frac{D_H(N)}N>0.01344370,\qquad
 \frac{U_H(N)}N<0.01066191.
\]
On each constant-$H$ interval, $D_H(N)/N$ increases and $U_H(N)/N$
decreases.  The two endpoint checks prove
\eqref{eq:systemic-pure-strict-margin}.
\end{proof}

\begin{theorem}[Systemic pure-$18$ parity matching]
\label{thm:systemic-pure-parity-matching}
For every $N\geq2\cdot10^9$, the squarefree-edge graph between
$C_N^{18}$ and $C_N$ has a matching that saturates $C_N^{18}$.
Consequently every outside set $B\subseteq C_N^{18}$ satisfies the original
Hall inequality \eqref{eq:hall-defect}.
\end{theorem}

\begin{proof}
The range $N\geq7.67769\cdot10^9$ is
Corollary~\ref{cor:opposite-base-matching}, so consider the preceding
interval.  Fix one block $L_i$--$R_i$ and suppose $S\subseteq L_i$ is a Hall
defect.  Write $m=|S|$.  If $m\leq D_H(N)$, the neighbourhood of one member
of $S$ already has size at least $m$, a contradiction.  Hence
\[
 m>D_H(N)>\frac N{100}.
\]
List $S$ increasingly as $x_1<\cdots<x_m$.  Since
\[
 \sum_{j=1}^{m-2}(x_{j+2}-x_j)\leq2(N-1),
\]
some three consecutive members have span less than $201$, and hence at most
$200$.  Lemma~\ref{lem:systemic-pure-medium-tail} bounds their common
non-neighbourhood by $U_H(N)<D_H(N)$.

On the other hand let $T=R_i\setminus\Gamma(S)$.  Every neighbour of a
vertex of $T$ lies in $L_i\setminus S$, so the reverse degree bound gives
$m\leq|L_i|-D_H(N)$.  The assumed Hall defect therefore implies
\[
 |T|>|R_i|-m\geq|R_i|-|L_i|+D_H(N)\geq D_H(N),
\]
contradicting the common-non-neighbourhood bound.  Hall's theorem supplies a
matching in each parity block.  Their disjoint union saturates
$C_N^{18}=L_0\sqcup L_1$.  Restricting the matching to any
$B\subseteq C_N^{18}$ proves the last statement.
\end{proof}

The exact rational constants and both root-jump endpoint margins are
reproduced by
\begin{center}\small
\path{scripts/systematic_pure_matching_close.py}.
\end{center}

We next retain the quadratic-residue screen in the one-pivot progression
itself.  This is the ingredient needed to make the preceding matching robust
against a genuinely mixed outside set.

Put
\[
 X_0(N)=\frac{501}{25}+\frac{501}{N},\qquad M_0=251002,
 \qquad \eta_0=0.002267575.
\]
For $0\leq k\leq8$, define $s_k$ by
\[
\begin{array}{c|rrrrrrrrr}
k&0&1&2&3&4&5&6&7&8\\ \hline
s_k&251002&251002&251002&125464&63250&32653&16510&8125&3679.
\end{array}
\]

\begin{proposition}[QR-correlated mixed degree]
\label{prop:systemic-mixed-degree}
For
\[
 2\cdot10^9\leq N<7.67769\cdot10^9,
\]
every outside vertex has more than
\begin{equation}\label{eq:systemic-mixed-degree}
 \frac N{50}+4535153
\end{equation}
squarefree neighbours in $C_N$.  The same lower bound holds in the reverse
progression from $C_N$ into $C_N^{18}$.  More precisely, the normalized
excess in \eqref{eq:systemic-mixed-degree} is greater than
$0.0022675765$ at the lower endpoint and increases on the interval.
\end{proposition}

\begin{proof}
Take $D=\lfloor N/500\rfloor$.  Then $D\geq N/501$, and hence
\[
 \frac XD\leq X_0(N),\qquad
 \frac{N^2+25N+1}{D^2}<251002=M_0.
\]
Let $k$ be the number of non-$5$ odd prime divisors of the effective root
modulus.  Nine such divisors are impossible, because
\[
 3\cdot7\cdot11\cdot13\cdot17\cdot19\cdot23\cdot29\cdot31
 =20056049013>N.
\]
If the $2$-adic exponent is $\nu$, the root count is at most
\[
 H_{\nu,k}=c_\nu2^{k+1},\quad
 c_0=c_1=1,\quad c_2=2,\quad c_\nu=4\ (\nu\geq3).
\]

For $k\geq3$, product forcing supplies the following small subsets of the
odd prime divisors:
\[
\begin{array}{c|c|c|c}
k&\text{cutoff}&\text{forced subset}&\max\#\{m\leq M_0:m\text{ survives}\}\\ \hline
3&2000&1&125464\\
4&1400&2&63250\\
5&1000&3&32653\\
6&360&4&16510\\
7&160&5&8125\\
8&60&6&3679.
\end{array}
\]
Indeed, failure of the first two assertions would force respectively the
products
\[
 2003\cdot2011\cdot2017=8124542561,\quad
 3\cdot1409\cdot1423\cdot1427=8583434967,
\]
and the remaining four excluded products are the ones displayed in the
proof of Proposition~\ref{prop:qr-root-hierarchy}; all exceed the upper
endpoint.  Every contributing complementary variable $m$ is a nonzero
quadratic residue modulo each selected prime divisor.  Exact intersections
of those residue masks give the last column.  For $k\leq2$ we simply retain
all $M_0$ values, giving the first three entries $s_k=M_0$.

The same M\"obius split used in Proposition~\ref{prop:qr-screened-degree}
therefore gives, for an odd pivot, the lower bound
\begin{align}
 G_{0,k}(N)={}&\frac{\kappa_5(N-43)}{25}-\frac N{500}-X_0(N)
 -H_{0,k}\left(s_k+\frac{X_0(N)}2
       \left(4+\frac{s_k}{3026}\right)\right)
 -\frac{501}{5N}-1.                         \label{eq:mixed-odd-degree}
\end{align}
For an even pivot the Euler product omits the factor at $2$, so the first
term is instead $4\kappa_5(N-43)/75$; use the corresponding $c_\nu$ in
$H_{\nu,k}$.  These bounds discard the modulus-$5$ and $2$-adic residue
screens and every further Euler gain.

The term $s_k/3026$ cannot be discarded.  The reciprocal sum over the
permitted transformed parameters through $3025$ is at most $4$, while each
remaining survivor contributes at most $1/3026$.  Therefore
\[
 \sum_{m\ {\rm surviving}}\frac1m\leq4+\frac{s_k}{3026}.
\]
The elementary transformed-parameter count then gives the root
term in \eqref{eq:mixed-odd-degree}.

There are only $4\cdot9$ displayed valuation/root strata.  Direct exact
substitution at $N=2\cdot10^9$ shows that the minimum is the odd $k=6$
stratum:
\[
 G_{0,6}(N)-\frac N{50}>4535153.1398.
\]
The next smallest odd strata retain more than $4.55\cdot10^6$, and every
even stratum retains more than $1.50\cdot10^7$.  In each fixed stratum the
affine main term increases, while $X_0(N)$ and all normalized constant
payments decrease.  Feasible prime supports only disappear as $N$ decreases.
Thus the lower endpoint controls the full interval.  The reverse
$18\pmod {25}$ progression has the identical root and residue-screen
calculation.  This proves the proposition.
\end{proof}

The support forcing, all exact QR-mask intersections, and all $36$ endpoint
comparisons are reproduced by
\begin{center}\small
\path{scripts/systematic_mixed_degree_close.py}.
\end{center}
The reciprocal arithmetic is checked by
\begin{center}\small
\path{scripts/verify_four_range_paper_arithmetic.py}.
\end{center}

\begin{corollary}[A systemic mixed residual]
\label{cor:systemic-mixed-residual}
Let $B$ be a Hall defect in the interval of
Proposition~\ref{prop:systemic-mixed-degree}, and put
$B^*=B\setminus C_N^{18}$.  Then
\begin{equation}\label{eq:systemic-mixed-residual}
 |B^*|>0.00453515N.
\end{equation}
Consequently $B^*$, and hence $A^*=(B\cup T_N(B))\setminus
(C_N\cup C_N^{18})$, contains a mod-$4$ class with three members of span
less than $1764$.
\end{corollary}

\begin{proof}
The pure case is already excluded by
Theorem~\ref{thm:systemic-pure-parity-matching}, so choose $x\in B^*$.
Compatibility and the reverse degree bound give
\[
 |B\cap C_N^{18}|<|C_N^{18}|-\frac N{50}-4535153,
\]
whereas Hall defect and the forward degree give
$|B|>N/50+4535153$.  Since $|C_N^{18}|\leq N/25+1$,
\[
 |B^*|>9070305>0.00453515N
\]
at the lower endpoint; the normalized lower bound increases thereafter by
the preceding proof.

Let $S$ be the largest mod-$4$ class in $B^*$.  Then
$|S|>0.0011337875N$.  If $x_1<\cdots<x_m$ are its members, then
\[
 \sum_{j=1}^{m-2}(x_{j+2}-x_j)\leq2(N-1).
\]
At the lower endpoint this forces a three-term span below $1764$, and the
quotient decreases thereafter.  Hence an integral span is at most $1763$.
\end{proof}

We now give the terminal three-pivot payment.  Set
\[
 \tau_{131}=0.001257623,\quad P_-=13,\quad H_\Sigma=384.
\]
The first constant is Lemma~\ref{lem:three-term-prime-tail}, the second is
Lemma~\ref{lem:sharp-negative-pell-gap}, and the last follows from the QR
hierarchy: for $Y=\lfloor N/77\rfloor$, each pivot's high-root tail is no
larger than the root-$128$ unscreened tail.  Indeed $m\leq5929$, and the
four higher-root mask maxima are $986,592,315,146$; in every case
\[
 H\left(\#m+\frac{N/25+2}{2Y}
       \left(4+\frac{\#m}{3026}\right)\right)
 \leq128\left(5929+\frac{N/25+2}{2Y}
       \left(4+\frac{5929}{3026}\right)\right).
\]
The four mask maxima remain valid; their comparison is checked by
\path{scripts/verify_four_range_paper_arithmetic.py}.

The exact restricted M\"obius support in the aggregate
diagonal estimate gives, uniformly on the present interval,
\begin{equation}\label{eq:systemic-mixed-diagonal}
 \Delta_{25}^{\rm mix}<0.002980825.
\end{equation}
If three pivots lie in one mod-$4$ class which is globally concentrated in
one mod-$9$ class, the allowed union of residue classes has density
$161/225$; its prefix discrepancy modulo $900$ is less than $5$, and
\begin{equation}\label{eq:systemic-mixed-concentrated-diagonal}
 \Delta_{900}^{\rm mix}<0.003766884.
\end{equation}
The maximum in \eqref{eq:systemic-mixed-diagonal} occurs at
\[
 N=2\cdot10^9,\quad R=3809763,\quad
 \sum_{s\leq R}2^{\omega(s)}=787461
\]
on the permitted squarefree semigroup.  The factor $13$ is used in both
Pell tails.

For primes through $131$, apply degree-$3$, degree-$2$, and degree-$3$
Bonferroni bounds to the one-, two-, and three-form survivors.  A prime common
to all three forms has multiplicity $1$; a pair coincidence has multiplicity
$1$ on that pair and $2$ in the three-form union.  Concentrating all allowed
pair coincidences on one edge maximizes the sum of the three quadratic pair
bounds.  Direct expansion gives the following terminal table.  The column
$\theta$ multiplies $C_{\rm quad}$; ``charge'' is an explicitly discarded
part of $A^*$.
\begingroup\small
\[
\begin{array}{c|c|c|c|c|c}
\text{branch}&\theta&\text{common primes}&\text{charge}&(I,\epsilon)&|A|/N\\ \hline
\text{even generic}&23/25&7,11,13,17&0&(0.041080150,45066)&<0.038927972\\
\text{two odd classes}&23/50&7,11,13,17&\eta_0&(0.115095250,49865)&<0.034542160\\
\text{two odd, common }3&23/50&3,7,11,13,17&\eta_0&(0.160187242,48180)&<0.038147835\\
\text{one odd class}&23/100&2,7,11,13,17&0&(0.290282102,48180)&<0.039998152\\
\text{one odd, common }2,3&23/900&2,3&0.0012&(0.344765994,53410)&<0.039971251\\
\text{global mod }9&23/225&3,7,11,13,17&0&(0.145156038,43495)&<0.024889106\\
\text{even concentrated}&161/225&3&0&(0.113917972,50119)&<0.039955268\\
\text{mixed parity}&23/100&7,11,13,17&\eta_0&(0.261748734,49285)&<0.039984162.
\end{array}
\]
\endgroup
Here every row includes
\[
 \frac{6\tau_{131}}{25}<0.000301830
\]
and the common intermediate-prime/root payment
\begin{align*}
 &\frac6{77}\left(\frac1L+\frac1{L^2}+\frac{2.51}{L^3}\right)
 +\frac{2H_\Sigma}{N}\left(5929+\frac{156}{25}+\frac{312}{N}
 +\left(\frac{39}{25}+\frac{78}{N}\right)\frac{5929}{3026}\right)\\
 &\hspace{35mm}<0.007155053,\quad L=\log(N/78).
\end{align*}
Thus the displayed values are strict upper bounds at $N=2\cdot10^9$;
every analytic term decreases thereafter.  The exact rational finite
polynomials and restricted diagonal jump scan are reproduced by
\path{scripts/systematic_mixed_full_close.py}; all branch totals
are checked by \path{scripts/verify_four_range_paper_arithmetic.py}.

\begin{theorem}[Systemic mixed three-pivot close]
\label{thm:systemic-mixed-close}
The Hall inequality \eqref{eq:hall-defect}, and hence the required bound for
Erd\H{o}s Problem~848, holds for
\[
 2\cdot10^9\leq N<7.67769\cdot10^9.
\]
Consequently it holds for every $N\geq2\cdot10^9$.
\end{theorem}

\begin{proof}
Suppose that $B$ is a Hall defect and put $A=B\cup T_N(B)$.  Then $A$ is
compatible, $|A|>|C_N|$, and $A^*$ contains $B^*$ from
Corollary~\ref{cor:systemic-mixed-residual}.  Let $e$ be the number of even
members of $A^*$ and split at $e=\eta_0N$.

First suppose $e>\eta_0N$.  One even mod-$4$ class contains more than
$0.0011337875N$ members.  If that class is not contained in one residue
modulo $9$, choose three of its members which are nonconstant modulo $9$.
The prime $2$ is inactive and $3$ is not common.  Since a common-prime
product $Q$ satisfies $4Q^2<N$, at most four further common primes occur;
replacing them by $7,11,13,17$ only enlarges the finite survivor bound.  This
is the first row of the table.

If the even class is contained in one residue modulo $9$, choose three
consecutive members.  Their span is at most $1763$, so $3$ is the only
possible common odd prime because
$4(3\cdot7)^2=1764$.  The other three mod-$4$ classes are unrestricted,
while the selected class is restricted modulo $9$; this is precisely the
$161/225$ residue union in
\eqref{eq:systemic-mixed-concentrated-diagonal}.  The seventh row closes this
case.

Now suppose $e\leq\eta_0N$.  Charge the even part.  If the odd part meets
both odd classes modulo $4$, choose a triple nonconstant modulo $4$; if the
odd part is also nonconstant modulo $9$, the same triple may be chosen
nonconstant for both projections.  These are respectively the second and,
when all odd members share one residue modulo $9$, the third rows.  The
diagonal set is restricted to odd non-base values, giving $\theta=23/50$.

It remains that all odd members of $A^*$ occupy one mod-$4$ class.  If
$e=0$ and this class is nonconstant modulo $9$, the prime $2$ is common but
$3$ is not; the fourth row applies.  If the class is constant modulo $9$,
three consecutive odd pivots have common primes only $2$ and $3$.  When
$e\leq0.0012N$, charge the even part and use the fifth row.
When instead $e>0.0012N$ and all of $A^*$ is constant modulo $9$, select
three even pivots; the prime $2$ is inactive, and the sixth row allows the
common prime $3$ and four further common primes.

Finally suppose $0<e\leq\eta_0N$ and $A^*$ is not constant modulo $9$.
Choose two odd pivots and one even pivot so that the three residues modulo
$9$ are nonconstant.  The $2$-square event is active on at most two forms
and is not common.  The asymmetric finite polynomial in the last row,
together with the charged even part, closes this final case.

These alternatives are exhaustive.  Every row gives
\[
 \frac{|A|}{N}<\frac1{25}-\frac7{25N}\leq\frac{|C_N|}{N},
\]
contradicting $|A|>|C_N|$.  This proves the stated interval.  The range
$N\geq7.67769\cdot10^9$ is Theorem~\ref{thm:767769e9-close}, completing the
last assertion.
\end{proof}

\subsection{Actual-support descent to two hundred million}

For $N_0\leq N<N_1$ we retain the ordered prime support.  Put
\[
 N_0=2\cdot10^8,\qquad N_1=2\cdot10^9,
 \qquad D=\lfloor N/300\rfloor,\qquad M=90001.
\]
If a pivot has $k$ distinct non-$5$ odd prime divisors, retain its smallest
$k-1$ of them and require that a larger $k$-th prime still fit under $N_1$.
Exact intersections of their nonzero quadratic-residue masks on $m\leq M$
give
\[
\begin{array}{c|rrrrrrrrr}
k&0&1&2&3&4&5&6&7&8\\ \hline
s_k&90001&90001&45180&22924&11414&5520&2513&1039&353.
\end{array}
\]
The maximizing supports for $k=2,\ldots,8$ are respectively
\[
 (38639),\ (971,1429),\ (137,193,199),\ (43,61,73,89),
\]
\[
 (23,29,31,41,43),\ (11,13,17,19,23,41),
 \ (3,7,13,17,19,23,29).
\]

\begin{proposition}[Uniform actual-support degree]
\label{prop:actual-support-degree}
For $N_0\leq N<N_1$, every outside vertex has more than
\[
 \frac N{50}+0.0001574N
\]
squarefree neighbours in $C_N$, and the identical reverse bound holds from
$C_N$ into $C_N^{18}$.  Consequently every mixed Hall defect satisfies
\begin{equation}\label{eq:actual-support-residual}
 |B\setminus C_N^{18}|>0.0003148N.
\end{equation}
\end{proposition}

\begin{proof}
Since $D\geq(N-299)/300$, throughout the interval
\[
 \frac{N^2+25N+1}{D^2}<90001.
\]
For a valuation factor $c_\nu\in\{1,2,4\}$ and $k$ odd prime
divisors, put $H=c_\nu2^{k+1}$.  The same M\"obius split as in
Proposition~\ref{prop:systemic-mixed-degree}, now with the actual-support
survivor $s_k$, gives
\begin{align*}
 \deg(b)-\frac N{50}\geq{}&
 \left(\frac43\right)^{[2\mid b]}
 \frac{\kappa_5(N-43)}{25}-\frac N{300}-\frac N{50}-1\\
 &-\frac{N/25+1}{(N-299)/300}
 -\frac1{5(N-299)/300}\\
 &-H\left(s_k+\frac{X_D}{2}
       \left(4+\frac{s_k}{3026}\right)\right),
 \qquad X_D=\frac{N/25+1}{(N-299)/300}.
\end{align*}
Every fixed row divided by $N$ is increasing: its leading coefficient is
positive, while each remaining positive payment divided by $N$ decreases.
It therefore suffices to use $N=N_0$ and the support superset at $N_1$.
Among all $4\cdot9$ rows the minimum is the odd $k=3$, $H=16$ row,
with survivor $22924$ and witness $(971,1429)$; its margin is
\[
 31482.2990360\ldots>0.0001574N_0.
\]
The reverse progression has the same support and root calculation.  The
argument of Corollary~\ref{cor:systemic-mixed-residual} then gives
$|B\setminus C_N^{18}|>2(0.0001574N)-1$, proving
\eqref{eq:actual-support-residual}.  The exact support enumeration is
reproduced by \path{scripts/systematic_mixed_interval_support.py}; the
degree rows are checked by
\path{scripts/verify_four_range_paper_arithmetic.py}.
\end{proof}

The transformed congruence contains one further uniform condition which was
not needed for Proposition~\ref{prop:actual-support-degree}.  In the notation
of the progression split, $d_1=25b/g$ with $g\in\{1,5\}$, so $5\mid d_1$;
moreover $c_1$ is a unit.  Hence
\[
 d^2m\equiv c_1\pmod {d_1}
\]
forces $m$ into one fixed nonzero square coset modulo $5$.  Taking the worse
of the two possible cosets and intersecting it with the same actual-support
masks gives
\[
\begin{array}{c|rrrrrrrrr}
k&0&1&2&3&4&5&6&7&8\\ \hline
s_k^{(5)}&36001&36001&18234&9314&4680&2290&1055&471&175.
\end{array}
\]
The maximizing retained supports for $k=2,\ldots,8$ are respectively
\[
 (44563),\ (1087,1117),\ (101,241,263),\ (53,59,71,83),
\]
\[
 (19,29,37,41,43),\ (11,13,17,19,29,31),
 \ (3,7,11,13,17,23,37).
\]

\begin{corollary}[Mod-$5$-coset actual-support degree]
\label{cor:mod5-actual-support-degree}
On $N_0\leq N<N_1$, every forward and reverse outside degree is greater
than
\[
 \frac N{50}+0.0012427N.
\]
Consequently every mixed Hall defect in this interval satisfies
\[
 |B\setminus C_N^{18}|>0.0024854N-1>0.002485N.
\]
\end{corollary}

\begin{proof}
Repeat Proposition~\ref{prop:actual-support-degree} with $s_k^{(5)}$.
The unique controlling row is now the odd $k=4$ row, with $H=32$,
survivor $4680$, and retained support $(101,241,263)$.  At $N=N_0$ its
margin above $N/50$ is
\[
 248552.6176834\ldots>0.0012427N_0.
\]
The same normalized monotonicity proves the degree bound on the full
interval.  Applying the forward--reverse defect argument gives the displayed
residual.  Both cosets and every feasible support prefix at the upper endpoint
are enumerated by \path{scripts/systematic_mod5_actual_support.py}; the
degree rows are checked by
\path{scripts/verify_four_range_paper_arithmetic.py}.
\end{proof}


The same support constraint is retained in the large-square tail.  On the
following four blocks use $Y=\lfloor N/k_0\rfloor$.  The integer box is
$m\leq k_0^2$: indeed
$(k_0^2+1)(N/k_0-1)^2>N^2+25N+1$ at the left endpoint, and the
difference increases.  Exact support intersections give the displayed
worst root row and the resulting three-pivot prime/root ceiling; the
intermediate-prime term uses the three-term Dusart bound and
$Y\geq N/(k_0+1)$.
\[
\begin{array}{c|c|c|c|c|c}
[L,U)&k_0&H&s&\text{support}&\text{three-pivot ceiling}\\ \hline
[2,3]\cdot10^8&75&128&805&(97,113,151)&0.008955907\\
{}[3,5]\cdot10^8&85&128&1016&(97,113,151)&0.007668494\\
{}[5,10]\cdot10^8&100&128&1391&(73,199,257)&0.006336343\\
{}[10,20]\cdot10^8&125&128&2124&(47,233,401)&0.004881897.
\end{array}
\]
All support maxima are reproduced by
\path{scripts/systematic_mod5_actual_support.py}.  The monotone
analytic payments are checked by
\path{scripts/verify_four_range_paper_arithmetic.py}.

\begin{theorem}[Actual-support pure-$18$ parity matching]
\label{thm:actual-support-pure-matching}
For every $N\geq2\cdot10^8$, the squarefree-edge graph between
$C_N^{18}$ and $C_N$ has a matching saturating $C_N^{18}$.
Consequently every $B\subseteq C_N^{18}$ satisfies the original Hall
inequality.
\end{theorem}

\begin{proof}
It remains only to treat $N_0\leq N<N_1$.  Work in the parity blocks
$L_i$--$R_i$ used in Theorem~\ref{thm:systemic-pure-parity-matching}.
The three-pivot ceiling above is six times a one-form intermediate-prime plus
root payment.  A parity-block progression has length $N/50+1$, rather than
$N/25+1$, so one sixth of that ceiling is a valid one-form bound.  Add the
square-density tail
$(N/50+1)0.000407315$ and retain the exact finite constants of
Lemma~\ref{lem:zero-coincidence-triple-block}.  At the four left endpoints,
the degree lower bound $D$ and common-non-neighbour upper bound $U$ satisfy
\[
\begin{array}{c|c|c|c}
N&D/N&U/N&D-U\\ \hline
2\cdot10^8&0.0152620317&0.0127212769&>508150\\
3\cdot10^8&0.0154767174&0.0093626504&>1834220\\
5\cdot10^8&0.0156988362&0.0065246389&>4587098\\
10^9&0.0159413140&0.0041684639&>11772850.
\end{array}
\]
On each block $D/N$ increases and $U/N$ decreases.  In particular
$D>N/100$ and $U<D$.  A Hall defect in either parity block therefore
contains three consecutive pivots of span at most $200$; the reverse degree
argument forces its common non-neighbourhood to exceed $D$, contradicting
$U<D$, exactly as in Theorem~\ref{thm:systemic-pure-parity-matching}.
The four endpoint calculations are reproduced by
\path{scripts/verify_four_range_paper_arithmetic.py}.  The already proved range
$N\geq N_1$ completes the assertion.
\end{proof}

We finally replace the absolute small-square rounding loss by an exact
finite prefix.  Let $R=2\cdot10^7$.  For a union $W$ of non-base residue
classes, put
\begin{align*}
 Q_W(U)={}&-\sum_{\substack{2\leq s\leq R\\
                p\mid s\Rightarrow p\equiv1\ (4)}}
 \mu(s)\#\{x\leq U:x\in W,\ s^2\mid x^2+1\}\\
 &+\frac{2\theta_WU(\log R+2)}R
 +13\frac{U^2+1}{R^2},
\end{align*}
where $\theta_W$ is the density of $W$.  M\"obius inversion, the harmonic
tail, and the sharp negative-Pell orbit bound show that $Q_W(U)$ is an upper
bound for the diagonal non-squarefree count through $U$.

Partition $[N_0,N_1)$ into 463 adjacent blocks $[L,U]$, each of relative
width less than $0.005$, with additional breaks at
$3\cdot10^8,5\cdot10^8,10^9$.  Since the actual diagonal count is
monotone, $Q_W(U)/L$ bounds its normalized value throughout the block.
The finite computation enumerates $1237599$ squarefree support nodes,
$4134027$ roots, and $51357374$ progression events.  The eighteen
concentrated sets are evaluated simultaneously from
\[
 U_{25}-E_e+E_{e,c}\qquad(e=0,2;\ c\bmod9).
\]
The controlling block is $[200000000,200999999]$, where the worst
concentrated diagonal is
\begin{equation}\label{eq:actual-support-diagonal}
 \frac{Q_W(U)}L<0.019673927941605.
\end{equation}
All integer events and block maxima are reproduced by
\path{scripts/systematic_concentrated_diagonal_blocks.cpp}.

\begin{theorem}[Actual-support mixed close]
\label{thm:actual-support-mixed-close}
The Hall inequality \eqref{eq:hall-defect}, and hence the exact bound for
Erd\H{o}s Problem~848, holds for every
\[
 N\geq2\cdot10^8.
\]
\end{theorem}

\begin{proof}
The range $N\geq N_1$ is Theorem~\ref{thm:systemic-mixed-close} and the
subsequent descent, so let $N_0\leq N<N_1$.  Suppose $B$ is a Hall defect
and put $A=B\cup T_N(B)$.  Theorem~\ref{thm:actual-support-pure-matching}
excludes $B\subseteq C_N^{18}$.  Corollary~
\ref{cor:mod5-actual-support-degree} therefore gives
\[
 |A^*|\geq|B\setminus C_N^{18}|>0.002485N.
\]
Let $e$ be the number of even members of $A^*$ and put
$\eta=0.00016$.

If $e>\eta N$, one even mod-$4$ class has density greater than
$\eta/2$.  A closest pair in it has gap less than
$2/\eta=12500$.  If the class is nonconstant modulo $9$, add a third pivot
which makes the triple nonconstant modulo $9$; the prime $2$ is inactive and
$3$ is not common.  If it is constant modulo $9$, the same short pair shows
that the only possible common odd primes are $3,7$, since
\[
 4(3\cdot7\cdot11)^2=213444>12500.
\]
These are the even-generic and even-concentrated rows below.

Now suppose $e\leq\eta N$ and charge all even members.  More than
$0.002325N$ odd members remain.  If both odd mod-$4$ classes occur, choose
a triple nonconstant modulo $4$ and split according as its odd support is
nonconstant or constant modulo $9$.  If only one odd mod-$4$ class occurs,
a closest pair has gap less than $6169.4$.  In the nonconstant-mod-$9$ case
add a third pivot of another mod-$9$ residue: $2$ is common, $3$ is not, and
at most one further odd prime is common because
$4(7\cdot11)^2=23716>6169.4$.  In the constant case the common list is at
most $2,3,7$, since adjoining $11$ would again force a gap at least
$4(3\cdot7\cdot11)^2$.

Use cutoff $131$, degree-$3/2/3$ Bonferroni bounds, the integrated tail
$6\tau_{131}/25$, the appropriate exact diagonal prefix, and the root
ceiling of the containing product block.  Exhausting all 463 blocks gives
\[
\begin{array}{c|c|c}
\text{branch}&\text{allowed common primes}&\max |A|/N\\ \hline
\text{even generic}&7,11,13,17&0.038287487920\\
\text{even concentrated}&3,7&0.039870783152\\
\text{two odd classes}&7,11,13,17&0.031773851295\\
\text{two odd, common }3&3,7,11,13,17&0.035364360630\\
\text{one odd class}&2,7&0.038964611892\\
\text{one odd, common }2,3&2,3,7&0.039088651986.
\end{array}
\]
Every maximum occurs in the first block.  Its exact target is
$1/25-7/(25N)>0.0399999986$, so all six exhaustive alternatives contradict
$|A|>|C_N|$.  This proves the new interval and hence the theorem.
\end{proof}

\begin{corollary}[A mixed-base strip]\label{cor:thin-nonbase-strip}
Let $N\geq7.67769\cdot10^9$ and let $B$ be a compatible outside set.  Put
\[
 B^*=B\setminus C_N^{18},\qquad
 \mathcal R(N)=2\mathcal D(N)-|C_N^{18}|.
\]
If $|B^*|\leq\mathcal R(N)$, then $B$ satisfies the Hall inequality.  Moreover
\[
 \mathcal R(N)>0.000000036N
 \qquad(2.2827\cdot10^{10}\leq N<3.4375\cdot10^{11}),
\]
and, more sharply,
\[
 \mathcal R(N)>0.0010038N
 \qquad(7.67769\cdot10^9\leq N<1.15\cdot10^{10}),
\]
and
\[
 \mathcal R(N)>0.0000268N
 \qquad(1.15\cdot10^{10}\leq N<1.2643\cdot10^{10}),
\]
and
\[
 \mathcal R(N)>0.000000228N
 \qquad(1.2643\cdot10^{10}\leq N<1.279\cdot10^{10}),
\]
and
\[
 \mathcal R(N)>0.00000079N
 \qquad(1.279\cdot10^{10}\leq N<1.29\cdot10^{10}),
\]
and
\[
 \mathcal R(N)>0.00000104N
 \qquad(1.29\cdot10^{10}\leq N<1.382\cdot10^{10}),
\]
and
\[
 \mathcal R(N)>0.00000236N
 \qquad(1.382\cdot10^{10}\leq N<2.2827\cdot10^{10}),
\]
and
\[
 \mathcal R(N)>0.000127N
 \qquad(2.3669\cdot10^{10}\leq N<3.4375\cdot10^{11}),
\]
and
\[
 \mathcal R(N)>0.001162N
 \qquad(3.22715\cdot10^{10}\leq N<3.4375\cdot10^{11}),
\]
and
\[
 \mathcal R(N)>0.001187N
 \qquad(3.2532\cdot10^{10}\leq N<3.4375\cdot10^{11}),
\]
and
\[
 \mathcal R(N)>0.001414N
 \qquad(3.4984\cdot10^{10}\leq N<3.4375\cdot10^{11}),
\]
and
\[
 \mathcal R(N)>0.001424N
 \qquad(3.51\cdot10^{10}\leq N<3.4375\cdot10^{11}),
\]
and
\[
 \mathcal R(N)>0.00032N
 \qquad(1.001\cdot10^{11}\leq N<N_{25}),
\]
and
\[
 \mathcal R(N)>0.000724N
 \qquad(N\geq5.5\cdot10^{11}).
\]
Thus every Hall defect in the new bottom interval must contain more than
$0.001627N$ genuinely non-base outside vertices.  The subsequent displayed
bounds give the strip separately on each later degree piece.
\end{corollary}

\begin{proof}
If $B^*=\varnothing$, use Corollary~\ref{cor:opposite-base-matching}.  Otherwise
fix $x\in B^*$.  Repeat the relevant one-pivot degree proof with the variable in the progression
$18\pmod {25}$.  Since $x\not\equiv18\pmod {25}$, the gcd
$(18x+1,25x)$ is $1$ or $5$, so the identical calculation gives at least
$\mathcal D(N)$ values $y\in C_N^{18}$ for which $xy+1$ is squarefree.
Compatibility forces all of them out of $B$.  Consequently
\[
 |B\cap C_N^{18}|\leq|C_N^{18}|-\mathcal D(N).
\]
Under the stated hypothesis,
\[
 |B|\leq |C_N^{18}|-\mathcal D(N)+|B^*|
 \leq\mathcal D(N).
\]
The corresponding one-pivot degree bound now proves the Hall inequality.

Finally $|C_N^{18}|\leq N/25+1$.  Direct substitution in the displayed
formula for $\mathcal M_{\rm QR2}$ at $N_0=7.67769\cdot10^9$ gives
\[
 \frac{2\mathcal M_{\rm QR2}(N_0)}{N_0}-\frac1{25}-\frac1{N_0}
 >0.0010038.
\]
Direct substitution in $\mathcal M_{\rm QR}$ at $1.15\cdot10^{10}$ gives
\[
 \frac{2\mathcal M_{\rm QR}(1.15\cdot10^{10})}{1.15\cdot10^{10}}
 -\frac1{25}-\frac1{1.15\cdot10^{10}}
 >0.0000268.
\]
Direct substitution in $\mathcal M_{1/32}$ at $1.2643\cdot10^{10}$ gives
\[
 \frac{2\mathcal M_{1/32}(1.2643\cdot10^{10})}{1.2643\cdot10^{10}}
 -\frac1{25}-\frac1{1.2643\cdot10^{10}}
 >0.000000228.
\]
Direct substitution in $\mathcal M_{1/16}$ at $1.279\cdot10^{10}$ gives
the stronger $0.00000079$ bound.
Direct substitution in $\mathcal M_{1/8}$ at $1.29\cdot10^{10}$ gives the
stronger $0.00000104$ bound.  Direct substitution in $\mathcal M_{1/4}$ at
$1.382\cdot10^{10}$ gives the
stronger $0.00000236$ bound.  Direct substitution in $\mathcal M_{1/2}$ at
$2.2827\cdot10^{10}$ gives
the weaker transition bound $0.000000036$, while substitution at
$2.3669\cdot10^{10}$ gives the sharper $0.000127$
bound.  Direct substitutions at $3.22715\cdot10^{10}$, $3.2532\cdot10^{10}$,
$3.4984\cdot10^{10}$, and $3.51\cdot10^{10}$ give the sharper
$0.001162$, $0.001187$, $0.001414$, and $0.001424$ bounds.
The quotient increases on each of the displayed degree
pieces.  Direct checks from $\mathcal M_{1/2}$ to $\mathcal M_1$ preserve
both displayed bounds on their stated intervals.  The later switches to $\mathcal M$, across its root-count jump, and
from $\mathcal M$ to $\mathcal L$ keep the stated weaker bounds.  Direct
evaluation at $1.001\cdot10^{11}$ gives more than
$0.00032N$, and at $5.5\cdot10^{11}$ gives more than $0.0012N$,
and the same switch checks prove the $0.000724N$ bound from that point onward.
\end{proof}

\section{The valid two-pivot CRT decomposition}

The large-prime tail is nonempty in general, but it can be placed into exact
CRT classes without changing the Hall target.

\begin{lemma}[Two-pivot square classes]\label{lem:two-pivot}
Let $b,c\not\equiv7\pmod {25}$ be distinct, and let
$a\equiv7\pmod {25}$.  Suppose both $ab+1$ and $ac+1$ are non-squarefree.
Choose primes $p,q$ with
\[
  p^2\mid ab+1,\qquad q^2\mid ac+1.
\]
Then $p,q\neq5$, and exactly one of the following alternatives applies:
\begin{enumerate}
\item $p=q$ and $p^2\mid b-c$;
\item $p\neq q$ and $a$ lies in one uniquely determined residue class modulo
      $25p^2q^2$.
\end{enumerate}
\end{lemma}

\begin{proof}
If $25\mid ab+1$, then $7b+1\equiv0\pmod {25}$.  Since
$7^{-1}\equiv18\pmod {25}$, this forces $b\equiv7\pmod {25}$, contrary to
the hypothesis.  Hence $p\neq5$, and similarly $q\neq5$.

If $p=q$, subtracting the two divisibilities gives
$p^2\mid a(b-c)$.  But $p\nmid a$, since $p\mid a$ would imply
$ab+1\equiv1\pmod p$.  Therefore $p^2\mid b-c$.

If $p\neq q$, the three congruences
\[
 a\equiv7\pmod {25},\qquad
 a\equiv-b^{-1}\pmod {p^2},\qquad
 a\equiv-c^{-1}\pmod {q^2}
\]
have pairwise coprime moduli.  The Chinese remainder theorem gives one class
modulo $25p^2q^2$.
\end{proof}

For a compatible pair $B=\{b,c\}$ define the same-prime and distinct-prime
parts of its common non-neighbour set by
\begin{align*}
 S_N(b,c)&=\{a\in\TN(B):\text{some prime }p\text{ has }
              p^2\mid ab+1\text{ and }p^2\mid ac+1\},\\
 D_N(b,c)&=\TN(B)\setminus S_N(b,c).
\end{align*}
This is an exact disjoint decomposition, not a new cut.  The same-prime debit
can be bounded independently of what occurs in $D_N(b,c)$.

\begin{proposition}[Two-point same-prime debit]
\label{prop:same-prime-debit}
Let $N\geq107$ and let $B=\{b,c\}$ be a compatible outside set.  Then
\[
 2+|S_N(b,c)|\leq|\CN|.
\]
For $N\geq500$ the stronger explicit estimate
\begin{equation}\label{eq:distinct-residual-debit}
 |\CN|-2-|S_N(b,c)|>
 \frac{0.549N}{25}-\log_4(N-1)-\frac{57}{25}
\end{equation}
holds.  Consequently the two-point Hall inequality follows whenever the
distinct-prime part fits inside the right-hand side of
\eqref{eq:distinct-residual-debit}.
\end{proposition}

\begin{proof}
By Lemma~\ref{lem:two-pivot}, every relevant prime
satisfies $p\neq5$ and $p^2\mid|b-c|$.  For a fixed $p$, the simultaneous
conditions $a\equiv7\pmod {25}$ and $p^2\mid ab+1$ determine one residue
class modulo $25p^2$.  Hence
\begin{equation}\label{eq:same-prime-count}
 |S_N(b,c)|\leq
 \sum_{\substack{p^2\mid|b-c|\\p\neq5}}
 \left(\frac{N}{25p^2}+1\right).
\end{equation}

First suppose $107\leq N\leq499$.  Since
$(2\cdot3\cdot7)^2=1764>N-1$, at most two distinct primes occur in the sum,
and the largest two class bounds are those for $p=2,3$.  For
$107\leq N\leq167$ the exact class counts are at most $2$ and $1$, so
$2+|S_N(b,c)|\leq5\leq|\CN|$.  For $168\leq N\leq499$,
\[
 2+|S_N(b,c)|\leq4+\frac N{100}+\frac N{225}
 =4+\frac{13N}{900}
 \leq\frac N{25}-\frac7{25}\leq|\CN|.
\]
The middle inequality holds already at $N=168$ and its right-minus-left
side increases with $N$.

Now let $N\geq500$, and let $k$ be the number of primes in
\eqref{eq:same-prime-count}.  Their squared product divides $|b-c|\leq N-1$,
so $4^k\leq N-1$ and $k\leq\log_4(N-1)$.  Also, every relevant prime is
different from $5$.  Removing $5$ and the displayed composite integers from
the sum over all integers gives the elementary numerical bound
\[
 \sum_{p\neq5}\frac1{p^2}
 <\frac{\pi^2}{6}-1-
 \sum_{m\in\{4,5,6,8,9,10,12,14,15,16,18,20\}}\frac1{m^2}
 <0.451.
\]
It follows that
\[
 2+|S_N(b,c)|<2+\frac{0.451N}{25}+\log_4(N-1)
 \leq\frac N{25}-\frac7{25}\leq|\CN|.
\]
At $N=500$ the middle inequality has positive margin, and its margin is
increasing because
\[
 \frac{0.549}{25}-\frac1{(N-1)\log4}>0
 \qquad(N\geq500).
\]
This closes every case.  Subtracting the last strict upper bound from
$|\CN|\geq N/25-7/25$ gives
\eqref{eq:distinct-residual-debit}.
\end{proof}

Thus the same-prime contribution has a global debit bound, while every fixed
finite distinct-prime skeleton is impossible.  What remains is to fit the
genuinely unbounded $D_N(b,c)$ contribution into the residual debit
\eqref{eq:distinct-residual-debit}.

The extreme one-form tail already has a coefficient-uniform square-root
bound.  This pays the branch in which only one witness is very large, without
requiring a second large witness.

\begin{lemma}[Uniform one-form large-square tail]
\label{lem:one-form-large-square-tail}
For $1\leq t\leq N$ and $Y\geq1$, put
\[
 \mathcal H_t(Y)=
 \{a\leq N:\text{some prime }p>Y\text{ satisfies }p^2\mid at+1\}
\]
and
\[
 K_t=\left\lfloor\frac{tN+1}{Y^2}\right\rfloor.
\]
Then
\begin{equation}\label{eq:one-form-large-square-tail}
 |\mathcal H_t(Y)|
 \leq
 2\sqrt{N+1}\sum_{s=1}^{K_t}\frac1{\sqrt s}
 +2K_t\sqrt t
 \leq4\sqrt{K_t(N+1)}+2K_t\sqrt t.
\end{equation}
The same bounds hold after restricting $a$ to any progression.
\end{lemma}

\begin{proof}
First note the elementary root bound
\begin{equation}\label{eq:unit-square-root-count}
 \#\{x\pmod m:x^2\equiv u\pmod m\}\leq2\sqrt m
 \qquad ((u,m)=1).
\end{equation}
Indeed, an odd prime power contributes at most two roots, while a power of
$2$ contributes $1,2$, or at most $4$ roots according as its exponent is
$1,2$, or at least $3$.  If the odd part of $m$ has $r$ distinct prime
divisors, then
\[
 \frac{2^r}{\sqrt{\operatorname{rad}(m_{\rm odd})}}
 \leq\max\left\{1,\frac2{\sqrt3}\right\};
\]
after the possible prime $3$, every new odd prime contributes a factor at
most $2/\sqrt5<1$.  Combining this with the three stated $2$-adic cases gives
\eqref{eq:unit-square-root-count}.  The case $m=1$ is immediate.

For every counted $a$, choose a witnessing prime and write
\[
 p^2s=at+1.
\]
Then $1\leq s\leq K_t$, $(s,t)=1$, and
\[
 p^2\equiv s^{-1}\pmod t,
 \qquad p\leq\sqrt{\frac{tN+1}{s}}.
\]
For fixed $s$, \eqref{eq:unit-square-root-count} bounds the number of possible
integers $p$---and hence the number of possible $a$---by
\[
 2\sqrt t\left(
 \frac1t\sqrt{\frac{tN+1}{s}}+1
 \right)
 \leq2\sqrt{\frac{N+1}{s}}+2\sqrt t.
\]
Summing over $s$ gives the first bound.  The second follows from
$\sum_{s\leq K}s^{-1/2}\leq2\sqrt K$.  Forgetting primality and progression
conditions only enlarges the count.
\end{proof}

\begin{corollary}[The very-large mixed tail is paid]
\label{cor:very-large-mixed-tail}
Let $N\geq41$ and $Y=\lfloor N/3\rfloor$.  For every $t\leq N$,
\[
 |\mathcal H_t(Y)|<27.42\sqrt{N+1}.
\]
Consequently, for any three pivots $t_1,t_2,t_3\leq N$, the number of
$a\leq N$ for which at least one of the three forms $at_i+1$ has a
prime-square witness exceeding $Y$ is less than
\begin{equation}\label{eq:three-form-very-large-tail}
 82.26\sqrt{N+1}.
\end{equation}
This includes, in particular, the case in which exactly one witness exceeds
$Y$.
\end{corollary}

\begin{proof}
The inequality $N^2+1<10\lfloor N/3\rfloor^2$ holds for $N\geq41$, so
$K_t\leq9$.  Since
\[
 \sum_{s=1}^9s^{-1/2}<4.705,
\]
the first bound in \eqref{eq:one-form-large-square-tail} is less than
$9.41\sqrt{N+1}+18\sqrt N<27.42\sqrt{N+1}$.  The union bound over the three
forms gives \eqref{eq:three-form-very-large-tail}.
\end{proof}

The root count and tail arithmetic are reproduced by
\path{scripts/uniform_large_square_tail_scan.py}.  On its displayed finite box
it checks $76116$ unit-congruence cases and $173880$ tail cases, with no root,
tail, or $K_t\leq9$ failure.

The large-square part of that obligation admits a coefficient-uniform Pell
reduction.  The point is to retain the determinant $|b-c|$ instead of
multiplying one coefficient into the norm, which would introduce a spurious
root count depending on the growing pivots.

\begin{lemma}[Determinant-uniform Pell orbits]
\label{lem:determinant-pell-orbits}
Let $A,B,d$ be positive integers.  For $X\geq\sqrt d$, let
$\mathcal P_{A,B,d}(X)$ be the set of positive integer solutions of
\[
 Ap^2-Bq^2=d,\qquad (q,d)=1,
 \qquad p\sqrt A+q\sqrt B\leq X.
\]
Put
\[
 R_d(A,B)=
 \#\{z\pmod d:Az^2\equiv B\pmod d\},
\]
with $R_1(A,B)=1$.  Then
\begin{equation}\label{eq:determinant-pell-orbit-bound}
 |\mathcal P_{A,B,d}(X)|
 \leq R_d(A,B)
 \left(1+\left\lfloor
 \log_{2+\sqrt3}\frac X{\sqrt d}
 \right\rfloor\right)
 \leq d\left(1+\left\lfloor
 \log_{2+\sqrt3}\frac X{\sqrt d}
 \right\rfloor\right).
\end{equation}
In particular, the orbit count depends on the determinant $d$, not on the
possibly growing product $AB$.
\end{lemma}

\begin{proof}
Every solution determines
\[
 z\equiv pq^{-1}\pmod d,
 \qquad Az^2\equiv B\pmod d.
\]
Take two solutions $(p_1,q_1)$ and $(p_2,q_2)$ with the same $z$, and put
\[
 u=\frac{Ap_1p_2-Bq_1q_2}{d},
 \qquad
 v=\frac{p_1q_2-p_2q_1}{d}.
\]
The common-root congruence shows that both numerators are divisible by $d$.
The identity
\[
 (Ap_1p_2-Bq_1q_2)^2
 -AB(p_1q_2-p_2q_1)^2
 =(Ap_1^2-Bq_1^2)(Ap_2^2-Bq_2^2)=d^2
\]
therefore gives
\[
 u^2-ABv^2=1.
\]
Writing $\alpha_i=p_i\sqrt A+q_i\sqrt B$ and
$\beta_i=p_i\sqrt A-q_i\sqrt B$, one also has
\[
 \frac{\alpha_1}{\alpha_2}=u-v\sqrt{AB}.
\]
Here $\alpha_i>0$, $\beta_i=d/\alpha_i>0$, and the conjugate of the displayed
quotient in $\mathbb Q(\sqrt{AB})$ is $\beta_1/\beta_2$.  Hence, after
ordering the two solutions by $\alpha_i$, the larger of a nontrivial quotient
and its inverse is a totally positive unit greater than $1$.  If $AB$ is a
square, the displayed unit equation forces $v=0$ and $u=1$, so the two
solutions coincide.  Otherwise every nontrivial integral totally positive
unit greater than $1$ is at least $2+\sqrt3$: indeed it has the form
$x+y\sqrt{AB}$ with $x\geq2$, $y\geq1$, and
$AB y^2=x^2-1\geq3$.

Thus the solutions in one root class form a geometric sequence of ratio at
least $2+\sqrt3$.  Since $\alpha_i$ has norm $d$ and is larger than its
positive conjugate, $\alpha_i\geq\sqrt d$.  Each root class consequently
contains at most the first parenthesis in
\eqref{eq:determinant-pell-orbit-bound}.  Summing over the $R_d(A,B)\leq d$
root classes proves the claim.
\end{proof}

\begin{corollary}[A one-payment two-large-witness tail]
\label{cor:two-large-witness-tail}
Let $1\leq b<c\leq N$, put $d=c-b$, and let $Y\geq d$.  Define
\[
 K=\left\lfloor\frac{N^2+1}{Y^2}\right\rfloor,
 \qquad
 L=1+\left\lfloor
 \log_{2+\sqrt3}
 \left(2\sqrt{\frac{N(N^2+1)}d}\right)
 \right\rfloor.
\]
The number of $a\leq N$ for which there are primes $p,q>Y$ satisfying
\[
 p^2\mid ab+1,\qquad q^2\mid ac+1
\]
is at most
\begin{equation}\label{eq:two-large-witness-tail}
 K^2dL.
\end{equation}
The same bound holds after interchanging $b$ and $c$, and after restricting
$a$ to any progression.
\end{corollary}

\begin{proof}
For each counted $a$, choose one pair of witnessing primes and write
\[
ab+1=p^2r,\qquad ac+1=q^2s.
\]
Then $1\leq r,s\leq K$, while
\[
 c p^2r-bq^2s=c(ab+1)-b(ac+1)=d.
\]
Since $q>Y\geq d$, one has $(q,d)=1$.  For fixed $(r,s)$ apply
Lemma~\ref{lem:determinant-pell-orbits} with $A=cr$ and $B=bs$.  Its height
parameter may be taken to be $2\sqrt{N(N^2+1)}$, because
\[
 p\sqrt{cr}=\sqrt{c(ab+1)}\leq\sqrt{N(N^2+1)}
\]
and the analogous inequality holds for $q\sqrt{bs}$.  Thus each of the
$K^2$ pairs $(r,s)$ contributes at most $dL$ solutions.  Forgetting primality
and progression conditions only enlarges the count, proving
\eqref{eq:two-large-witness-tail}.
\end{proof}

The factor $d$ in \eqref{eq:two-large-witness-tail} was paid separately for
every quotient pair.  For an actual-support argument the quotient sets are
sparse, and the root classes should instead be summed before the Pell-orbit
payment.

\begin{lemma}[Determinant square-sieve over quotient supports]
\label{lem:determinant-square-sieve}
Let $1\leq b<c\leq N$, $d=c-b$, and $Y\geq d$.  Let
$\gamma_b,\gamma_c$ be positive integers and let
\[
 \mathcal R\subseteq[1,K_b],\qquad
 \mathcal S\subseteq[1,K_c].
\]
Consider the integers $a\leq N$ for which there are primes $p,q>Y$ and
$r\in\mathcal R,s\in\mathcal S$ satisfying
\[
ab+1=\gamma_b p^2r,\qquad ac+1=\gamma_c q^2s.
\]
Put
\[
 A_0=c\gamma_b,\qquad B_0=b\gamma_c,\qquad
 g_A=(A_0,d),\qquad g_B=(B_0,d)
\]
and
\[
 L=1+\left\lfloor
 \log_{2+\sqrt3}
 \left(2\sqrt{\frac{N(N^2+1)}d}\right)
 \right\rfloor.
\]
Then their number is at most
\begin{align}
 L\sum_{r\in\mathcal R}\sum_{s\in\mathcal S}
 R_d^\times(A_0r,B_0s)
 \ \leq\ L\min\{&d|\mathcal R||\mathcal S|,\notag\\
 &|\mathcal R|(g_BK_c+d),
 |\mathcal S|(g_AK_b+d)\}.
 \label{eq:determinant-square-sieve}
\end{align}
Here $R_d^\times(A,B)$ counts only roots $z\pmod d$ with $(z,d)=1$.
This remains valid after restricting $a$ to any progression.  In the
transformed $q_0=25$ progression one may take $\gamma_t\in\{1,5\}$ and let
$\mathcal R,\mathcal S$ be the exact actual-support/mod-$5$ quotient masks.
\end{lemma}

\begin{proof}
For a counted $a$ choose one witness pair.  Eliminating $a$ gives
\[
 A_0rp^2-B_0sq^2=d.
\]
Since $p,q>Y\geq d$, one has $(p,d)=(q,d)=1$.  For fixed $(r,s)$,
Lemma~\ref{lem:determinant-pell-orbits} therefore gives at most
$R_d^\times(A_0r,B_0s)L$ solutions, since the associated root
$z=pq^{-1}$ is a unit.  Its height bound is unchanged because
\[
 p\sqrt{A_0r}=\sqrt{c(ab+1)},\qquad
 q\sqrt{B_0s}=\sqrt{b(ac+1)}.
\]
Summing first proves the first inequality in
\eqref{eq:determinant-square-sieve}.

For the second, interchange the quotient and root sums.  There are at most
$d$ possible unit roots $z$, and a root contributes precisely when
\[
 B_0s\equiv A_0rz^2\pmod d.
\]
For fixed $(z,r)$ this linear congruence is either insoluble or is one class
modulo $d/g_B$.  Hence among $s\in\mathcal S\subseteq[1,K_c]$ it has at most
\[
 \min\left\{|\mathcal S|,\frac{g_BK_c}{d}+1\right\}
\]
solutions.  Summing over $z$ and $r$ gives both
$d|\mathcal R||\mathcal S|$ and
$|\mathcal R|(g_BK_c+d)$.  Interchanging the two sides gives the third
bound because $z^2$ is a unit, so $(A_0z^2,d)=g_A$.  Taking their minimum
proves the assertion.  Discarding a progression
condition only enlarges the count.
\end{proof}

For
$Y=\lfloor N/3\rfloor$, Corollary~\ref{cor:very-large-mixed-tail} also pays
the branch in which exactly one witness exceeds $Y$.  The remaining
large-prime obligation is therefore the medium window between the finite
sieve cutoff and $N/3$, including its endpoint discrepancy; it is still part
of the same distinct-prime debit.

For the pure $18$-class parity blocks, the determinant lemma has the following
finite very-large specialization.

\begin{corollary}[Pure parity determinant debit at the half split]
\label{cor:pure-half-determinant-debit}
Let $107\leq N\leq2\cdot10^8$, put $Y=\lfloor N/2\rfloor$, and let
$t_1<t_2<t_3$ be members of $C_N^{18}$ of one parity and of span at most
$200$.  The number of members $a$ of the opposite base parity for which at
least two of the three forms $at_i+1$ have prime-square witnesses greater
than $Y$ is at most
\begin{equation}\label{eq:pure-half-determinant-debit}
 16\left(1+\left\lfloor
 \log_{2+\sqrt3}\left(2\sqrt{\frac{N(N^2+1)}{50}}\right)
 \right\rfloor\right).
\end{equation}
At $N=2\cdot10^6$ this is at most $256$; retaining the actual pivot supports
lowers the exhaustive maximum to $224$.
\end{corollary}

\begin{proof}
For a pure pivot $t\equiv18\pmod {25}$ one has
\[
 at+1=p^2r,
 \qquad r\equiv2\text{ or }3\pmod5.
\]
The product $at$ is even in the opposite-parity block, so $at+1$ and hence
$r$ are odd.  Since $p>Y$ and $N\geq107$, one has
$(tN+1)/p^2<5$.  Thus the only possible quotient is $r=3$; every prime
divisor of $t$ may only delete that quotient by its quadratic-residue
condition.

Three same-parity members of the $18$ class and span at most $200$ have one
of the six offset patterns in the first column below.  For each pair apply
Lemma~\ref{lem:determinant-square-sieve} with its exact quotient set contained
in $\{3\}$.  Summing the unit root classes over the three pairs, and then
maximizing over the finitely many residues $t_1\equiv18\pmod {25}$ modulo
the three gaps, gives
\[
\begin{array}{c|rrrrrr}
(t_2-t_1,t_3-t_1)&(50,100)&(50,150)&(50,200)&(100,150)&(100,200)&(150,200)\\
\hline
\text{root-class sum}&8&10&14&10&16&14.
\end{array}
\]
Every gap is at least $50$, so its orbit factor is at most the parenthesis in
\eqref{eq:pure-half-determinant-debit}.  A union bound over the three pairs
proves the assertion.  The residue table and the actual-support scan over all
$479960$ eligible triples at $N=2\cdot10^6$ are reproduced by
\path{scripts/systematic_pure_determinant_tail.py}.
\end{proof}

There is an even sharper one-form observation at this particular split.

\begin{corollary}[The pure very-large branch is terminally paid]
\label{cor:pure-half-one-form-debit}
Under the hypotheses of Corollary~\ref{cor:pure-half-determinant-debit}, the
number of opposite-parity base vertices for which at least one of the three
forms has a prime-square witness greater than $Y$ is at most $192$.
Consequently the determinant debit is not the controlling term in the pure
half-split continuation.
\end{corollary}

\begin{proof}
The preceding quotient argument leaves only
\[
 3p^2=at+1.
\]
It forces $t>(3(Y+1)^2-1)/N>3N/4-2$.  The interval
\[
 Y<p\leq\sqrt{(tN+1)/3}
\]
has length less than $0.079N<t$.  Thus each root class of
$3x^2\equiv1\pmod t$ contains at most one possible $p$.  If $3\mid t$ or
$4\mid t$ there are no roots.  Otherwise the prime $5$ is absent and every
odd prime-power factor contributes at most two roots.  Since
\[
 7\cdot11\cdot13\cdot17\cdot19\cdot23\cdot29
 =215656441>2\cdot10^8,
\]
$t$ has at most six relevant odd prime divisors.  Hence there are at most
$2^6=64$ witnesses for one pivot, and the union over three pivots is at most
$192$.
\end{proof}

Corollaries~\ref{cor:pure-half-determinant-debit} and
\ref{cor:pure-half-one-form-debit} unconditionally remove the very-large
branch from the pure matching obstruction.  The remaining medium window
$Q<p\leq N/2$ is bounded jointly by separating ``at least two small'' from
``at least two tail'' witnesses.  Let
\[
 \mathcal P_{47}=\{p:3\leq p\leq47,\ p\ne5\},\qquad z_p=p^{-2}.
\]
Put
\begin{align*}
 \lambda_3&=1-e_1(z)+e_2(z)-e_3(z)
             >0.847622705884,\\
 \alpha_2&=1-3\{1-e_1(2z)\}
             +2\{1-e_1(3z)+e_2(3z)\}
             <0.109696396217.
\end{align*}
The first polynomial is the degree-$3$ lower Bonferroni bound for one form.
For the second, if $S_i$ is the event that form $i$ avoids every finite
square class, the indicator that at least two forms have a finite witness is
\[
 1-e_2(S_1,S_2,S_3)+2e_3(S_1,S_2,S_3).
\]
The degree-$1$ lower bound for each two-form survivor and the degree-$2$
upper bound for the three-form survivor give precisely $\alpha_2$.  Their
respective endpoint budgets are
\begin{equation}\label{eq:pure-joint-finite-errors}
 377,\qquad
 3e_1(2\mathbf1)+2\{e_1(3\mathbf1)+e_2(3\mathbf1)\}=1560.
\end{equation}

Let $Y=\lfloor N/25\rfloor$ and
\[
 \Pi(N)=\frac{Y}{\log Y-1.1}+1.34,\qquad \tau_{47}=0.003887.
\]
For $G>0$ define
\begin{align}
 T_G(N)&=\left(\frac N{50}+1\right)\tau_{47}
          +\Pi(N)-15+G,\label{eq:pure-joint-tail}\\
 D_G(N)&=\lambda_3\left(\frac N{50}-1\right)-377-T_G(N),
          \label{eq:pure-joint-degree}\\
 U_G(N)&=\alpha_2\left(\frac N{50}+1\right)+1560
          +\frac32T_G(N).\label{eq:pure-joint-codegree}
\end{align}
The prime bound $\pi(Y)\leq\Pi(N)$ on the range used below follows from
\eqref{eq:dusart-pi} for $Y>60184$; the added $1.34$ covers the remaining
integer interval $60000\leq Y\leq60184$, which is checked directly.

\begin{lemma}[Joint pure-tail certificate]
\label{lem:pure-joint-tail-certificate}
For $1.5\cdot10^6\leq N<2\cdot10^8$, every vertex in a pure opposite-parity
block has degree at least $D_G(N)$, and every eligible span-$200$ triple has
common non-neighbourhood at most $U_G(N)$, where
\[
 G=\begin{cases}
 1541.126,&1.5\cdot10^6\leq N<2.3\cdot10^6,\\
 3605,&2.3\cdot10^6\leq N<2\cdot10^8.
 \end{cases}
\]
On both intervals
\begin{equation}\label{eq:pure-joint-strict}
 D_G(N)>U_G(N),\qquad D_G(N)>\frac N{125}+2.
\end{equation}
\end{lemma}

\begin{proof}
The one-form finite sieve gives the first term of
\eqref{eq:pure-joint-degree}.  For $47<p\leq Y$, the square-density is
included in $\tau_{47}$ and there is one terminal class per prime.  For
$p>Y$, write $at+1=p^2m$.  Since $Y\geq N/26$ and $t\leq N$, one has
$m\leq676$.  Opposite parity makes $m$ odd, the mod-$5$ calculation makes it
a nonsquare modulo $5$, and every odd prime divisor of $t$ imposes its exact
quadratic-residue condition.  The transformed root payment is
\[
 H_t\left(s_t+\frac{2(N/50+1)}Y\right),
 \qquad \frac{2(N/50+1)}Y<1.04004,
\]
where $s_t$ is the exact survivor count.  Exhausting all $92000$ pure pivots
$t\leq2.3\cdot10^6$ gives maximum $1541.12512$, attained at
\[
 t=2038168,\qquad \nu_2(t)=3,\qquad H_t=128,\qquad s_t=11.
\]

For the full upper interval, retain the smallest $k-1$ of the $k$ non-$5$
odd prime divisors of $t$.  Requiring that one further prime still fit below
$2\cdot10^8$ gives the following unconditional relaxation:
\[
\begin{array}{c|rrrr}
k&0&1&2&3\\ \hline
s_k&135&135&83&48\\
H_k&8&16&32&64\\
H_k(s_k+1.04004)&1088.321&2176.641&2689.282&3138.563
\end{array}
\]
\[
\begin{array}{c|rrrr}
k&4&5&6&7\\ \hline
s_k&26&13&6&1\\
H_k&128&256&512&1024\\
H_k(s_k+1.04004)&3461.126&3594.251&3604.501&2089.001
\end{array}
\]
Thus $G=3605$ is valid.  This proves the degree bound.

Now let $a$ be a common non-neighbour of three pivots.  If at least two forms
have witnesses in $\mathcal P_{47}$, the finite polynomial $\alpha_2$ and
the second error in \eqref{eq:pure-joint-finite-errors} pay it.  Otherwise at
least two forms have witnesses exceeding $47$.  If $H_i$ denotes the latter
tail event for form $i$, then
\[
 \#\{a:\text{at least two }H_i\}
 \leq\frac12\sum_{i=1}^3|H_i|\leq\frac32T_G(N).
\]
This proves \eqref{eq:pure-joint-codegree}; in particular the medium terminal
events are paid jointly rather than independently.

It remains only to check the two scalar margins.  Both are increasing on
each displayed interval.  Between jumps of $Y$ this is immediate; at a jump,
differentiating $y/(\log y-1.1)$ gives an increment less than $0.091$, which
is dominated by the linear increment in both margins.  More explicitly, an
increment of $25$ in $N$ increases $D_G-U_G$ by more than
\[
 \frac{\lambda_3-\alpha_2}{2}-\frac54\tau_{47}-2.5(0.091)>0.136,
\]
and increases $D_G-(N/125+2)$ by more than
\[
 \frac{\lambda_3}{2}-\frac15-\frac12\tau_{47}-0.091>0.130.
\]
At the two left
endpoints the rigorous rounded values are
\[
\begin{array}{c|c|c|c|c}
N&G&D_G&U_G&D_G-U_G\\ \hline
1500000&1541.126&>17347.432&<16406.103&>941.330\\
2300000&3605&>25936.158&<25621.101&>315.058\\
\end{array}
\]
and $D_G-(N/125+2)$ is respectively greater than $5345.43$ and $7534.15$.
The support enumeration, finite polynomials, prime counts, and endpoint
margins are reproduced by
\path{scripts/systematic_pure_joint_tail_close.py}.
\end{proof}

\begin{theorem}[Pure matching down to one and a half million]
\label{thm:pure-joint-tail-matching}
For every $N\geq1.5\cdot10^6$, the squarefree-edge graph between
$C_N^{18}$ and $C_N$ has a matching saturating $C_N^{18}$.  Hence every
$B\subseteq C_N^{18}$ satisfies the original Hall inequality.
\end{theorem}

\begin{proof}
Theorem~\ref{thm:actual-support-pure-matching} handles
$N\geq2\cdot10^8$.  In the remaining range work in one parity block and
suppose $S$ is a Hall defect, with $m=|S|$.  The forward degree bound gives
$m>D_G(N)>N/125+2$.  If no three consecutive members of $S$ had span at
most $200$, their spans, being multiples of $50$, would all be at least
$250$, whence
\[
 250(m-2)\leq\sum_{j=1}^{m-2}(x_{j+2}-x_j)\leq2(N-1),
\]
contrary to $m>N/125+2$.  Choose such a triple.

As in Theorem~\ref{thm:systemic-pure-parity-matching}, the reverse degree
bound and the assumed Hall defect force the unused right-hand set to have
more than $D_G(N)$ members.  Every one is a common non-neighbour of the
chosen triple, contradicting $U_G(N)<D_G(N)$.  Hall's theorem supplies the
matching in both parity blocks, and restricting it to $B$ proves the last
assertion.
\end{proof}

\subsection{Joint mixed tail descent to forty million}

The same joint split has a stronger mixed form.  Averaging the
all-three-finite intersection with the two-of-three event preserves the mixed
diagonal balance.

\begin{lemma}[Half finite-intersection, half two-of-three]
\label{lem:mixed-half-finite-tail}
For three pairs of events $F_i,H_i$ one has pointwise
\begin{align}
 \mathbf 1_{\bigcap_{i=1}^3(F_i\cup H_i)}
 \leq{}&\frac12\mathbf 1_{F_1\cap F_2\cap F_3}
 +\frac12\mathbf 1_{\{\#\{i:F_i\}\geq2\}}
 +\frac12\sum_{i=1}^3\mathbf 1_{H_i}.
 \label{eq:mixed-half-finite-tail}
\end{align}
\end{lemma}

\begin{proof}
If exactly $k$ of the $F_i$ occur and all three unions occur, the right side
is at least $1$ for $k=3,2,1,0$, respectively: its minimum in the four cases
is $1,1,1,3/2$.  This also proves the assertion when some $F_i$ and $H_i$
occur simultaneously.
\end{proof}

We apply the lemma with $F_i$ the existence of a prime-square witness through
$47$ and $H_i$ the existence of a witness above $47$.  Put
\[
 \eta=0.00015,\qquad \tau_{47}=0.003887.
\]
Retaining the nonzero square coset modulo $5$ in both the finite and
actual-support degree arguments gives the following low block.

\begin{proposition}[Mixed residual below two hundred million]
\label{prop:mixed-joint-low-residual}
For $4\cdot10^7\leq N<2\cdot10^8$, every mixed Hall defect satisfies
\[
 |A^*|>0.00030627N>2\eta N.
\]
\end{proposition}

\begin{proof}
First let $4\cdot10^7\leq N<8\cdot10^7$.  Use the one-form
degree-$3$ Bonferroni lower bound through $47$, split the remaining primes at
$Y=\lfloor N/70\rfloor$, and retain the actual support and worse nonzero
mod-$5$ coset in the root tail.  The global support at $8\cdot10^7$ fits in
$M=4902$.  Evaluated at the lower endpoint, the least row is
\[
 \text{odd},\qquad k=4,\qquad H=32,\qquad s_k=286,
 \qquad (p_1,p_2,p_3)=(61,73,79).
\]
After the reciprocal-survivor correction the normalized mixed residual is
greater than
\[
 0.0077080.
\]

For $8\cdot10^7\leq N<2\cdot10^8$, use
$D=\lfloor N/300\rfloor$ and the common box $M=90001$ exactly as in
Proposition~\ref{prop:actual-support-degree}, but intersect every
complementary mask with the worse of the two nonzero square cosets modulo
$5$.  To make the calculation uniform, take the prime-support superset at
$2\cdot10^8$ and evaluate every fixed row at $N=8\cdot10^7$.  The least row
is
\[
 \text{odd},\qquad k=3,\qquad H=16,\qquad s_k=9176,
 \qquad (p_1,p_2)=(449,491).
\]
Its excess over $N/50$ is greater than $12542.5712308$.  The same reverse
degree subtraction, after including the reciprocal-survivor correction
$<291.111$, therefore gives
\[
 \frac{2(12542.5712308-291.111)-1}{8\cdot10^7}
 >0.00030627.
\]
In both blocks the global support is fixed, the normalized affine main term
increases, and all normalized positive payments decrease.  Thus the lower
 endpoints control their intervals.  The support enumeration is reproduced by
\path{scripts/systematic_mixed_interval_support.py}; the endpoint
subtraction is checked by
\path{scripts/verify_four_range_paper_arithmetic.py}.
\end{proof}

We next record the finite payment.  On an odd-prime progression let $S_i$ be
the event that the $i$th form survives every prime through $47$.  The
indicator of at least two finite-bad forms is
\[
 1-\sum_{i<j}\mathbf 1_{S_iS_j}
   +2\mathbf 1_{S_1S_2S_3}.
\]
Use degree-$1$ lower bounds for the three pair survivors and a degree-$2$
upper bound for the triple survivor.  At a common prime the pair
multiplicities sum to $3$ and the triple multiplicity is $1$.  At a
pair-coincident non-common prime they are $5$ and $2$; without a coincidence
they are $6$ and $3$.  If $z_p=p^{-2}$, declaring a non-common prime
pair-coincident changes the resulting upper polynomial by
\[
 z_p(1-2S_p)>0,
 \qquad S_p<3\sum_{\substack{3\leq q\leq47\\q\ne5}}q^{-2}<\frac12.
\]
Hence all non-common odd primes may be declared pair-coincident for this
upper bound.  A common prime only decreases the two-of-three polynomial, so
the two finite polynomials are evaluated jointly.

Prime $2$ is exact.  For a triple meeting both odd classes modulo $4$, one
mod-$4$ class makes the close pair automatically bad, the opposite class
makes the third form automatically bad, and the other two classes use the
odd-prime two-of-three polynomial.  For three pivots in one odd mod-$4$
class, one class makes all three forms automatically bad and the other three
use the odd-prime polynomial.

If the even mass exceeds $\eta N$, a close pair in one even mod-$4$ class has
gap less than $13334$.  After charging the even mass in the other case, the
odd mass is greater than $0.0001635517807N$.  A close pair in the denser of
two odd classes has gap less than $12229$; if only one odd class occurs, its
close pair has gap less than $6115$.  For any common odd-prime product $Q$,
the close pair gives $4Q^2<d$.  Exhausting the resulting cutoff-prime
supports gives
\[
\begin{array}{c|c|c|c}
\text{branch}&\text{maximizing common support}&I_3+I_2&\epsilon_3+\epsilon_2\\ \hline
\text{even generic}&7&<0.223675978&4874\\
\text{even concentrated}&3,7&<0.309830838&4510\\
\text{two odd classes}&\varnothing&<0.541995982&7113\\
\text{two odd, common }3&3&<0.543144010&6615\\
\text{one odd class}&2&<0.679312680&7435\\
\text{one odd, common }2,3&2,3,7&<0.739301322&6423.
\end{array}
\]
Here $I_3$ is the degree-$3/2/3$ all-three-finite bound used above, and
$I_2$ is the new two-of-three bound.  By
Lemma~\ref{lem:mixed-half-finite-tail}, the normalized finite payment across
the two base progressions is therefore
\begin{equation}\label{eq:mixed-joint-finite-payment}
 \frac{I_3+I_2}{25}+\frac{\epsilon_3+\epsilon_2}{N}.
\end{equation}

The finite polynomial \eqref{eq:mixed-joint-finite-payment} is combined with
the ten-branch envelopes below, each of which includes the reciprocal mass
beyond $3025$.  This keeps each root bound paired with the finite row at the
same cutoff.

Finally, the eventwise M\"obius--CRT diagonal scan with $R=2\cdot10^7$
uses the following entirely rational tail estimate.  If $P$ is the exact
truncated signed count at the right endpoint $U$, and $\theta$ is the
density of the permitted residue set, then throughout $[L,U]$
\begin{equation}\label{eq:rational-diagonal-envelope}
 \frac{\Delta(N)}N<
 \frac1L\left\{
 P+\frac{38\theta U}{R}
 +\frac{13(U^2+1)}{R^2}\right\}.
\end{equation}
Indeed $\log(2\cdot10^7)+2<19$: the positive Taylor series gives
$e>2.7$, while $2.7^{17}>2\cdot10^7$.  Every comparison in
\eqref{eq:rational-diagonal-envelope} is made by unsigned-integer
cross-multiplication.  On all nine blocks from $5\cdot10^6$ to
$2\cdot10^8$ this gives the uniform strict ceilings
\[
\begin{gathered}
\begin{array}{c|rrrr}
\text{support}&\text{all}&\text{even cell}&
\text{odd or }8\mid a&E_2\text{ cell}\\ \hline
10^6\Delta/N<&25294&19681&15812&16164
\end{array}
\\[3pt]
\begin{array}{c|rrrr}
\text{support}&\text{odd union}&\text{odd+cell}&
\text{one odd}&\text{one odd cell}\\ \hline
10^6\Delta/N<&12652&7036&6330&711.
\end{array}
\end{gathered}
\]
The parameterized scan is reproduced by
\begin{center}\small
\path{scripts/systematic_concentrated_diagonal_blocks.cpp}.
\end{center}

\begin{theorem}[Joint mixed close to forty million]
\label{thm:mixed-joint-tail-close}
The Hall inequality \eqref{eq:hall-defect}, and hence the exact bound in
Erd\H{o}s Problem~848, holds for every
\[
 N\geq4\cdot10^7.
\]
\end{theorem}

\begin{proof}
Theorem~\ref{thm:actual-support-mixed-close} handles $N\geq2\cdot10^8$.
In the remaining interval take a mixed defect and put
$A=B\cup T_N(B)$.  Proposition~\ref{prop:mixed-joint-low-residual} supplies
more than three sparse-class charges.  Apply the ten exhaustive
valuation/parity alternatives used in
Proposition~\ref{prop:twenty-million-interval-close}; their selection proof
uses only the residual inequality and the gap lemma, and is therefore valid
unchanged on this larger interval.

For the six blocks
\[
[40,50),[50,70),[70,80),[80,100),[100,150),[150,200)
\quad\text{million},
\]
the transformed-root ceilings for $E_1,E_2,E_3$, and the odd
classes are
\[
\begin{array}{c|rrrr}
{[L,U)}&E_1&E_2&E_3&O\\ \hline
{[40,50)\mathrm M}&.009677836&.009973185&.010991872&.010396089\\
{[50,70)\mathrm M}&.009350108&.009648325&.010405463&.009926965\\
{[70,80)\mathrm M}&.008921868&.009135839&.009679089&.009335765\\
{[80,100)\mathrm M}&.008771555&.008959044&.009511773&.009134226\\
{[100,150)\mathrm M}&.008534837&.008684996&.009127671&.008825297\\
{[150,200)\mathrm M}&.008175215&.008275070&.008569448&.008368371.
\end{array}
\]
These ceilings include the reciprocal-survivor correction.  Combine them
with the exact diagonal grid, the ten finite payments
\[
 .008685,\ .012616,\ .019420,\ .020878,\ .026643,\ .029459
\]
in their respective generic/common branches, the sparse-class charge
$1/20001+1/L$, and the cutoff-$23$ or cutoff-$19$ square tail.  Exact
rational comparison shows that $E_2$-generic is the largest row in every
block:
\[
\begin{array}{c|c|c}
{[L,U)}&\max |A|/N&
  1/25-7/(25L)-\max |A|/N\\ \hline
{[40,50)\mathrm M}&.039972160176&>.0000278328\\
{[50,70)\mathrm M}&.039810601176&>.0001893932\\
{[70,80)\mathrm M}&.039554376176&>.0004456198\\
{[80,100)\mathrm M}&.039466546676&>.0005334498\\
{[100,150)\mathrm M}&.039330914676&>.0006690825\\
{[150,200)\mathrm M}&.039127498676&>.0008724994.
\end{array}
\]
For example, with $\widehat\tau_{23}$ as in the ten-million proof, the
controlling lower-endpoint total is
\[
 C_{40}=\frac{25288306}{10^9}+\frac{8685}{10^6}
 +\frac3{25}\widehat\tau_{23}
 +\frac12\frac{9973185}{10^9}
 <\frac1{25}-\frac7{25(4\cdot10^7)}-\frac{27}{10^6}.
\]
Thus every branch contradicts the Hall defect.  The rows and exact rational
inequalities are checked by
\path{scripts/verify_four_range_paper_arithmetic.py}.
\end{proof}

\subsection{Valuation, cell, and fibre descent to five million}

The remaining factor-of-eight descent refines the same mixed defect by the
$2$-adic valuation of its even pivots and by the
actual residue cells occupied by the odd pivots.  We first isolate the two
finite combinatorial facts used below.

\begin{lemma}[Valuation support and the four-pivot inequality]
\label{lem:valuation-four-pivot}
Let $a$ be even, $p$ odd, and $p^2m=at+1$.  If
$\nu_2(a)=1,2$, or at least $3$, then respectively
\[
 m\equiv1\pmod2,\qquad m\equiv1\pmod4,\qquad m\equiv1\pmod8.
\]
Moreover, for four pairs of events $F_i,H_i$,
\begin{equation}\label{eq:four-pivot-boolean}
 \mathbf1_{\cap_{i=1}^4(F_i\cup H_i)}
 \leq \frac12\sum_{|S|=3}\mathbf1_{\cap_{i\in S}F_i}
       +\frac12\sum_{i=1}^4\mathbf1_{H_i}.
\end{equation}
More generally, for $n$ pairs and $1\leq k\leq n$,
\begin{equation}\label{eq:n-pivot-boolean}
 \mathbf1_{\cap_{i=1}^n(F_i\cup H_i)}
 \leq \mathbf1_{\#\{i:F_i\}\geq k}
       +\frac1{n-k+1}\sum_{i=1}^n\mathbf1_{H_i}.
\end{equation}
\end{lemma}

\begin{proof}
Every odd square is $1$ modulo $8$.  Since $2^{\nu_2(a)}\mid at$ and
$p^2\equiv1\pmod8$, division by $p^2$ gives the three displayed quotient
restrictions.  For
\eqref{eq:four-pivot-boolean}, if exactly $j$ of the $F_i$ occur, the right
side is at least $2,3/2,1,1,2$ for $j=4,3,2,1,0$, respectively.  For
\eqref{eq:n-pivot-boolean}, either at least $k$ finite events occur, or at
least $n-k+1$ of the unions must be supplied by their $H_i$ terms.
\end{proof}

\begin{lemma}[Two-coordinate fibres and capacitated cells]
\label{lem:cell-fibre-capacity}
Let $X$ be a finite set and let $r_q:X\to\mathbb Z/q^2\mathbb Z$.
\begin{enumerate}
\item If every pair of elements of $X$ agrees in at least one of the two
coordinates $r_7,r_{11}$, then $X$ is contained in one mod-$49$ fibre or in
one mod-$121$ fibre.
\item Let $S$ be a set of $s$ occupied mod-$9$ cells, joined to the mod-$49$
residues which they contain, and give every right vertex capacity two.  If
there is no matching of size eight, then some $Y\subseteq S$ satisfies
\begin{equation}\label{eq:capacity-defect}
 |Y|-2|N(Y)|\geq s-7.
\end{equation}
\end{enumerate}
\end{lemma}

\begin{proof}
For the first assertion, choose two points with different first coordinates.
They have the same second coordinate.  Any third point which had a different
second coordinate would have to share both distinct first coordinates, which
is impossible.  The second assertion is the max-flow/min-cut formula for a
bipartite matching with right capacities two:
\[
 \nu=\min_{Y\subseteq S}\bigl(s-|Y|+2|N(Y)|\bigr).
\]
If $\nu\leq7$, rearranging gives \eqref{eq:capacity-defect}.
\end{proof}

\begin{lemma}[Gap selection and cell charge]
\label{lem:gap-selection-charge}
Let $X\subseteq[1,N]$ and let $G\geq1$.  If no two members of $X$ have
difference at most $G$, then
\[
 |X|\leq\frac{N-1}{G+1}+1
 <N\left(\frac1{G+1}+\frac1N\right).
\]
Consequently, if $X$ is split into $r$ residue cells, each nondense cell of
cardinality at most $(N-1)/(G+1)+1$ may be discarded at total normalized
cost at most
\[
 r\left(\frac1{G+1}+\frac1N\right).
\]
\end{lemma}

\begin{proof}
If $x_1<\cdots<x_m$ are the members of $X$, then
$x_{i+1}-x_i\geq G+1$, so $(m-1)(G+1)\leq N-1$.  Summing this estimate over
the cells proves the second assertion.
\end{proof}

We now close the two upper intervals.  In both proofs split the even pivots
into $E_1,E_2,E_3$ according as $\nu_2=1,2$, or at least $3$.  The quotient
restrictions in Lemma~\ref{lem:valuation-four-pivot} reduce their transformed
root supports.

\begin{proposition}[Valuation close on the twenty-million block]
\label{prop:twenty-million-interval-close}
The Hall inequality \eqref{eq:hall-defect} holds for
\[
 2\cdot10^7\leq N<4\cdot10^7.
\]
\end{proposition}

\begin{proof}
The cutoff-$19$ degree-$3$ one-form certificate gives
\[
 |A^*|>0.006910733259914470N.
\]
Put
\[
 \delta_{20}=\frac1{20001}+\frac1{2\cdot10^7}.
\]
The four six-form root ceilings for $E_1,E_2,E_3$, and an odd pivot are
\[
 .008622089,\qquad .010006474,\qquad .012915119,
 \qquad .010177869.
\]

If $|E_1|>\delta_{20}N$, Lemma~\ref{lem:gap-selection-charge} supplies a
pair of gap at most $20000$.  Add a third $E_1$ pivot.  If the class is
nonconstant modulo $9$, choose the third pivot outside the pair's cell and
use the generic row; otherwise use the common-$3$ row.  This is the
three-pivot inequality \eqref{eq:mixed-half-finite-tail}; no four-pivot
claim is used in this interval.  If $|E_1|\leq\delta_{20}N$, charge it and
repeat with $E_2$.  If both have been charged, repeat with $E_3$; the
off-base diagonal is then supported on the odd-or-divisible-by-$8$ set.
If all three even classes have been charged, their total cost is at most
$3\delta_{20}N$.

The remaining odd mass is greater than
\[
 (0.006910733259914470-3\delta_{20})N.
\]
If both odd mod-$4$ classes occur, its denser class contains a pair of gap
less than twice the reciprocal of this density; add a pivot from the other
class.  If only one odd class occurs, its close-pair bound is the reciprocal
of the same density and the third pivot is taken in that class.  Splitting
each alternative according as the three pivots are constant modulo $9$
gives the final four rows.  Thus the following ten rows are exhaustive:
\[
\begin{array}{c|ccccc}
\text{row}&E_1\mathrm{G}&E_1{}_3&E_2\mathrm{G}&E_2{}_3&E_3\mathrm{G}\\ \hline
\text{total}<&.039257774&.037068015&.039949967&.034297022&.032021074
\end{array}
\]
\[
\begin{array}{c|ccccc}
\text{row}&E_3{}_3&O_2\mathrm{G}&O_2{}_3&O_1\mathrm{G}&O_1{}_3\\ \hline
\text{total}<&.035449069&.039924070&.034577881&.039438565&.036137568.
\end{array}
\]

The controlling row is now the generic $E_2$ row.  With the exact finite
and square-tail payments, its total is
\begin{align}
C_{20}={}&\frac{25289550}{10^9}
 +\frac{1572362642548502371982744227834515457}
 {171081196711186531200844968164610000000}\\
&+\frac{11661}{25000000}
 +\frac12\frac{10006474}{10^9}.                 \label{eq:twenty-million-total}
\end{align}
Exact common-denominator arithmetic gives
\[
 C_{20}<\frac1{25}-\frac7{25(2\cdot10^7)}
 -\frac{50}{10^6}.
\]
All other rows have larger slack, and every normalized endpoint payment
decreases while the Hall target increases.  This proves the full interval.
The support exhaustion is reproduced by the accompanying finite
enumerations, and the ten rational totals are checked by
\path{scripts/verify_four_range_paper_arithmetic.py}.
\end{proof}

\begin{proposition}[Cell close on the ten-million block]
\label{prop:ten-million-interval-close}
The Hall inequality \eqref{eq:hall-defect} holds for
\[
 10^7\leq N<2\cdot10^7.
\]
\end{proposition}

\begin{proof}
Call a mod-$9$ cell dense if its cardinality is more than
$(N-1)/1000001+1$, and put
\[
 \delta_{10}=\frac1{1000001}+\frac1{10^7}.
\]
The cutoff-$17$ degree certificate gives
$|A^*|>.005826196348744825N$.  The relevant root ceilings are
\[
 .010545872,\qquad .011845571,\qquad .014894092,\qquad .013231580.
\]

In $E_1$, two dense cells supply two close pairs and hence the four-pivot
row \eqref{eq:four-pivot-boolean}; one dense cell supplies three pivots and
the other eight cells cost at most $8\delta_{10}N$.  If there is no dense
cell, retain the whole class in the unrestricted diagonal and continue to
$E_2$.  The same alternatives for $E_2$ give the unrestricted two-cell row
or, after charging the nine $E_1$ cells and eight unused $E_2$ cells, the
$E_2$-cell row with charge $17\delta_{10}N$.  If neither class supplies a
row, charge both and use a dense $E_3$ cell; this costs
$18\delta_{10}N$ and leaves the odd-or-divisible-by-$8$ diagonal.  If no
even class supplies a row, all $27$ even cells cost at most
$27\delta_{10}N$.

For the odd remainder use the cutoff-$7$ threshold allocation described
below for the five-million block.  Discarding odd cells of size at most five
costs at most $90$ pivots.  The alternatives are: disjoint triples of active
cells in the two odd mod-$4$ classes; at most eleven active cells; one
concentrated residue; or one odd mod-$4$ class with $s=1,\ldots,9$ active
cells.  Lemma~\ref{lem:cell-fibre-capacity} and the exact recurrence
\eqref{eq:collision-dp} cover the $s=8,9$ matching failures.  This avoids the
invalid inference that crossed common supports at $7$ and $11$ must put all
pivots into one fixed prime-square fibre.

Using the five-million diagonal ceilings, which also dominate the direct
scan on this interval, every odd row is at most
$0.039247000$.  The five even rows are, respectively,
\[
\begin{array}{c|ccccc}
\text{row}&E_1\text{ two}&E_1\text{ one}&E_2\text{ two}&E_2\text{ one}&E_3\\ \hline
\text{total}<&.039077928&.039259999&.039944394&.036406507&.037579342.
\end{array}
\]
Thus the controlling row is $E_2$ with two dense cells.  Its four
three-form finite intersections admit the conservative cutoff-$23$ density
\[
 I_{10}=\frac{5355507302216}{165901428816123},
 \qquad \epsilon_{10}=607.
\]
Put also
\[
 \widehat\tau_{23}=
 \frac{64081802747648035629863}
 {7596668444022826249000000}.
\]
This upper bound uses cutoff $23$.  The decimal-free controlling total is
\begin{align}
C_{10}={}&\frac{25289862}{10^9}
 +\frac{4I_{10}}{25}+\frac{4\epsilon_{10}}{10^7}
 +\frac4{25}\widehat\tau_{23}
 +\frac23\frac{11845571}{10^9}.                 \label{eq:ten-million-rational-row}
\end{align}
Direct common-denominator arithmetic yields
\begin{equation}\label{eq:ten-million-total}
 C_{10}<\frac1{25}-\frac7{25\cdot10^7}
 -\frac{55}{10^6}.
\end{equation}
The exact rational calculation above is reproduced by
\begin{center}
\path{scripts/verify_four_range_paper_arithmetic.py}.
\end{center}
The diagonal ceilings come from
\begin{center}
\path{scripts/systematic_concentrated_diagonal_blocks.cpp}.
\end{center}
\end{proof}

For the lower block the support constraints are retained jointly.  On
$5\cdot10^6\leq N<10^7$, the cutoff-$17$
degree calculation gives
\begin{equation}\label{eq:five-million-residual}
 |A^*|>0.003997632216030162N.
\end{equation}
The relevant cutoff root ceilings are
\begin{equation}\label{eq:five-million-roots}
 T_{E_1,23}<.013878159,\quad T_{E_2,7}<.015472245,\quad
 T_{E_3,23}<.020270750,\quad T_{O,7}<.017398287.
\end{equation}

For $E_1$, call a paired mod-$9$ cell good if it is contained in neither one
mod-$49$ fibre nor one mod-$121$ fibre.  Lemma~
\ref{lem:cell-fibre-capacity}(1) supplies a pair which breaks both primes.
Otherwise label the cell $7$-only, $11$-only, or both.  Every cutoff-$23$
common support compatible with the label is enumerated; in particular,
the primes $13,17,19,23$ are not silently declared non-common.  If there are
two good cells, use four pivots.  If there is one, pair it with the best bad
cell.  If there is none, every $7$-only or both cell is inserted into the
eventwise diagonal through its actual mod-$49$ fibre.  The maxima for
$E_{1,0}$ plus $k=0,\ldots,9$ such fibres have the uniform strict rational
ceiling
\[
 .019113,
\]
instead of the incompatible raw-fibre/unrestricted hybrid.  The one-paired
cell case uses three pivots.  The resulting four $E_1$ envelopes are
\[
 .039267556,\qquad .039360535,\qquad .035990495,\qquad .039569683.
\]
For $E_2$, the two-cell and one-cell rows are below $.037494782$ and
$.038211102$.  If $E_1,E_2$ are both charged, nineteen $E_3$ pivots force
three in one mod-$9$ cell; its envelope is below $.039335952$.

For the odd part, discard every mod-$9$ cell which contains at most five
pivots and call the other cells active.  Together with the preceding even charges this
costs at most $128$ pivots.  If both odd mod-$4$ classes admit disjoint triples
of active mod-$9$ cells, choose three pivots from each.  Otherwise an elementary
three-set check gives at most eleven active odd cells: if one class has at
most two cells the assertion is immediate, while if both have at least three,
the failure of disjoint triples forces both cell sets to have size at most
five.  A pair from each class
in different mod-$9$ cells then gives a four-pivot row below $.037135620$; if all
active cells have the same mod-$9$ residue, the concentrated three-pivot row
is below $.033200482$.

For one odd mod-$4$ class, let $s$ be its number of active mod-$9$ cells.
When $s=8$ or $9$, give the mod-$49$ residues capacity two.  A matching of
size eight selects eight pivots with no mod-$49$ collision block larger than
two.  If it fails, \eqref{eq:capacity-defect} restricts the diagonal.  The
eventwise maxima in the two failure cases are
\[
 .003572,\qquad .003587.
\]
When $s=1$, either four mod-$49$ fibres each supply a pair, or all but at most
$49$ pivots lie in three such fibres; the latter diagonal is below
$.000050$.  The cases $2\leq s\leq7$ use six pivots distributed as evenly
as possible among the active cells.

Here the finite calculation is genuinely finite, with no hidden limiting
step.  Put $M=4\cdot9\cdot49=1764$ and let $V_p(S)$ be the worst conditional
probability, after the primes above $p$ have been exposed, that the terminal
set of finite-bad forms has the required cardinality.  If the collision
partition at $p$ is $\mathcal Q$, then backwards induction uses
\begin{equation}\label{eq:collision-dp}
 V_{p^-}(S)=\max_{\mathcal Q}
 \left(\frac{p^2-|\mathcal Q|}{p^2}V_p(S)
 +\frac1{p^2}\sum_{Q\in\mathcal Q}V_p(S\cup Q)\right).
\end{equation}
Allowing the maximizing partition to depend on $S$ only enlarges the true
density.  If the resulting density is $R/M$, then in either base progression
every prefix has excess at most
\begin{equation}\label{eq:period-prefix}
 R\left(1-\frac RM\right),
\end{equation}
because a prefix of length $u$ contains at most $\min(u,R)$ favourable
residues.  Thus both the main density and the integer endpoint are certified.

The controlling periodic rows are
\[
\begin{array}{c|c|c|c|c}
\text{branch}&R/M&\text{tail factor}&\text{diagonal}&\text{total}\\ \hline
O_1,\ s=2&40/147&2/3&.001415&.039866076\\
O_1,\ s=7&55/196&1/2&.004928&.039932809\\
O_1,\ s=9,\ \text{matching}&1/4&8/15&.006330&.039701523\\
O_1,\ s=9,\ \text{failure}&55/196&1/2&.003587&.038591809\\
O_2,\ \text{balanced}&1/42&1&.012652&.039374127.
\end{array}
\]
The controlling row admits a decimal-free exact comparison.  Define
\[
 \widehat\tau_7=\frac{3887}{10^6}
  +\sum_{p\in\{11,13,17,19,23,29,31,37,41,43,47\}}\frac1{p^2}.
\]
The cutoff-$7$ dynamic programme gives exactly $R/M=55/196$ and prefix
excess $69795/196$.  With the upward rational ceilings
\[
 D_{5,7}<\frac{4928}{10^6},
 \qquad T_{O,7}\leq\frac{17398287}{10^9},
\]
the complete $O_1,s=7$ row is at most
\begin{align}
 C_5={}&\frac{4928}{10^6}
 +\frac2{25}\frac{55}{196}
 +\frac2{5\cdot10^6}\left(\frac{55}{196}+\frac{69795}{196}\right)
 +\frac12\frac{17398287}{10^9}\notag\\
 &+\frac3{25}\widehat\tau_7+\frac{128}{5\cdot10^6}.
                                                        \label{eq:five-million-rational-row}
\end{align}
Exact common-denominator arithmetic yields
\[
 C_5<\frac1{25}-\frac7{25(5\cdot10^6)}-\frac{67}{10^6}.
\]
Consequently the controlling five-million row retains rational slack greater
than $6.7\cdot10^{-5}$ after outward rounding.

Every omitted subrow is smaller.  The exact rational dynamic programme,
prefix payments, root jumps, fibre classifications, and all diagonal maxima
are reproduced by
\begin{center}\small
\path{scripts/systematic_mixed_cell_5m_close.py},\qquad
\path{scripts/systematic_concentrated_diagonal_blocks.cpp}.
\end{center}

\begin{proposition}[Cell--fibre close on the five-million block]
\label{prop:five-million-interval-close}
The Hall inequality \eqref{eq:hall-defect} holds for
\[
 5\cdot10^6\leq N<10^7.
\]
\end{proposition}

\begin{proof}
Suppose that a strict Hall defect exists.  The residual estimate
\eqref{eq:five-million-residual} supplies more than
$0.003997632216030162N$ pivots.  Call an $E_1$ mod-$9$ cell paired when it
contains at least two pivots.  If at least two cells are paired, classify
them as good, $7$-only, $11$-only, or both.  Two good cells, one good cell,
and no good cell give respectively the first three $E_1$ envelopes displayed
above; Lemma~\ref{lem:cell-fibre-capacity}(1) proves that these cases exhaust
the possible $7$--$11$ collision patterns.  If exactly one cell is paired
and contains at least three pivots, the four one-cell rows apply.  Otherwise
$|E_1|\leq10$, which is charged.

After that charge, the identical paired-cell split in $E_2$ gives its
two-cell or one-cell row.  If neither applies, $|E_2|\leq10$ and it too is
charged.  Nineteen remaining $E_3$ pivots force three into one of the nine
mod-$9$ cells; if there are fewer, $|E_3|\leq18$.  Thus failure to select an
even row costs at most $10+10+18=38$ pivots.

In the odd part, cells containing at most five pivots cost at most
$2\cdot9\cdot5=90$ more.  Hence the total raw charge is at most $128$.
If both odd mod-$4$ classes have active cells and admit disjoint triples of
cell residues, use the balanced six-pivot row.  If not, the argument above
gives at most eleven active cells.  Two different residues give the
four-pivot row, while one residue gives the concentrated three-pivot row.
These are all possibilities with both odd classes active.

It remains that only one odd mod-$4$ class is active.  Let $s$ be its number
of active cells.  For $2\leq s\leq7$, six pivots distributed as evenly as
possible give the listed periodic rows.  For $s=8,9$, a capacity-two matching
either selects eight pivots with collision blocks of size at most two or
Lemma~\ref{lem:cell-fibre-capacity}(2) gives the displayed restricted
diagonal.  For $s=1$, either four mod-$49$ fibres supply pairs, or at most
three fibres are nonsingletons and the remaining at most $49$ pivots are
charged.  Therefore the branch list used in
\path{scripts/systematic_mixed_cell_5m_close.py} is exhaustive.

For every branch, \eqref{eq:n-pivot-boolean}, the exact periodic recurrence
\eqref{eq:collision-dp}, the prefix bound \eqref{eq:period-prefix}, and the
appropriate diagonal and root ceilings bound the Hall completion by its
listed total.  The largest is the $O_1,s=7$ row, and
\eqref{eq:five-million-rational-row} puts it strictly below
$1/25-7/(25N)$.  Hence $|B|+|T_N(B)|\leq|C_N|$, contradicting the strict Hall
defect.
\end{proof}

\begin{theorem}[Direct interval splice at five million]
\label{thm:mixed-cell-five-million-close}
The Hall inequality \eqref{eq:hall-defect}, and hence the exact upper bound
in Erd\H{o}s Problem~848, holds for every
\[
 N\geq5\cdot10^6.
\]
\end{theorem}

\begin{proof}
Proposition~\ref{prop:five-million-interval-close} handles
$5\cdot10^6\leq N<10^7$,
Proposition~\ref{prop:ten-million-interval-close} handles
$10^7\leq N<2\cdot10^7$,
Proposition~\ref{prop:twenty-million-interval-close} handles
$2\cdot10^7\leq N<4\cdot10^7$, and
Theorem~\ref{thm:mixed-joint-tail-close} handles $N\geq4\cdot10^7$.
The four intervals are disjoint and exhaustive.
\end{proof}


\subsection{Full-support continuation below five million}

The low degree calculation can be strengthened because every prime divisor
of a pivot, rather than only its smallest prefix, may be retained in the
quadratic-residue mask.  For $M=4902$ the exact survivor maxima at product
bound $3\cdot10^6$ are
\[
\begin{array}{c|rrrrrrr}
 &k=0&1&2&3&4&5&6\\ \hline
\text{odd}&1961&1194&612&312&153&73&28\\
\nu_2=1&981&594&312&156&80&38&15\\
\nu_2=2&491&312&171&92&54&28&13\\
\nu_2\geq3&246&169&102&58&37&22&12.
\end{array}
\]
These maxima include the worse nonzero square coset modulo $5$ and the
quotient congruences modulo $2,4,8$.  For example, the controlling support
is $(73,149,211)$ in the odd $k=3$ row.

\begin{proposition}[Full-support low residual and bounded clusters]
\label{prop:full-support-low-cluster}
On $1.5\cdot10^6\leq N<3\cdot10^6$, every mixed Hall defect has more than
\begin{equation}\label{eq:full-support-low-residual}
 \alpha N,
 \qquad
 \alpha=0.000249932448083113\ldots
\end{equation}
non-base residual pivots.  Among them there are three of span at most $8045$
and eight of span at most $28540$.
\end{proposition}

\begin{proof}
Put $Y=\lfloor N/70\rfloor$, $X=N/25+1$, and sieve through $17$.
For each displayed support row retain the active degree-$3$ Bonferroni lower
bound, subtract
\[
 X\tau_{17}+\pi(Y)-\pi(17)
 +H\left(s+\frac{2X}{Y}\right)+\frac N{50},
\]
including its exact finite discrepancy.  The normalized minimum is the odd
$k=3$, $H=16$, $s=312$ row at $N=1.5\cdot10^6$; before division by $N$ its
reverse-degree residual is
\[
 374.898672124669531\ldots.
\]
For fixed $Y$ every row is affine in $1/N$ and increases.  Between consecutive
prime jumps the normalized positive payments decrease; exact rational checks
at every prime jump show that the same lower endpoint is the global minimum.

If $x_1<\cdots<x_m$ are the residual pivots, then for $r=3$ or $8$,
\[
 \sum_{i=1}^{m-r+1}(x_{i+r-1}-x_i)\leq(r-1)(N-1).
\]
Using $m>\alpha N$ and exhausting the fewer than $376$ possible integer
values of $\lfloor\alpha N\rfloor$ gives the two stated spans.  The support
search, rational degree rows, prime jumps, and integer span exhaustion are
reproduced by \path{scripts/systematic_full_support_scan.cpp},
\path{scripts/systematic_one_prime_support_scan.cpp}, and
\path{scripts/systematic_full_support_degree_low.py}.
\end{proof}

We use the following exact one-form tail.

\begin{lemma}[Exact complete one-form tail]
\label{lem:exact-low-complete-tail}
For $1.5\cdot10^6\leq N<3\cdot10^6$, $1\leq t\leq N$, and
\[
 \mathcal M_t(N)=\left\{a\leq N:
 \begin{array}{l}
 a\equiv7\text{ or }18\pmod {25},\text{ and some prime }p\text{ satisfies}\\
 D<p\leq N,\quad p^2\mid at+1
 \end{array}\right\},
\]
one has
\begin{equation}\label{eq:exact-low-complete-tail}
 |\mathcal M_t(N)|\leq40.
\end{equation}
\end{lemma}

\begin{proof}
Every event on the stated interval embeds in the upper box $N_*=3\cdot10^6$.
Writing $at+1=p^2r$, every event occurs in the finite list
\[
 D<p\leq N_*,\qquad
 1\leq r\leq\left\lfloor\frac{N_*^2+1}{p^2}\right\rfloor,\qquad
 p^2r-1=at.
\]
There are $213713$ primes and exactly $27898239$ pairs $(p,r)$ in this list.
For each pair, factor $p^2r-1$ and retain every divisor
$a\leq N_*$ with $a\equiv7,18\pmod {25}$ whose complementary divisor
$t$ is also at most $N_*$.  This is an equivalence, not a relaxation.
Moreover $D^4>N_*^2+1$, so two different retained primes cannot represent
the same pair $(a,t)$.  Exact counters over $1\leq t\leq N_*$ have maximum
\[
 40\quad\hbox{at }t=10920.
\]
The prime sieve, complete quotient range, all $3063864$ retained divisor
events, deterministic $64$-bit factorization with factor-product verification,
and divisor enumeration are implemented by
\begin{center}
 \path{scripts/systematic_one_form_medium_exhaustion.cpp}.
\end{center}
Since every smaller-$N$ event is in the enumerated upper box,
\eqref{eq:exact-low-complete-tail} follows.
\end{proof}

The resulting uniform one-form bounds are as follows.

\begin{lemma}[Low one-form incidence envelopes]
\label{lem:low-one-form-incidence-envelopes}
For $1.5\cdot10^6\leq N<3\cdot10^6$, put
\begin{align*}
 E(N)&=\sum_{\substack{3\leq p\leq300\\p\ne5}}
       2\left\lceil\frac{N}{25p^2}\right\rceil+182+40,\\
 O(N)&=\sum_{\substack{3\leq p\leq300\\p\ne5}}
       2\left\lceil\frac{N}{100p^2}\right\rceil+268+40,
\end{align*}
where the sums are over primes.  If $t$ is even and non-base, then
\begin{equation}\label{eq:low-even-one-form-degree}
 \#\{a\leq N:a\equiv7,18\pmod {25},\ at+1\text{ is nonsquarefree}\}
 \leq E(N).
\end{equation}
If $t$ is odd and non-base and $r\pmod4$ is fixed, then the number of such
$a\equiv r\pmod4$ for which $at+1$ has an odd-prime-square divisor is at
most $O(N)$.  At the lower endpoint,
\[
 E(1.5\cdot10^6)=19694,
 \qquad O(1.5\cdot10^6)=5224.
\]
\end{lemma}

\begin{proof}
For a fixed odd prime $p\ne5$, the congruence $p^2\mid at+1$ has at most
one $a$-class modulo $p^2$.  Intersecting it with the two base classes gives
two progressions modulo $25p^2$, and imposing one class modulo $4$ gives two
progressions modulo $100p^2$.  These are the two ceiling sums.

On the complete upper box $a,t\leq3\cdot10^6$, exhaustive modular inversion
over $300<p\leq28540$ gives total prime-incidence maximum $268$, and maximum
$182$ when $t$ is even and non-base.  The latter is attained at $t=120$.
Lemma~\ref{lem:exact-low-complete-tail} supplies the final $40$ for
$p>28540$.  If $t$ is even, prime $2$ is inactive.  Prime $5$ is always
inactive for a non-base pivot: from $25\mid at+1$ and
$a\equiv7$ or $18\pmod {25}$ one obtains respectively
$t\equiv7$ or $18\pmod {25}$.  This proves both bounds.  The upper-box
enumeration and its exhaustive check are reproduced by
\path{scripts/systematic_low_high_prime_degree.cpp}.
\end{proof}

\begin{lemma}[First-low parity diagonal ceilings]
\label{lem:first-low-parity-diagonal}
On $1.5\cdot10^6\leq N<3\cdot10^6$, the nonsquarefree diagonal among the
non-base classes has the following uniform upper bounds:
\[
 \begin{array}{c|c}
 \text{allowed support}&\text{count}\\ \hline
 \text{unrestricted}&0.025297N\\
 \text{the two odd classes modulo }4&0.012657N\\
 \text{one odd class modulo }4&0.006337N.
 \end{array}
\]
\end{lemma}

\begin{proof}
The eventwise M\"obius--CRT diagonal enumeration on $139$ blocks applies
\eqref{eq:rational-diagonal-envelope} by exact integer
cross-multiplication.  Decimal renderings of the largest rational envelopes
are
\[
 0.0252968815411\ldots,\qquad
 0.0126562333156\ldots,\qquad
 0.0063363175457\ldots.
\]
They occur respectively on the blocks beginning at $1665621$, $1812996$,
and $1724796$.  The squarefree semigroup enumeration has $1237599$ nodes,
$4134027$ root terms, and $77041$ retained events; the harmonic and
negative-Pell tails are then added monotonically at the upper block endpoint.
The complete producer and the three asserted upward decimal ceilings are in
\path{scripts/systematic_concentrated_diagonal_blocks.cpp}.
\end{proof}

\begin{theorem}[Direct close of the first low block]
\label{thm:first-low-parity-close}
The Hall inequality~\eqref{eq:hall-defect}, and hence the exact bound in
Erd\H{o}s Problem~848, holds for every
\[
 1.5\cdot10^6\leq N<3\cdot10^6.
\]
\end{theorem}

\begin{proof}
Theorem~\ref{thm:pure-joint-tail-matching} closes the pure opposite-base
case.  Suppose therefore that a mixed Hall defect exists, put
$A=B\cup\TN(B)$, and let $A^*$ be its non-base residual part.  By
Proposition~\ref{prop:full-support-low-cluster}, the residual pivot set
contains more than
\[
 0.000249932448083113N>374
\]
members throughout the interval.  Every base member of $A$ is a common
non-neighbour of every selected residual pivot, while every member of $A^*$
belongs to the nonsquarefree diagonal.  We use these two facts in three
exhaustive parity branches.

If there are at least eight even residual pivots, select eight of them.  For
each common base non-neighbour, all eight associated forms are
nonsquarefree.  Averaging their eight indicators and then taking the union
bound over prime-square witnesses gives the common-base count at most $E(N)$.
Lemma~\ref{lem:first-low-parity-diagonal} therefore gives
\[
 |A|\leq0.025297N+E(N).
\]

Otherwise charge the at most seven even residual pivots.  Suppose first that
both odd classes modulo $4$ contain at least eight residual pivots, and select
four from each class.  Split the base variable $a$ into its four classes
modulo $4$.  When $a$ is even, prime $2$ witnesses none of the eight forms.
When $a$ is odd, it witnesses exactly one of the two four-pivot classes, so
the other four forms still require odd-prime-square witnesses.  Averaging the
required four or eight indicators on each $a$-class and applying
Lemma~\ref{lem:low-one-form-incidence-envelopes} gives common-base count at
most $4O(N)$.  Hence
\[
 |A|\leq0.012657N+7+4O(N).
\]

Finally, one odd pivot class has at most seven members.  Charge it together
with the at most seven even pivots and select eight pivots from the remaining
odd class.  On exactly one odd $a$-class modulo $4$, prime $2$ witnesses all
eight forms; that class contains at most
$2\lceil N/100\rceil$ base values.  On the other three $a$-classes prime $2$
is inactive for all selected pivots, so the averaged odd-prime count is at
most $3O(N)$.  Thus
\[
 |A|\leq0.006337N+14+2\left\lceil\frac N{100}\right\rceil+3O(N).
\]

Each right side is strictly smaller than $N/25-7/25$ on the full interval.
Indeed it is affine and increasing between ceiling jumps, so it suffices to
check the lower endpoint and the points immediately following a multiple of
$25p^2$, $100p^2$, or $100$.  The exact scan of all $21999$ such points gives
\[
\begin{array}{c|c|c}
\text{branch}&\text{worst }N&\text{count margin}\\ \hline
\text{eight even}&1500751&2359.261953\\
\text{two odd classes}&1500401&20106.184543\\
\text{one odd class}&1500401&4799.718863.
\end{array}
\]
This finite arithmetic is reproduced by
\path{scripts/systematic_low_parity_close.py}.  Every branch contradicts the
assumed Hall defect, proving the theorem.
\end{proof}

\begin{lemma}[Second-low support and one-form envelopes]
\label{lem:second-low-support-incidence}
For $3\cdot10^6\leq N<5\cdot10^6$, every mixed Hall defect has more than
\begin{equation}\label{eq:second-low-residual}
 0.00157645N
\end{equation}
non-base residual pivots.  Put
\begin{align*}
 E_5(N)&=\sum_{\substack{3\leq p\leq300\\p\ne5}}
       2\left\lceil\frac{N}{25p^2}\right\rceil+278+49,\\
 O_5(N)&=\sum_{\substack{3\leq p\leq300\\p\ne5}}
       2\left\lceil\frac{N}{100p^2}\right\rceil+383+49,
\end{align*}
where the sums are over primes.  If $t$ is even and non-base, its
nonsquarefree degree into the two base classes is at most $E_5(N)$.  If
$t$ is odd and non-base and one class of $a$ modulo $4$ is fixed, its
odd-prime-square degree into that class and the two base classes is at most
$O_5(N)$.  At the lower endpoint,
\[
 E_5(3\cdot10^6)=39211,\qquad O_5(3\cdot10^6)=10192.
\]
\end{lemma}

\begin{proof}
For the residual assertion, repeat the reverse-degree calculation of
Proposition~\ref{prop:full-support-low-cluster} with upper support
$5\cdot10^6$, cutoff $17$, split $70$, and support modulus $4902$.
For each possible number $0\leq k\leq6$ of retained prime divisors and each
$2$-adic type, take independently the smallest degree-$3$ Bonferroni density,
the largest endpoint error, and the largest full-support quadratic-residue
survivor count.  This enlarges every actual debit.  On a fixed cell
$Y=\lfloor N/70\rfloor$, the resulting normalized residual has the form
$A+B/N$, so its minimum is at a cell endpoint.  Exact rational exhaustion of
all cells and all $28$ relaxed rows has minimum
\[
 0.001576458682760064\ldots>0.00157645
\]
at $N=3000130$.  The controlling relaxed row is odd, has $k=3$, root
factor $16$, and survivor maximum $528$.  This proves
\eqref{eq:second-low-residual}; in particular the residual contains more
than $4729$ pivots throughout the interval.

The low-prime terms follow from the same two CRT progressions as in
Lemma~\ref{lem:low-one-form-incidence-envelopes}.  Exhaustive modular
inversion on the complete box $a,t\leq5\cdot10^6$ gives, over
$300<p\leq28540$, unrestricted prime-incidence maximum $383$ and even
non-base maximum $278$.  They occur at $t=168$ and $t=120$, respectively.
For $p>28540$, enumerate the complete quotient identity
\[
 p^2r-1=at,\qquad a,t\leq5\cdot10^6.
\]
There are $77829785$ prime--quotient terms and $7742361$ retained divisor
events.  Since $28540^4>(5\cdot10^6)^2+1$, distinct witness primes cannot
represent the same $(a,t)$; the exact maximum degree is $49$, at $t=9240$.
Prime $2$ is inactive for even $t$, and prime $5$ is inactive for every
non-base pivot exactly as in the preceding lemma.  The support calculation,
middle-prime inversion, and complete tail enumeration are reproduced by
\path{scripts/systematic_second_low_parity_close.py},
\path{scripts/systematic_low_high_prime_degree.cpp}, and
\path{scripts/systematic_one_form_medium_exhaustion.cpp}.
\end{proof}

\begin{lemma}[Second-low parity diagonal ceilings]
\label{lem:second-low-parity-diagonal}
On $3\cdot10^6\leq N<5\cdot10^6$, the nonsquarefree diagonal among the
non-base classes is bounded by
\[
 \begin{array}{c|c}
 \text{allowed support}&\text{count}\\ \hline
 \text{unrestricted}&0.025295N\\
 \text{the two odd classes modulo }4&0.012655N\\
 \text{one odd class modulo }4&0.006333N.
 \end{array}
\]
\end{lemma}

\begin{proof}
The complete eventwise M\"obius--CRT producer on this interval again uses
\eqref{eq:rational-diagonal-envelope} and exact integer comparisons.
Decimal renderings of its three largest rational envelopes are
\[
 0.0252947682147\ldots,\qquad
 0.0126547423852\ldots,\qquad
 0.0063325725290\ldots.
\]
The asserted decimals are strict upward ceilings.  The enumeration and the
three assertions are in
\path{scripts/systematic_concentrated_diagonal_blocks.cpp}.
\end{proof}

\begin{theorem}[Direct close of the second low block]
\label{thm:second-low-parity-close}
The Hall inequality~\eqref{eq:hall-defect}, and hence the exact bound in
Erd\H{o}s Problem~848, holds for every
\[
 3\cdot10^6\leq N<5\cdot10^6.
\]
\end{theorem}

\begin{proof}
The pure opposite-base case is already closed by
Theorem~\ref{thm:pure-joint-tail-matching}.  In a mixed defect,
Lemma~\ref{lem:second-low-support-incidence} supplies more than $4729$
non-base residual pivots, in particular more than the $22$ needed below.

Apply the same exhaustive parity trichotomy as in
Theorem~\ref{thm:first-low-parity-close}.  Eight even pivots give
\[
 |A|\leq0.025295N+E_5(N).
\]
If there are at most seven even pivots and both odd classes modulo $4$ have
at least eight members, choose four from each class; then
\[
 |A|\leq0.012655N+7+4O_5(N).
\]
Otherwise charge the at most seven even pivots and the at most seven pivots
in the smaller odd class.  More than $22$ residual pivots ensure that eight
remain in the other odd class.  Selecting them gives
\[
 |A|\leq0.006333N+14+
 2\left\lceil\frac N{100}\right\rceil+3O_5(N).
\]
Thus no parity allocation is omitted.

Each right side is affine and increasing between its ceiling jumps.
Checking the lower endpoint and every point immediately after a multiple of
$25p^2$, $100p^2$, or $100$ gives $29322$ exact candidates and
\[
\begin{array}{c|c|c}
\text{branch}&\text{worst }N&\text{count margin}\\ \hline
\text{eight even}&3000151&4901.940455\\
\text{two odd classes}&3000000&41259.720000\\
\text{one odd class}&3000801&10407.687267.
\end{array}
\]
All three margins are against the unchanged Hall target
$N/25-7/25$ and are strictly positive.  The exact scan is
\path{scripts/systematic_second_low_parity_close.py}, so every possible mixed
defect is contradicted and the block is closed.
\end{proof}

The diagonal large-square tail has a further uniform improvement which is
specific to norm $-1$.  It is stronger than the generic totally positive
unit gap used above.

\begin{lemma}[Sharp negative-Pell gap]
\label{lem:sharp-negative-pell-gap}
Let $d\geq2$ be nonsquare.  The number of positive solutions of
\[
 x^2-dy^2=-1,\qquad x\leq N,
\]
is at most
\begin{equation}\label{eq:sharp-negative-pell-count}
 P_-(N):=\left\lfloor
 1+\log_{3+2\sqrt2}N
 \right\rfloor.
\end{equation}
In particular,
\begin{equation}\label{eq:negative-pell-thirteen}
 P_-(N)=13
 \qquad(2\cdot10^9\leq N<7.67769\cdot10^9).
\end{equation}
Thus the last factor in every aggregate diagonal Pell tail on this interval
may be replaced by $13$.
\end{lemma}

\begin{proof}
Write $\alpha=x+y\sqrt d$ for a solution.  Its conjugate is
$-\alpha^{-1}$.  If $\alpha_2>\alpha_1$ are two solutions, then
$\alpha_2/\alpha_1=u+v\sqrt d$ is an integral totally positive unit of norm
$1$.  Hence $u\geq2$.  The case $u=2$ would force $dv^2=3$, hence $d=3$;
but $x^2\equiv-1\pmod3$ has no solution.  Therefore $u\geq3$, and
\[
 \frac{\alpha_2}{\alpha_1}
 =u+\sqrt{u^2-1}\geq3+2\sqrt2.
\]
Also
\[
 1+\sqrt2\leq\alpha=x+\sqrt{x^2+1}\leq(1+\sqrt2)x,
\]
so the ratio of the largest to the smallest allowed $\alpha$ is at most
$N$.  The geometric gap proves \eqref{eq:sharp-negative-pell-count}.
Finally direct squaring gives
\[
 (3+2\sqrt2)^{12}<2\cdot10^9
 <7.67769\cdot10^9<(3+2\sqrt2)^{13},
\]
which proves \eqref{eq:negative-pell-thirteen}.
\end{proof}

For an actual residual Hall defect, the mixed-base strip forces a much more
rigid pair of pivots.

\begin{proposition}[Bounded-gap residual pivots]
\label{prop:bounded-gap-residual}
Let $N\geq3.7\cdot10^{11}$, and suppose that a compatible outside set $B$
violates the Hall inequality.  Then there are distinct
\[
 b,c\in B\setminus C_N^{18},\qquad 1\leq |b-c|\leq5000.
\]
If $N\geq5.5\cdot10^{11}$, the pair may be chosen with
$d:=|b-c|\leq1381$; if also $N<N_{25}$, one may take $d\leq834$.  For this
sharper pair,
\begin{equation}\label{eq:bounded-gap-same-prime}
 |S_N(b,c)|\leq\frac{13N}{900}+2.
\end{equation}
Moreover, for every $a\in D_N(b,c)$,
\[
 \gcd(ab+1,ac+1)
\]
is a squarefree divisor of $d$.  The original Hall inequality follows if
\begin{equation}\label{eq:bounded-gap-distinct-debit}
 |\TN(B)\cap D_N(b,c)|
 \leq\frac{23N}{900}-|B|-\frac{57}{25}.
\end{equation}
Thus above $2\cdot10^{12}$ the distinct-prime contribution has determinant at
most $1381$ and linear debit $23N/900$.
\end{proposition}

\begin{proof}
Put $B^*=B\setminus C_N^{18}$ and $m=|B^*|$.  By
Corollary~\ref{cor:thin-nonbase-strip},
\[
 m>\mathcal R(N).
\]
List the elements of $B^*$ increasingly.  First suppose
$3.7\cdot10^{11}\leq N<5.5\cdot10^{11}$.  Then
\[
 m-1>0.0002N-1>\frac{N-1}{5001}.
\]
The least consecutive gap is therefore less than $5001$, hence at most
$5000$.

If $5.5\cdot10^{11}\leq N<N_{25}$, then instead
\[
 m-1>0.0012N-1>\frac{N-1}{835},
\]
so the least gap is at most $834$.

Finally suppose $N\geq N_{25}$.  The second part of
Corollary~\ref{cor:thin-nonbase-strip} gives
\[
 m>0.000724N=\frac{181N}{250000}.
\]
The least consecutive gap is at
most $(N-1)/(m-1)$.  The identity
\[
 \frac{181}{250000}-\frac1{1382}=\frac{142}{345500000}
\]
and $N\geq2\cdot10^{12}$ give
\[
 m-1>\frac{181N}{250000}-1>\frac{N-1}{1382}.
\]
Hence the least gap is strictly smaller than $1382$, and, being integral, is
at most $1381$.

If a prime contributes to $S_N(b,c)$, then its square divides $d$ and the
prime is not $5$.  Three distinct such primes would force
\[
 (2\cdot3\cdot7)^2=1764\mid d,
\]
which is impossible.  There are therefore at most two, and the two largest
residue-class budgets occur for $p=2,3$.  This proves
\eqref{eq:bounded-gap-same-prime}.

For $u=ab+1$ and $v=ac+1$ one has
\[
 cu-bv=c-b,
\]
so $(u,v)\mid d$.  If $(u,v)$ were divisible by a prime square, then $a$
would belong to $S_N(b,c)$, contrary to $a\in D_N(b,c)$.  Thus the gcd is
squarefree.

Finally $\TN(B)\subseteq\TN(\{b,c\})$ and the latter is the disjoint union
of its same- and distinct-prime parts.  If
\eqref{eq:bounded-gap-distinct-debit} holds, then
\begin{align*}
 |B|+|\TN(B)|
 &\leq |B|+\left(\frac{13N}{900}+2\right)
       +\left(\frac{23N}{900}-|B|-\frac{57}{25}\right)\\
 &=\frac N{25}-\frac7{25}\leq|\CN|.
\end{align*}
This is exactly the original Hall inequality.
\end{proof}

\begin{theorem}[Parity-pair descent to the matching threshold]
\label{thm:matching-threshold-close}
The exact bound \eqref{eq:target}, and equivalently the Hall inequality \eqref{eq:hall-defect},
holds for every $N\geq1.409\cdot10^{12}$.
\end{theorem}

\begin{proof}
Theorem~\ref{thm:8858-large-close} handles $N\geq8.858\cdot10^{12}$, so
assume
\[
 1.409\cdot10^{12}\leq N<8.858\cdot10^{12}
\]
and suppose that $B$ is a Hall defect.  Put $A=B\cup T_N(B)$ and
$A^*=B\setminus C_N^{18}$.  Corollary~\ref{cor:thin-nonbase-strip} and
Proposition~\ref{prop:bounded-gap-residual} give
\[
 |A^*|>\left(\frac2{9018009}+5\cdot10^{-7}\right)N
 \quad\text{and a pair in $A^*$ at distance at most $1400000$}.
\]

We use the same thirteen-prime set $\mathcal P$ as in
Lemma~\ref{lem:finite-two-pivot-base-sieve}.  Put
\[
 J_1=1-\prod_{p\in\mathcal P}\left(1-\frac1{p^2}\right)<0.152377,
 \qquad I_2<0.129757,\qquad I_3<0.136690.
\]
Here $I_2$ is the two-form intersection density when at most two residue
pairs coincide, and $I_3$ is the corresponding density for at most three.
If the difference of two pivots is arbitrary but at most
$N$, no more than six pairs can coincide, because
\[
 (3\cdot7\cdot11\cdot13\cdot17\cdot19\cdot23)^2
 =497704286448369>N.
\]
The corresponding finite product gives
\[
 I_6<0.146927.
\]
Write $\epsilon_1=2^{13}-1=8191$ and
$\epsilon_2=2(2^{13}-1)+(3^{13}-1)=1610704$ for the one-form and two-form
inclusion--exclusion discrepancies.

The large-prime tail can be reused with parity-sensitive root counts.  For
$k\geq2$ and root bounds $H_1,H_2$, set
\begin{align*}
 \mathfrak T_{H_1,H_2,k}(N)={}&\frac{4(0.004751)}{25}
 +\frac{4(1.25506)}{k\log(N/(k+1))}\\
 &+\frac{2(H_1+H_2)}N\left(
 \frac{2(N/25+2)}{N/(k+1)}
 +\frac{N^2+25N+1}{(N/(k+1))^2}\right).
\end{align*}
This is the bound obtained from
\eqref{eq:finite-two-pivot-envelope} with
$Y=\lfloor N/k\rfloor$: the first two terms pay the intermediate primes in
both forms and both base classes, and the last pays their transformed
progression tails.  For an odd pivot the effective root modulus is odd and,
after removal of the redundant factor $5$, is at most $5N<P_{12}$; hence
$H=2048$.  For an even pivot we retain $H=8192$.  Direct substitution at the
left endpoint gives
\[
 \mathfrak T_{8192,8192,168}(N)<0.002732802,
 \qquad
 \mathfrak T_{2048,2048,269}(N)<0.002018195.
\]
Both displayed envelopes decrease.  At the same endpoint,
\[
 \Delta_{25}^{(0)}<0.001098525,\quad
 \Delta_{50}^{(0)}<0.001098460,\quad
 \Delta_{100}^{(0)}<0.001098428.
\]

Let $e$ be the number of even elements of $A^*$.  First suppose
\[
 e-1>\frac{N-1}{9018009}.
\]
Two even pivots then lie at distance at most $9018008$.  The prime $2$ cannot
witness either linear form, and at most three odd-prime residue pairs coincide.
The unrestricted aggregate diagonal estimate and the two-pivot base sieve
give, after division by $N$,
\[
 \frac{|A|}{N}<
 \frac{23}{25}C_{\rm quad}+\Delta_{25}^{(0)}
 +\frac{2I_3}{25}+\frac{2\epsilon_2}{N}
 +\mathfrak T_{8192,8192,168}<0.039928.
\]

It remains to suppose
\[
 e\leq\frac{N-1}{9018009}+1<\frac N{9018009}+1.
\]
Let $O$ be the odd part of $A^*$.  If $O$ meets both odd residue classes
modulo $4$, choose one pivot from each.  On each of the two base classes
modulo $25$, two of the four progressions modulo $100$ have an automatic
factor $4$ for exactly one pivot, and two have none.  Thus the small-prime
base density is at most $4(J_1+I_6)/100$.  Charging the even part by $e$ and
using the odd diagonal estimate gives
\[
 \frac{|A|}{N}<
 \frac1{9018009}+\frac1N
 +\frac{23}{50}C_{\rm quad}+\Delta_{50}^{(0)}
 +\frac{4(J_1+I_6)}{100}
 +\frac{4\epsilon_1+4\epsilon_2}{N}
 +\mathfrak T_{2048,2048,269}<0.027674.
\]

Finally suppose that $O$ lies in one odd residue class modulo $4$.  Since
\[
 |O|>|A^*|-e>
 N\left(\frac1{9018009}+5\cdot10^{-7}\right)-2,
\]
two elements of $O$ lie at distance at most $9018008$; the extra
$5\cdot10^{-7}N$ absorbs the two endpoint losses.  At most three odd-prime
residue pairs therefore coincide.  In each base class one progression modulo
$100$ has an automatic factor $4$ for both pivots, while the other three are
governed by the two-form density $I_2$.  The restricted diagonal estimate
now gives
\[
 \frac{|A|}{N}<
 \frac1{9018009}+\frac1N
 +\frac{23}{100}C_{\rm quad}+\Delta_{100}^{(0)}
 +\frac2{100}+\frac{6I_3}{100}
 +\frac{2+6\epsilon_2}{N}
 +\mathfrak T_{2048,2048,269}<0.037615.
\]

All three bounds are strictly below $1/25-7/(25N)$ on the stated range,
contradicting $|A|>|C_N|$.  This closes the interval and proves the theorem.
\end{proof}

\begin{theorem}[Primorial matching descent]
\label{thm:550b-close}
The exact bound \eqref{eq:target}, and equivalently the Hall inequality \eqref{eq:hall-defect},
holds for every $N\geq5.5\cdot10^{11}$.
\end{theorem}

\begin{proof}
It remains only to treat
\[
 5.5\cdot10^{11}\leq N<1.409\cdot10^{12}.
\]
Suppose $B$ is a Hall defect and put $A=B\cup T_N(B)$ and
$A^*=B\setminus C_N^{18}$.  Corollary~\ref{cor:thin-nonbase-strip} gives
\[
 |A^*|>0.0012N.
\]
On this interval an even pivot has root count at most $4096$, since
$5N<8P_{11}=29682952539240$; an odd pivot has root count at most $2048$,
since $5N<P_{12}=152125131763605$.

Retain $J_1,I_2,I_6,\epsilon_1,\epsilon_2$ and the tail function
$\mathfrak T$ from the preceding proof.  At the left endpoint,
\[
 \mathfrak T_{4096,4096,154}<0.002958674,
 \qquad
 \mathfrak T_{2048,2048,194}<0.002515972,
\]
and
\[
 \Delta_{25}^{(0)}<0.001465437,\quad
 \Delta_{50}^{(0)}<0.001465317,\quad
 \Delta_{100}^{(0)}<0.001465258.
\]
All these envelopes decrease.

Let $e$ count the even elements of $A^*$.  If
$e-1>(N-1)/53361$, two even pivots lie within $53360$ and at most two
odd-prime residue pairs coincide.  Therefore
\[
 \frac{|A|}{N}<
 \frac{23}{25}C_{\rm quad}+\Delta_{25}^{(0)}
 +\frac{2I_2}{25}+\frac{2\epsilon_2}{N}
 +\mathfrak T_{4096,4096,154}<0.039970.
\]

Otherwise $e<N/53361+1$.  If the odd part of $A^*$ meets both odd classes
modulo $4$, choose one pivot from each.  As before at most six small-prime
residue pairs can coincide, and
\[
 \frac{|A|}{N}<
 \frac1{53361}+\frac1N
 +\frac{23}{50}C_{\rm quad}+\Delta_{50}^{(0)}
 +\frac{4(J_1+I_6)}{100}
 +\frac{4\epsilon_1+4\epsilon_2}{N}
 +\mathfrak T_{2048,2048,194}<0.028564.
\]
If instead the odd part lies in one class modulo $4$, its size is greater
than $(0.0012-1/53361)N-2$.  Thus two odd pivots lie within $53360$, and
\[
 \frac{|A|}{N}<
 \frac1{53361}+\frac1N
 +\frac{23}{100}C_{\rm quad}+\Delta_{100}^{(0)}
 +\frac2{100}+\frac{6I_2}{100}
 +\frac{2+6\epsilon_2}{N}
 +\mathfrak T_{2048,2048,194}<0.038093.
\]
Each bound is below $1/25-7/(25N)$, a contradiction.  The preceding theorem
handles all larger $N$.
\end{proof}

\begin{theorem}[Two-adic primorial descent]
\label{thm:370b-close}
The exact bound \eqref{eq:target}, and equivalently the Hall inequality \eqref{eq:hall-defect},
holds for every $N\geq3.7\cdot10^{11}$.
\end{theorem}

\begin{proof}
By Theorem~\ref{thm:550b-close}, it remains to consider
\[
 3.7\cdot10^{11}\leq N<5.5\cdot10^{11}.
\]
Suppose that $B$ is a Hall defect, put $A=B\cup T_N(B)$ and
$A^*=B\setminus C_N^{18}$, and let $e$ be the number of even elements of
$A^*$.  Corollary~\ref{cor:thin-nonbase-strip} gives
\begin{equation}\label{eq:380-nonbase-mass}
 |A^*|>0.0002N.
\end{equation}

We first determine the required root counts.  Let
\[
 m=2^\nu m_0\leq5N,\qquad m_0\ \text{odd},
\]
be the effective modulus in Lemma~\ref{lem:primorial-progression}.
Since
\[
 P_{11}=3710369067405,\qquad P_{12}=152125131763605,
\]
an odd pivot has root count at most $2048$.  For an even pivot, respectively,
\[
 \begin{array}{c|c|c}
 \nu_2(b)&m_0\text{ is at most}&\rho(m)\text{ is at most}\\ \hline
 1&5N/2<P_{11}&1024\\
 2&5N/4<P_{11}&2048\\
 \geq3&5N/8<P_{11}&4096.
 \end{array}
\]
Indeed these are exactly the factors $c_1=1,c_2=2,c_\nu=4$ in the proof of
Lemma~\ref{lem:primorial-progression}, multiplied by the largest possible
$2^{\omega(m_0)}$.

Retain $J_1,I_2,I_6,\epsilon_1,\epsilon_2$ and $\mathfrak T$ from the proof of
Theorem~\ref{thm:matching-threshold-close}.  At $N=3.7\cdot10^{11}$,
\[
 \begin{array}{lll}
 \Delta_{25}^{(0)}<0.001654479,&
 \Delta_{50}^{(0)}<0.001654326,&
 \Delta_{100}^{(0)}<0.001654249,\\
 \Delta_8^{(0)}<0.002795228,&
 \mathfrak T_{1024,1024,218}<0.002375113,&
 \mathfrak T_{2048,2048,172}<0.002781711,\\
 &&\mathfrak T_{4096,4096,136}<0.003291527.
 \end{array}
\]
All these envelopes decrease on the interval.

First suppose $e>3N/100000$.  Partition the even pivots as
\[
 E_1=\{b:\nu_2(b)=1\},\qquad E_2=\{b:\nu_2(b)=2\},\qquad
 E_3=\{b:\nu_2(b)\geq3\}.
\]
If
\[
 |E_i|-1>\frac{N-1}{100001}
 \qquad(i=1\text{ or }2),
\]
two pivots of that type have distance at most $100000$.  The prime $2$ cannot
witness either associated linear form.  Moreover at most two odd-prime
residue pairs can coincide: three would make the even difference divisible by
\[
 2(3\cdot7\cdot11)^2=106722.
\]
The unrestricted diagonal estimate and the two-pivot base sieve therefore
give, for $i=1$ and $i=2$, respectively,
\begin{align*}
 \frac{|A|}{N}
 &<\frac{23}{25}C_{\rm quad}+\Delta_{25}^{(0)}
   +\frac{2I_2}{25}+\frac{2\epsilon_2}{N}
   +\mathfrak T_{1024,1024,218}<0.039578,\\
 \frac{|A|}{N}
 &<\frac{23}{25}C_{\rm quad}+\Delta_{25}^{(0)}
   +\frac{2I_2}{25}+\frac{2\epsilon_2}{N}
   +\mathfrak T_{2048,2048,172}<0.039985.
\end{align*}

It remains in the even case to suppose that neither $E_1$ nor $E_2$ is
dense.  Then
\[
 |E_1|+|E_2|\leq\frac{2(N-1)}{100001}+2
\]
and, using $e>3N/100000$ and $N\geq3.7\cdot10^{11}$,
\[
 |E_3|-1>\frac{3N}{100000}-\frac{2(N-1)}{100001}-3
          >\frac{N-1}{100001}.
\]
Thus two members of $E_3$ again lie within $100000$.  Apart from the raw
charge for $E_1\cup E_2$, every member of $A^*$ is odd or divisible by $8$,
so Corollary~\ref{cor:low-two-adic-diagonal} applies.  Consequently
\begin{align*}
 \frac{|A|}{N}<{}&\frac{23}{40}C_{\rm quad}+\Delta_8^{(0)}
  +\frac{2}{100001}+\frac2N
  +\frac{2I_2}{25}+\frac{2\epsilon_2}{N}
  +\mathfrak T_{4096,4096,136}\\
 &<0.032221.
\end{align*}

We may therefore assume $e\leq3N/100000$.  Let $O$ be the odd part of
$A^*$.  By \eqref{eq:380-nonbase-mass},
\[
 |O|>0.00017N.
\]
If $O$ meets both odd classes modulo $4$, select one pivot from each.  At most
six small-prime residue pairs coincide, just as in
Theorem~\ref{thm:matching-threshold-close}; charging all even pivots directly
gives
\begin{align*}
 \frac{|A|}{N}<{}&\frac3{100000}
  +\frac{23}{50}C_{\rm quad}+\Delta_{50}^{(0)}
  +\frac{4(J_1+I_6)}{100}
  +\frac{4\epsilon_1+4\epsilon_2}{N}
  +\mathfrak T_{2048,2048,172}\\
 &<0.029036.
\end{align*}

Finally suppose that $O$ lies in one odd class modulo $4$.  The inequality
$|O|>17N/100000$ implies, throughout the interval, that two of its elements
lie at distance at most $5882$.  Their difference is even, so at most two
odd-prime residue pairs coincide.  In each base class one progression modulo
$100$ has an automatic factor $4$ for both pivots, while the other three are
governed by $I_2$.  Hence
\begin{align*}
 \frac{|A|}{N}<{}&\frac3{100000}
  +\frac{23}{100}C_{\rm quad}+\Delta_{100}^{(0)}
  +\frac2{100}+\frac{6I_2}{100}
  +\frac{2+6\epsilon_2}{N}
  +\mathfrak T_{2048,2048,172}\\
 &<0.038568.
\end{align*}

These five cases exhaust the parity and two-adic possibilities.  Every bound
is strictly below $1/25-7/(25N)$ on the stated interval, contradicting
$|A|>|C_N|$.  Thus the interval closes inside the original Hall inequality.
\end{proof}

\begin{theorem}[Three-class matching descent]
\label{thm:34375b-close}
The exact bound \eqref{eq:target}, and equivalently the Hall inequality \eqref{eq:hall-defect},
holds for every $N\geq3.4375\cdot10^{11}$.
\end{theorem}

\begin{proof}
Theorem~\ref{thm:370b-close} leaves only
\[
 3.4375\cdot10^{11}\leq N<3.7\cdot10^{11}.
\]
Suppose that $B$ is a Hall defect, put $A=B\cup T_N(B)$ and
$A^*=B\setminus C_N^{18}$, and let $e$ count the even elements of $A^*$.
Corollary~\ref{cor:thin-nonbase-strip} gives
\begin{equation}\label{eq:344-nonbase-mass}
 |A^*|>2.5\cdot10^{-7}N.
\end{equation}
On this interval every odd pivot has root count at most $1024$, because its
effective modulus is odd and at most $5N<P_{11}$.  The three even valuation
classes retain the respective bounds $1024,2048,4096$ from the preceding
theorem.

In addition to the constants already defined, let $I_1$ denote the two-form
small-prime intersection density when at most one residue pair coincides.  Its
worst case is the coincident prime $3$, and direct rational multiplication in
the product defining $I_J$ gives
\[
 I_1<0.112570.
\]
At the left endpoint,
\[
 \begin{array}{lll}
 \Delta_{25}^{(0)}<0.001692156,&
 \Delta_{50}^{(0)}<0.001691995,&
 \Delta_{100}^{(0)}<0.001691915,\\
 \Delta_W<0.002859240,&
 \mathfrak T_{1024,1024,211}<0.002417843,&
 \mathfrak T_{4096,4096,132}<0.003358600.
 \end{array}
\]
Every displayed envelope decreases on the interval.

First suppose $e>0.004N$.  One of the three even valuation classes contains
more than $N/750$ elements, so two of its members lie at distance at most
$750$.  Two coincident odd-prime residue pairs would make this even difference
divisible by
\[
 2(3\cdot7)^2=882;
\]
hence at most one pair coincides.  Using the worst even root count $4096$ gives
\[
 \frac{|A|}{N}<
 \frac{23}{25}C_{\rm quad}+\Delta_{25}^{(0)}
 +\frac{2I_1}{25}+\frac{2\epsilon_2}{N}
 +\mathfrak T_{4096,4096,132}<0.039225.
\]

We may assume $e\leq0.004N$.  Let $O$ be the odd part of $A^*$.  If $O$
meets both odd classes modulo $4$, choose one pivot from each.  Charging the
even part directly and using the cross-class calculation from
Theorem~\ref{thm:matching-threshold-close} gives
\begin{align*}
 \frac{|A|}{N}<{}&0.004+\frac{23}{50}C_{\rm quad}+\Delta_{50}^{(0)}
 +\frac{4(J_1+I_6)}{100}
 +\frac{4\epsilon_1+4\epsilon_2}{N}
 +\mathfrak T_{1024,1024,211}\\
 &<0.032681.
\end{align*}

It remains to suppose that $O$ occupies at most one odd class modulo $4$.
If $e>10^{-7}N$, two even pivots lie at distance at most $10^7$.  Four
coincident odd-prime residue pairs would force their even difference to be at
least
\[
 2(3\cdot7\cdot11\cdot13)^2=18036018,
\]
so at most three pairs coincide.  Choose an odd class
$\varepsilon\in\{1,3\}$ not occupied by $O$.  Then $A^*\subseteq W_\varepsilon$,
and Lemma~\ref{lem:three-mod-four-diagonal} gives
\[
 \frac{|A|}{N}<
 \frac{69}{100}C_{\rm quad}+\Delta_W
 +\frac{2I_3}{25}+\frac{2\epsilon_2}{N}
 +\mathfrak T_{4096,4096,132}<0.036032.
\]

Finally let $e\leq10^{-7}N$.  By \eqref{eq:344-nonbase-mass}, the single odd
class contains more than $1.5\cdot10^{-7}N$ elements.  Two lie at distance at
most $6699999$, again allowing at most three coincident odd-prime residue
pairs.  The one-class odd calculation therefore gives
\begin{align*}
 \frac{|A|}{N}<{}&10^{-7}
 +\frac{23}{100}C_{\rm quad}+\Delta_{100}^{(0)}
 +\frac2{100}+\frac{6I_3}{100}
 +\frac{2+6\epsilon_2}{N}
 +\mathfrak T_{1024,1024,211}\\
 &<0.038630.
\end{align*}

The four alternatives are exhaustive, and every upper bound is strictly less
than $1/25-7/(25N)$.  This contradicts the original Hall defect and closes the
whole interval without introducing another cut.
\end{proof}

\begin{theorem}[Five-free Euler-product descent]
\label{thm:18263b-close}
The exact bound \eqref{eq:target}, and equivalently the Hall inequality \eqref{eq:hall-defect},
holds for every $N\geq1.8263\cdot10^{11}$.
\end{theorem}

\begin{proof}
Theorem~\ref{thm:34375b-close} leaves the interval
\[
 1.8263\cdot10^{11}\leq N<3.4375\cdot10^{11}.
\]
Suppose that $B$ is a Hall defect, put $A=B\cup T_N(B)$ and
$A^*=B\setminus C_N^{18}$, and split its even part into
\[
 E_1=\{b:\nu_2(b)=1\},\qquad E_2=\{b:\nu_2(b)=2\},\qquad
 E_3=\{b:\nu_2(b)\geq3\}.
\]
Write $e=|E_1|+|E_2|+|E_3|$.  Corollary~\ref{cor:thin-nonbase-strip} gives
\begin{equation}\label{eq:18263-nonbase-mass}
 |A^*|>2.5\cdot10^{-7}N.
\end{equation}
As in the preceding theorem, odd pivots have root count at most $1024$, and
the three even classes have respective root bounds $1024,2048,4096$.  Retain
$I_1<0.112570$ and all the other finite-sieve constants defined above.

At $N=1.8263\cdot10^{11}$,
\[
 \begin{array}{lll}
 \Delta_{25}^{(0)}<0.002053324,&
 \Delta_{50}^{(0)}<0.002053084,&
 \Delta_{100}^{(0)}<0.002052964,\\
 \Delta_W<0.003472940,&
 \Delta_8^{(0)}<0.003472880,&
 \mathfrak T_{1024,1024,172}<0.002836481,\\
 \mathfrak T_{2048,2048,136}<0.003359437,&
 \mathfrak T_{4096,4096,108}<0.004015367.
 \end{array}
\]
All these envelopes decrease.

If $|E_1|>0.0012N$, two members of $E_1$ lie at distance at most $834$.
Since two coincident odd-prime residue pairs would force the even difference
to be at least $2(3\cdot7)^2=882$, at most one pair coincides.  Hence
\[
 \frac{|A|}{N}<\frac{23}{25}C_{\rm quad}+\Delta_{25}^{(0)}
 +\frac{2I_1}{25}+\frac{2\epsilon_2}{N}
 +\mathfrak T_{1024,1024,172}<0.039072.
\]

Suppose next that $|E_1|\leq0.0012N$ but $|E_2|>0.0012N$.  Choose the same
kind of pair in $E_2$.  After charging $E_1$ directly, every remaining
non-base diagonal element lies in $W_2$.  Thus
\[
 \frac{|A|}{N}<0.0012+\frac{69}{100}C_{\rm quad}+\Delta_W
 +\frac{2I_1}{25}+\frac{2\epsilon_2}{N}
 +\mathfrak T_{2048,2048,136}<0.035925.
\]

Now suppose $|E_1|,|E_2|\leq0.0012N$ but $|E_3|>0.0012N$.  A pair in
$E_3$ again has distance at most $834$ and at most one coincident residue
pair.  Apart from the raw charge for $E_1\cup E_2$, the diagonal lies in
$V_8$, so
\[
 \frac{|A|}{N}<0.0024+\frac{23}{40}C_{\rm quad}+\Delta_8^{(0)}
 +\frac{2I_1}{25}+\frac{2\epsilon_2}{N}
 +\mathfrak T_{4096,4096,108}<0.034636.
\]

We may therefore assume $e\leq0.0036N$.  Let $O$ be the odd part of $A^*$.
If $O$ meets both odd classes modulo $4$, the cross-class pair calculation,
with the even part charged directly, gives
\[
 \frac{|A|}{N}<0.0036+\frac{23}{50}C_{\rm quad}+\Delta_{50}^{(0)}
 +\frac{4(J_1+I_6)}{100}
 +\frac{4\epsilon_1+4\epsilon_2}{N}
 +\mathfrak T_{1024,1024,172}<0.033077.
\]

It remains that $O$ occupies at most one odd class.  If $e>10^{-7}N$, two
even pivots lie within $10^7$, so at most three odd-prime residue pairs
coincide.  Choosing an unoccupied odd class $\varepsilon$ gives
$A^*\subseteq W_\varepsilon$, and therefore
\[
 \frac{|A|}{N}<\frac{69}{100}C_{\rm quad}+\Delta_W
 +\frac{2I_3}{25}+\frac{2\epsilon_2}{N}
 +\mathfrak T_{4096,4096,108}<0.037311.
\]

Finally, if $e\leq10^{-7}N$, equation~\eqref{eq:18263-nonbase-mass} puts more
than $1.5\cdot10^{-7}N$ pivots in the single odd class.  Two lie within
$6699999$, allowing at most three coincident residue pairs.  Hence
\[
 \frac{|A|}{N}<10^{-7}+\frac{23}{100}C_{\rm quad}+\Delta_{100}^{(0)}
 +\frac2{100}+\frac{6I_3}{100}
 +\frac{2+6\epsilon_2}{N}
 +\mathfrak T_{1024,1024,172}<0.039434.
\]

The six cases exhaust all possible mass allocations.  Each bound is strictly
below $1/25-7/(25N)$, contradicting the same original Hall inequality.
\end{proof}

\begin{theorem}[Root-product correlated descent]
\label{thm:1001b-close}
The exact bound \eqref{eq:target}, and equivalently the Hall inequality \eqref{eq:hall-defect},
holds for every $N\geq1.001\cdot10^{11}$.
\end{theorem}

\begin{proof}
Theorem~\ref{thm:18263b-close} leaves
\[
 1.001\cdot10^{11}\leq N<1.8263\cdot10^{11}.
\]
Suppose that $B$ is a Hall defect and use the decomposition
$A^*=O\sqcup E_1\sqcup E_2\sqcup E_3$ from the preceding theorem.  Put
$e=|E_1|+|E_2|+|E_3|$.  The correlated degree and
Corollary~\ref{cor:thin-nonbase-strip} give
\begin{equation}\label{eq:1001-nonbase-mass}
 |A^*|>0.00032N.
\end{equation}
The parity-sensitive root bounds remain $1024$ for odd pivots and
$1024,2048,4096$ for $E_1,E_2,E_3$.

At the left endpoint,
\[
 \begin{array}{lll}
 \Delta_{25}^{(0)}<0.002467727,&
 \Delta_{50}^{(0)}<0.002467376,&
 \Delta_{100}^{(0)}<0.002467200,\\
 \Delta_W<0.004178034,&
 \Delta_8^{(0)}<0.004177947,&
 \mathfrak T_{1024,1024,142}<0.003333255,\\
 \mathfrak T_{2048,2048,112}<0.003981577,&
 \mathfrak T_{4096,4096,88}<0.004795138.
 \end{array}
\]
All envelopes decrease.

The first four mass cases are identical to those in the preceding theorem:
if $|E_1|>0.0012N$, if instead $|E_2|>0.0012N$, if instead
$|E_3|>0.0012N$, or if all three are at most $0.0012N$ and $O$ meets both odd
classes, the respective bounds are
\begin{align*}
 \frac{|A|}{N}&<\frac{23}{25}C_{\rm quad}+\Delta_{25}^{(0)}
 +\frac{2I_1}{25}+\frac{2\epsilon_2}{N}
 +\mathfrak T_{1024,1024,142}<0.039998,\\
 \frac{|A|}{N}&<0.0012+\frac{69}{100}C_{\rm quad}+\Delta_W
 +\frac{2I_1}{25}+\frac{2\epsilon_2}{N}
 +\mathfrak T_{2048,2048,112}<0.037267,\\
 \frac{|A|}{N}&<0.0024+\frac{23}{40}C_{\rm quad}+\Delta_8^{(0)}
 +\frac{2I_1}{25}+\frac{2\epsilon_2}{N}
 +\mathfrak T_{4096,4096,88}<0.036136,\\
 \frac{|A|}{N}&<0.0036+\frac{23}{50}C_{\rm quad}+\Delta_{50}^{(0)}
 +\frac{4(J_1+I_6)}{100}+\frac{4\epsilon_1+4\epsilon_2}{N}
 +\mathfrak T_{1024,1024,142}<0.034017.
\end{align*}

It remains that $O$ occupies at most one odd class.  If $e>10^{-5}N$, two
even pivots lie within $100000$, so at most two odd-prime residue pairs
coincide.  Choosing an unoccupied odd class gives
\[
 \frac{|A|}{N}<\frac{69}{100}C_{\rm quad}+\widetilde\Delta_W
 +\frac{2I_2}{25}+\frac{2\epsilon_2}{N}
 +\mathfrak T_{4096,4096,88}<0.038256.
\]
If $e\leq10^{-5}N$, equation~\eqref{eq:1001-nonbase-mass} leaves more than
$0.00031N$ pivots in the single odd class.  Two lie within $3226$, again
allowing at most two coincident residue pairs, and
\[
 \frac{|A|}{N}<10^{-5}+\frac{23}{100}C_{\rm quad}+\Delta_{100}^{(0)}
 +\frac2{100}+\frac{6I_2}{100}
 +\frac{2+6\epsilon_2}{N}
 +\mathfrak T_{1024,1024,142}<0.039983.
\]
All six exhaustive bounds are below $1/25-7/(25N)$, closing the interval in
the original Hall balance.
\end{proof}

The last descent uses a sharper elementary payment for the same prime tail.
Exact rational summation over the primes through $10^5$, followed by the
integral bound for all remaining integers, gives
\[
 \sum_{47<p\leq10^5}\frac1{p^2}<0.003876557,
 \qquad
 \sum_{n>10^5}\frac1{n^2}<10^{-5}.
\]
Consequently
\[
 \sum_{p>47}\frac1{p^2}<0.003887=: \tau_\sharp.
\]
Define $\mathfrak T^\sharp_{H_1,H_2,k}$ by the formula for $\mathfrak T$ with
$0.004751$ replaced by $\tau_\sharp$.

For the intermediate-prime count, Dusart proved the explicit estimate
\cite{DusartPrimes}
\[
 \pi(y)\leq\frac{y}{\log y-1.1}\qquad(y>60184)
 \tag{D}\label{eq:dusart-pi}
\]
If $Y=\lfloor N/k\rfloor$, then
$N/(k+1)\leq Y\leq N/k$ in every use below.  Consequently define
\begin{align*}
 \mathfrak T^{D,\sharp}_{H_1,H_2,k}(N)={}&\frac{4\tau_\sharp}{25}
 +\frac{4}{k\bigl(\log(N/(k+1))-1.1\bigr)}\\
 &+\frac{2(H_1+H_2)}N\left(
 \frac{2(N/25+2)}{N/(k+1)}
 +\frac{N^2+25N+1}{(N/(k+1))^2}\right).
\end{align*}
This is the same two-pivot union bound as $\mathfrak T^\sharp$, with only
the prime-count payment replaced by \eqref{eq:dusart-pi}.  For the three
values of $k$ used below, already at the left endpoint
$N/(k+1)>6\cdot10^8$, so the hypothesis in \eqref{eq:dusart-pi} holds with
ample room.

\begin{theorem}[Short-block correlated descent]
\label{thm:670b-close}
The exact bound \eqref{eq:target}, and equivalently the Hall inequality \eqref{eq:hall-defect},
holds for every $N\geq6.7\cdot10^{10}$.
\end{theorem}

\begin{proof}
It remains after Theorem~\ref{thm:1001b-close} to treat
\[
 6.7\cdot10^{10}\leq N<1.001\cdot10^{11}.
\]
Use the same decomposition $A^*=O\sqcup E_1\sqcup E_2\sqcup E_3$.
Corollary~\ref{cor:thin-nonbase-strip} gives
\begin{equation}\label{eq:670-nonbase-mass}
 |A^*|>0.00127N.
\end{equation}
At the left endpoint,
\[
 \begin{array}{lll}
 \widetilde\Delta_{25}<0.002514345,&
 \widetilde\Delta_{50}<0.002513892,&
 \widetilde\Delta_{100}<0.002513665,\\
 \widetilde\Delta_W<0.004175714,&
 \widetilde\Delta_8<0.004175602,&
 \mathfrak T^{D,\sharp}_{1024,1024,117}<0.003267694,\\
 \mathfrak T^{D,\sharp}_{2048,2048,93}<0.003933514,&
 \mathfrak T^{D,\sharp}_{4096,4096,73}<0.004768991.
 \end{array}
\]
All envelopes decrease.

The first four allocation cases from Theorem~\ref{thm:1001b-close}, now using
Lemma~\ref{lem:short-block-diagonal} and $\mathfrak T^{D,\sharp}$, give
\[
 \begin{array}{c|cccc}
 \text{case}&E_1&E_2&E_3&\text{two odd classes}\\ \hline
 |A|/N&<0.039995&<0.037233&<0.036123&<0.034030.
 \end{array}
\]
respectively.  For clarity, these are the cases $|E_1|>0.0012N$, then
$|E_2|>0.0012N$, then $|E_3|>0.0012N$, and finally both odd classes occupied
after all three even classes are sparse; their formulas are unchanged.

It remains that the odd part occupies at most one odd class.  If
$e>10^{-5}N$, an even pair lies within $100000$, at most two residue pairs
coincide, and
\[
 \frac{|A|}{N}<\frac{69}{100}C_{\rm quad}+\Delta_W
 +\frac{2I_2}{25}+\frac{2\epsilon_2}{N}
 +\mathfrak T^{D,\sharp}_{4096,4096,73}<0.038243.
\]
Otherwise \eqref{eq:670-nonbase-mass} puts more than $0.00126N$ pivots in
the single odd class.  Two lie within $794$.  Their difference is divisible
by $4$; if two residue pairs coincided, then $(3\cdot7)^2$ would also divide
it, forcing a difference of at least $1764$.  Thus at most one residue pair
coincides, and
\[
 \frac{|A|}{N}<10^{-5}+\frac{23}{100}C_{\rm quad}+\widetilde\Delta_{100}
 +\frac2{100}+\frac{6I_1}{100}
 +\frac{2+6\epsilon_2}{N}
 +\mathfrak T^{D,\sharp}_{1024,1024,117}<0.038980.
\]
All six bounds directly contradict the original Hall balance.
\end{proof}

For the last short interval put
\[
 h_\dagger:=\sum_{\substack{1\leq a\leq6293\\(a,43890)=1}}\frac1a
 <2.302594,
 \qquad
 K_\dagger(R)=2\eta_*Rh_\dagger+128\sqrt R.
\]
This is again an exact finite rational inequality.  For
$6.45\cdot10^{10}\leq N\leq6.7\cdot10^{10}$ define
\begin{align*}
 \Delta_q^\dagger&=\widetilde\Delta_q
 -\frac{3K_*(r_N+1)}N+\frac{3K_\dagger(r_N+1)}N,\\
 \Delta_W^\dagger&=\widetilde\Delta_W
 -\frac{6K_*(r_N+1)}N+\frac{6K_\dagger(r_N+1)}N,\\
 \Delta_8^\dagger&=\widetilde\Delta_8
 -\frac{6K_*(r_N+1)}N+\frac{6K_\dagger(r_N+1)}N.
\end{align*}
Since $\sqrt{r_N+1}<6293$ on this interval, the proof of
Lemma~\ref{lem:short-block-diagonal} gives the same three diagonal conclusions
with the daggered errors.

There is also a much smaller rounding budget for the one-coincidence finite
sieve.  For a vector $u=(u_p)_{p\in\mathcal P}$, write $e_j(u)$ for its
$j$th elementary symmetric function.  Put $x_p=p^{-2}$ and, for a possible
coincident prime $p_0$, put
\[
 z_p(p_0)=\begin{cases}p^{-2},&p=p_0,\\2p^{-2},&p\ne p_0.\end{cases}
\]
Define
\[
 L_5=\sum_{j=0}^5(-1)^je_j(x),\qquad
 U_4(p_0)=\sum_{j=0}^4(-1)^je_j(z(p_0)).
\]

\begin{lemma}[Truncated one-coincidence sieve]
\label{lem:bonferroni-one-coincidence}
If at most one small-prime residue pair coincides, the two-form intersection
density in one base progression is at most
\[
 \widehat I_1:=1-2L_5+\max_{p_0\in\mathcal P}U_4(p_0)
 <0.112569312,
\]
and its endpoint discrepancy is at most
\[
 \widehat\epsilon_1=18824.
\]
Thus in two base classes the corresponding contribution is at most
\[
 \frac{2\widehat I_1}{25}N+2\widehat\epsilon_1.
\]
\end{lemma}

\begin{proof}
For either individual form, the degree-$5$ Bonferroni polynomial is a lower
bound for avoiding all thirteen forbidden residues.  For the union of the two
residue sets at each prime, the degree-$4$ polynomial is an upper bound for
avoiding them.  Subtracting the two individual lower bounds from this joint
upper bound proves the displayed density.  Exact rational evaluation over the
thirteen possible $p_0$ shows that the maximum occurs at $p_0=3$ and is below
$0.112569312$; the no-coincidence case is smaller.

The two individual degree-$5$ polynomials contain
$\sum_{j=1}^5\binom{13}{j}=2379$ nonconstant residue counts each.  In the
zero-coincidence joint case, which has the largest endpoint count, the numbers
of residue classes in degrees $0,1,2,3,4$ are
$1,26,312,2288,11440$, so there are $14066$ nonconstant counts.  Charging one
endpoint error to each gives the uniform budget
\[
 2(2379)+14066=18824.
\]
The progression has length at most $N/25+1$, and repeating it in the two base
classes proves the last assertion.
\end{proof}

\begin{lemma}[Truncated two-coincidence sieve]
\label{lem:bonferroni-two-coincidence}
If at most two small-prime residue pairs coincide, the same truncation gives
\[
 \widehat I_2<0.129756863,\qquad
 \widehat\epsilon_2=18824.
\]
In six classes modulo $100$, the resulting contribution is at most
\[
 \frac{6\widehat I_2}{100}N+6\widehat\epsilon_2.
\]
\end{lemma}

\begin{proof}
Use the same individual degree-$5$ lower bounds as in
Lemma~\ref{lem:bonferroni-one-coincidence}.  In the joint degree-$4$ upper
bound let two local residue counts be $1$ and the other eleven be $2$.  Exact
rational evaluation over the possible coincident pairs is maximal for
$\{3,7\}$ and gives the displayed density.  The joint numbers of residue
classes in degrees $0,1,2,3,4$ are
\[
 1,24,265,1782,8140.
\]
Thus the exact-two-coincidence endpoint count is
\[
 2(2379)+(24+265+1782+8140)=14969<14972.
\]
For the stated ``at most two'' assertion, however, zero coincidences have the
larger uniform endpoint budget $18824$ from the preceding lemma.  Repeating
in the six base progressions modulo $100$ proves the claim.
\end{proof}


\begin{theorem}[Bonferroni short-block descent]
\label{thm:645b-close}
The exact bound \eqref{eq:target}, and equivalently the Hall inequality \eqref{eq:hall-defect},
holds for every $N\geq6.45\cdot10^{10}$.
\end{theorem}

\begin{proof}
Theorem~\ref{thm:670b-close} leaves only
\[
 6.45\cdot10^{10}\leq N<6.7\cdot10^{10}.
\]
At the left endpoint the daggered diagonal errors and optimized tails satisfy
\[
 \begin{array}{lll}
 \Delta_{25}^\dagger<0.002526953,&
 \Delta_{50}^\dagger<0.002526489,&
 \Delta_{100}^\dagger<0.002526257,\\
 \Delta_W^\dagger<0.004191286,&
 \Delta_8^\dagger<0.004191171,&
 \mathfrak T^{D,\sharp}_{1024,1024,116}<0.003304054,\\
 \mathfrak T^{D,\sharp}_{2048,2048,92}<0.003979107,&
 \mathfrak T^{D,\sharp}_{4096,4096,72}<0.004826133.
 \end{array}
\]
All these envelopes decrease.  In the first three allocation cases replace
$(I_1,\epsilon_2)$ by
$(\widehat I_1,\widehat\epsilon_1)$ using
Lemma~\ref{lem:bonferroni-one-coincidence}; in the fourth case retain the
two-odd-class formula.  Their respective bounds are
\[
 \begin{array}{c|cccc}
 \text{case}&E_1&E_2&E_3&\text{two odd classes}\\ \hline
 |A|/N&<0.039997&<0.037246&<0.036148&<0.034083.
 \end{array}
\]

It remains that the odd part occupies at most one odd class.  If
$e>10^{-5}N$, the even-pair formula gives
\[
 \frac{|A|}{N}<\frac{69}{100}C_{\rm quad}+\Delta_W^\dagger
 +\frac{2I_2}{25}+\frac{2\epsilon_2}{N}
 +\mathfrak T^{D,\sharp}_{4096,4096,72}<0.038318.
\]
Otherwise the mixed-strip bound leaves more than $0.00114N$ pivots in one odd
class, so two lie within $878$.  As in Theorem~\ref{thm:670b-close}, at most
one residue pair coincides.  Hence
\[
 \frac{|A|}{N}<10^{-5}+\frac{23}{100}C_{\rm quad}+\Delta_{100}^\dagger
 +\frac2{100}+\frac{6\widehat I_1}{100}
 +\frac{2+6\widehat\epsilon_1}{N}
 +\mathfrak T^{D,\sharp}_{1024,1024,116}<0.038886.
\]
All six exhaustive bounds are strictly below $1/25-7/(25N)$, so they close
directly in the original Hall balance.
\end{proof}

The dense $E_1$ branch admits a further root--small-prime correlation.  For an
$E_1$ pivot $b$, let $\rho_b$ denote the effective root count in its
transformed $q_0=25$ progression.

\begin{lemma}[At most one high-root valuation-one pivot]
\label{lem:one-high-root-e1}
Let $2.2827\cdot10^{10}\leq N\leq6.45\cdot10^{10}$ and let distinct
$b,c\in E_1$ satisfy $|b-c|\leq834$.  Then at most one of
$\rho_b,\rho_c$ exceeds $512$.

If, say, $\rho_b>512$, at least nine primes of $\mathcal P$ divide $b$.
Consequently at most four small-prime witness events remain active for that
form.  In a high--low pair the finite intersection density and endpoint error
may be bounded by
\[
 I_{HL}<0.001457,\qquad \epsilon_{HL}=50000.
\]
\end{lemma}

\begin{proof}
Because $\nu_2(b)=1$, the root formula has $c_1=1$.  A root count above $512$
therefore requires at least ten distinct odd prime factors in the effective
modulus.  At most one of them is the modulus prime $5$, so at least nine
distinct primes other than $5$ divide $b$.  If one were at least $53$, then
\[
 b\geq2\cdot53\cdot3\cdot7\cdot11\cdot13\cdot17\cdot19\cdot23\cdot29
 =68578748238>6.45\cdot10^{10}.
\]
Thus at least nine members of $\mathcal P$ divide $b$.

If both $b$ and $c$ had high root count, their two nine-element divisor sets
inside the thirteen-element set $\mathcal P$ would intersect in at least five
primes.  Hence
\[
 3\cdot7\cdot11\cdot13\cdot17=51051\mid b-c,
\]
contrary to $0<|b-c|\leq834$.

For the last assertion, enumerate the nine-element subsets
$D\subseteq\mathcal P$ satisfying $2\prod_{p\in D}p\leq6.45\cdot10^{10}$.
There are exactly eight.  Delete those inactive events from the high form and
retain all thirteen for the low form.  Exact rational enumeration over these
eight sets and the possible single coincident residue pair gives
$I_{HL}<0.001457$; the maximum has active set $\{29,37,41,47\}$ and coincident
prime $29$.  Full inclusion--exclusion has at most
\[
 (2^4-1)+(2^{13}-1)+(3^4 2^9-1)=49677<50000
\]
nonconstant CRT residue counts, proving the endpoint bound.
\end{proof}

It follows that the normalized small-prime-plus-large-tail contribution in
the dense $E_1$ branch is at most the larger of
\begin{align}
 &\frac{2\widehat I_1}{25}+\frac{2\widehat\epsilon_1}{N}
   +\mathfrak T^{D,\sharp}_{512,512,k}(N),\label{eq:e1-low-low}\\
 &\frac{2I_{HL}}{25}+\frac{100000}{N}
   +\mathfrak T^{D,\sharp}_{512,1024,k'}(N).
   \label{eq:e1-high-low}
\end{align}
Indeed the first line treats two low-root pivots, while
Lemma~\ref{lem:one-high-root-e1} reduces every remaining pair to the second
line.

\begin{theorem}[Root--small-prime correlated descent]
\label{thm:55e10-close}
The exact bound \eqref{eq:target}, and equivalently the Hall inequality \eqref{eq:hall-defect},
holds for every $N\geq5.5\cdot10^{10}$.
\end{theorem}

\begin{proof}
Theorem~\ref{thm:645b-close} leaves
\[
 5.5\cdot10^{10}\leq N<6.45\cdot10^{10}.
\]
At the left endpoint,
\[
 \begin{array}{lll}
 \Delta_{25}^\dagger<0.002661441,&
 \Delta_{50}^\dagger<0.002660928,&
 \Delta_{100}^\dagger<0.002660671,\\
 \Delta_W^\dagger<0.004418674,&
 \Delta_8^\dagger<0.004418546,&
 \mathfrak T^{D,\sharp}_{512,512,139}<0.002891955,\\
 \mathfrak T^{D,\sharp}_{512,1024,121}<0.003209719,&
 \mathfrak T^{D,\sharp}_{1024,1024,110}<0.003462020,\\
 \mathfrak T^{D,\sharp}_{2048,2048,87}<0.004176882,&
 \mathfrak T^{D,\sharp}_{4096,4096,69}<0.005074210.
 \end{array}
\]
All envelopes decrease.  If $|E_1|>0.0012N$, select the usual pair within
$834$ and use \eqref{eq:e1-low-low}--\eqref{eq:e1-high-low}; both alternatives,
including the diagonal term, give
\[
 \frac{|A|}{N}<0.039999.
\]
The next three allocation cases and the one-odd-class even-pair case give,
respectively,
\[
 \frac{|A|}{N}<0.037671,\quad0.036624,\quad0.034393,\quad0.038802.
\]

Finally suppose $e\leq10^{-5}N$ and the odd part lies in one class modulo
$4$.  The mixed-strip bound leaves more than $0.00063N$ odd pivots, so two lie
within $1588<1764$.  At most one small-prime residue pair therefore coincides,
and Lemma~\ref{lem:bonferroni-one-coincidence} gives
\[
 \frac{|A|}{N}<10^{-5}+\frac{23}{100}C_{\rm quad}+\Delta_{100}^\dagger
 +\frac2{100}+\frac{6\widehat I_1}{100}
 +\frac{2+6\widehat\epsilon_1}{N}
 +\mathfrak T^{D,\sharp}_{1024,1024,110}<0.039179.
\]
These six cases are exhaustive and each is strictly below
$1/25-7/(25N)$, giving a direct contradiction in the original Hall cut.
\end{proof}

The same correlation can be used when the last pivots are odd, but now a
single prime above $47$ is arithmetically possible.  The resulting finite
family can be enumerated exactly.

\begin{lemma}[Odd high-root finite patterns]
\label{lem:odd-high-root-patterns}
Let $N\leq4.7\cdot10^{10}$ and let $b$ be odd with effective root count above
$512$.  Then at least eight members of $\mathcal P$ divide $b$, so at most
five small-prime events remain active.  For a pair with at most two coincident
residue pairs, exact finite enumeration gives
\[
 I_{OL}<0.004200\quad\hbox{in the high--low case},\qquad
 I_{OO}<0.004054\quad\hbox{in the high--high case},
\]
with endpoint discrepancy at most $71000$ in either case.
\end{lemma}

\begin{proof}
A high root count again requires nine non-$5$ prime divisors of $b$.  Two of
them cannot both exceed $47$, since
\[
 53\cdot59\cdot3\cdot7\cdot11\cdot13\cdot17\cdot19\cdot23
 =69761140449>4.7\cdot10^{10}.
\]
Thus at least eight lie in $\mathcal P$.  Enumerate subsets
$D\subseteq\mathcal P$ of size at least nine with $\prod_{p\in D}p\leq N$,
and subsets of size eight with $53\prod_{p\in D}p\leq N$.  There are thirty
possible active complements.  Exact rational evaluation of the two-form local
products over these complements and zero, one, or two coincident pairs gives
the displayed bounds.  The high--low maximum has active set
$\{17,41,43,47\}$ with coincidences $\{17,41\}$; the high--high maximum has
the same two active sets and coincidences.  With at most five active events in
a high form, full inclusion--exclusion uses at most
\[
 (2^5-1)+(2^{13}-1)+(3^5 2^8-1)=70429<71000
\]
nonconstant CRT classes even against all thirteen events in the other form.
\end{proof}

\begin{theorem}[Two-sided root--small-prime descent]
\label{thm:47e10-close}
The exact bound \eqref{eq:target}, and equivalently the Hall inequality \eqref{eq:hall-defect},
holds for every $N\geq4.7\cdot10^{10}$.
\end{theorem}

\begin{proof}
Theorem~\ref{thm:55e10-close} leaves
\[
 4.7\cdot10^{10}\leq N<5.5\cdot10^{10}.
\]
At the left endpoint,
\[
 \begin{array}{lll}
 \Delta_{25}^\dagger<0.002801247,&
 \Delta_{50}^\dagger<0.002800681,&
 \Delta_{100}^\dagger<0.002800397,\\
 \Delta_W^\dagger<0.004655324,&
 \Delta_8^\dagger<0.004655183,&
 \mathfrak T^{D,\sharp}_{512,512,132}<0.003023853,\\
 \mathfrak T^{D,\sharp}_{512,1024,115}<0.003360097,&
 \mathfrak T^{D,\sharp}_{1024,1024,105}<0.003627133,\\
 \mathfrak T^{D,\sharp}_{2048,2048,83}<0.004383725,&
 \mathfrak T^{D,\sharp}_{4096,4096,65}<0.005333451.
 \end{array}
\]
All envelopes decrease.  In the dense $E_1$ case,
Lemma~\ref{lem:one-high-root-e1} and
\eqref{eq:e1-low-low}--\eqref{eq:e1-high-low}, now with the sharpened
$I_{HL}<0.001457$, give
\[
 \frac{|A|}{N}<0.039991.
\]
The next four allocations give
\[
 \frac{|A|}{N}<0.038115,\quad0.037120,\quad0.034718,\quad0.039308.
\]

In the final single-odd-class case, the mixed strip leaves more than
$0.00009N$ odd pivots, so two lie within $11112$.  Three coincident residue
pairs would force
\[
 \operatorname{lcm}\bigl(4,(3\cdot7\cdot11)^2\bigr)=213444
\]
to divide their difference; hence at most two coincide.  If both roots are at
most $512$, use Lemma~\ref{lem:bonferroni-two-coincidence} and the
$H=512,512$ tail.  If one or both roots are high, use
Lemma~\ref{lem:odd-high-root-patterns} with the corresponding $H=512,1024$ or
$H=1024,1024$ tail.  The low--low alternative is the largest and gives
\[
 \frac{|A|}{N}<10^{-5}+\frac{23}{100}C_{\rm quad}+\Delta_{100}^\dagger
 +\frac2{100}+\frac{6\widehat I_2}{100}
 +\frac{2+6\widehat\epsilon_2}{N}
 +\mathfrak T^{D,\sharp}_{512,512,132}<0.039912.
\]
Thus all six exhaustive allocations lie strictly below
$1/25-7/(25N)$ and close directly in the original Hall balance.
\end{proof}

The last rounding loss in the short block can itself be evaluated exactly.
Put
\[
 h_\ddagger:=\sum_{\substack{1\leq a\leq5591\\(a,43890)=1}}\frac1a
 <2.279170,
 \qquad
 K_\ddagger(R)=2\eta_*Rh_\ddagger+6.744\sqrt R.
\]
For $4.6\cdot10^{10}\leq N\leq4.7\cdot10^{10}$ define
\begin{align*}
 \Delta_q^\ddagger&=\Delta_q^{(0)}
   -\frac{3K_0(r_N+1)}N+\frac{3K_\ddagger(r_N+1)}N
   \qquad(q=25,50,100),\\
 \Delta_W^\ddagger&=\Delta_W
   -\frac{6K_0(r_N+1)}N+\frac{6K_\ddagger(r_N+1)}N,\\
 \Delta_8^\ddagger&=\Delta_8^{(0)}
   -\frac{6K_0(r_N+1)}N+\frac{6K_\ddagger(r_N+1)}N.
\end{align*}

\begin{lemma}[Exact short-block discrepancy]
\label{lem:exact-short-block-discrepancy}
On this interval the three aggregate diagonal conclusions hold with the
double-daggered errors.  These envelopes decrease.  At the left endpoint,
\[
 \begin{array}{lll}
 \Delta_{25}^\ddagger<0.002758521,&
 \Delta_{50}^\ddagger<0.002757946,&
 \Delta_{100}^\ddagger<0.002757658,\\
 \Delta_W^\ddagger<0.004563837,&
 \Delta_8^\ddagger<0.004563695.
 \end{array}
\]
\end{lemma}

\begin{proof}
Write $C(y)=\#\{a\leq y:(a,43890)=1\}$.  The exact finite recurrence
\[
 C(y)=C(y-1)+\boldsymbol 1_{(y,43890)=1}
\]
over $0\leq y\leq5591$ gives
\[
 \max_{0\leq y\leq5591}
 \left(C(y)-\frac{288}{1463}y\right)
 =828-\frac{288}{1463}\,4189
 =\frac{4932}{1463}<3.372.
\]
The harmonic inequality defining $h_\ddagger$ is likewise an exact rational
summation.  Since $\sqrt{r_N+1}<5591$ throughout the interval, the hyperbola
step in Lemma~\ref{lem:short-block-diagonal} now gives
\[
 \sum_{\substack{s\leq R\\(s,43890)=1}}d(s)
 \leq2\eta_*Rh_\ddagger+2(3.372)\sqrt R=K_\ddagger(R).
\]
Substitution in the three diagonal arguments proves the formulas.
Elementary differentiation gives monotonicity, and direct left-endpoint
evaluation gives the displayed strict bounds.
\end{proof}

\begin{theorem}[Exact-discrepancy descent]
\label{thm:46e10-close}
The exact bound \eqref{eq:target}, and equivalently the Hall inequality \eqref{eq:hall-defect},
holds for every $N\geq4.6\cdot10^{10}$.
\end{theorem}

\begin{proof}
Theorem~\ref{thm:47e10-close} leaves
\[
 4.6\cdot10^{10}\leq N<4.7\cdot10^{10}.
\]
At the left endpoint, optimizing the same integer split parameter gives
\[
 \begin{array}{lll}
 \mathfrak T^{D,\sharp}_{512,512,131}<0.003042501,&
 \mathfrak T^{D,\sharp}_{512,1024,115}<0.003381367,&
 \mathfrak T^{D,\sharp}_{1024,1024,104}<0.003650423,\\
 \mathfrak T^{D,\sharp}_{2048,2048,82}<0.004412878,&
 \mathfrak T^{D,\sharp}_{4096,4096,65}<0.005370064.
 \end{array}
\]
All terms decrease.  Use Lemma~\ref{lem:exact-short-block-discrepancy} in the
same six exhaustive allocations as Theorem~\ref{thm:47e10-close}.  The first
five endpoint bounds are, respectively,
\[
 \frac{|A|}{N}<0.039967,\quad0.0382,\quad0.0372,\quad0.0349,\quad0.0394.
\]

For completeness, the matching strip now satisfies
$\mathcal R(N)>0.0000268N$.  In the final branch $e\leq10^{-5}N$, it leaves
more than $0.0000168N$ odd pivots in one class modulo $4$, so two lie within
$59524$.  Three coincident residue pairs would again force $213444$ to divide
their nonzero difference.  Hence at most two coincide, and the low--low
alternative remains the largest.  Lemma~\ref{lem:bonferroni-two-coincidence}
and the $H=512,512$ tail give
\[
 \frac{|A|}{N}<10^{-5}+\frac{23}{100}C_{\rm quad}
 +\Delta_{100}^\ddagger+\frac2{100}+\frac{6\widehat I_2}{100}
 +\frac{2+6\widehat\epsilon_2}{N}
 +\mathfrak T^{D,\sharp}_{512,512,131}<0.039888.
\]
Every bound is strictly below $1/25-7/(25N)$, giving a direct contradiction
inside the original Hall cut.
\end{proof}

One more short block is available on the interval just below this endpoint.
Put
\[
 h_\natural:=\sum_{\substack{1\leq a\leq5551\\(a,43890)=1}}\frac1a
 <2.277735,
 \qquad
 K_\natural(R)=2\eta_*Rh_\natural+6.744\sqrt R,
\]
and for $4.5715\cdot10^{10}\leq N\leq4.6\cdot10^{10}$ define
\begin{align*}
 \Delta_q^\natural&=\Delta_q^{(0)}
  -\frac{3K_0(r_N+1)}N+\frac{3K_\natural(r_N+1)}N
  \qquad(q=25,50,100),\\
 \Delta_W^\natural&=\Delta_W
  -\frac{6K_0(r_N+1)}N+\frac{6K_\natural(r_N+1)}N,\\
 \Delta_8^\natural&=\Delta_8^{(0)}
  -\frac{6K_0(r_N+1)}N+\frac{6K_\natural(r_N+1)}N.
\end{align*}
The proof of Lemma~\ref{lem:exact-short-block-discrepancy} applies unchanged:
$\sqrt{r_N+1}<5551$, while the exact maximum coprime-count discrepancy is
still attained at $4189$.  The naturalled envelopes decrease, and at the left
endpoint
\[
 \begin{array}{lll}
 \Delta_{25}^\natural<0.002762881,&
 \Delta_{50}^\natural<0.002762304,&
 \Delta_{100}^\natural<0.002762016,\\
 \Delta_W^\natural<0.004570809,&
 \Delta_8^\natural<0.004570666.
 \end{array}
\]

\begin{theorem}[Vanishing-strip descent]
\label{thm:45715e10-close}
The exact bound \eqref{eq:target}, and equivalently the Hall inequality \eqref{eq:hall-defect},
holds for every $N\geq4.5715\cdot10^{10}$.
\end{theorem}

\begin{proof}
Theorem~\ref{thm:46e10-close} leaves
\[
 4.5715\cdot10^{10}\leq N<4.6\cdot10^{10}.
\]
At the left endpoint the optimized tails are
\[
 \begin{array}{lll}
 \mathfrak T^{D,\sharp}_{512,512,131}<0.003047890,&
 \mathfrak T^{D,\sharp}_{512,1024,114}<0.003387517,&
 \mathfrak T^{D,\sharp}_{1024,1024,104}<0.003657224,\\
 \mathfrak T^{D,\sharp}_{2048,2048,82}<0.004421371,&
 \mathfrak T^{D,\sharp}_{4096,4096,65}<0.005380780.
 \end{array}
\]
The first four allocations from Theorem~\ref{thm:47e10-close}, with the
naturalled diagonal errors, are respectively below
\[
 0.039977,\qquad0.0382,\qquad0.0372,\qquad0.0349.
\]

It remains to split according to the total even mass $e$.  If
$e>0.000000056N$, two even pivots lie within $17857143$.  Four coincident
small-prime pairs would force
\[
 2(3\cdot7\cdot11\cdot13)^2=18036018
\]
to divide their nonzero difference.  Thus at most three pairs coincide.  The
full finite product $I_3<0.136690$, its endpoint budget $\epsilon_2=1610704$,
and the $H=4096,4096$ tail give
\[
 \frac{|A|}{N}<\frac{69}{100}C_{\rm quad}+\Delta_W^\natural
 +\frac{2I_3}{25}+\frac{2\epsilon_2}{N}
 +\mathfrak T^{D,\sharp}_{4096,4096,65}<0.039827.
\]

Otherwise the mixed strip $\mathcal R(N)>0.0000047808N$ leaves more than
$0.0000047248N$ odd pivots in one class modulo $4$.  Two lie within $211650$,
which is below
\[
 4(3\cdot7\cdot11)^2=213444.
\]
Hence at most two pairs coincide.  The low--low alternative is again largest,
and gives
\[
 \frac{|A|}{N}<0.000000056+\frac{23}{100}C_{\rm quad}
 +\Delta_{100}^\natural+\frac2{100}+\frac{6\widehat I_2}{100}
 +\frac{2+6\widehat\epsilon_2}{N}
 +\mathfrak T^{D,\sharp}_{512,512,131}<0.039888.
\]
All six exhaustive bounds are strictly below $1/25-7/(25N)$ and therefore
close directly in the original Hall balance.
\end{proof}


We now take the first of those two terminal alternatives: enlarge the finite
prime window while subtracting the added primes from the same explicit tail.
Let
\[
 \mathcal P_{113}=\{p:3\leq p\leq113,\ p\ne5\},
 \qquad |\mathcal P_{113}|=28.
\]
Define the degree-$5$ individual and degree-$4$ joint Bonferroni polynomials
exactly as before, but over $\mathcal P_{113}$.

\begin{lemma}[Extended finite Bonferroni window]
\label{lem:extended-bonferroni-window}
For at most $r$ coincident residue pairs, $r=1,3,4$, the corresponding
two-form intersection densities satisfy
\[
 \widehat I^{113}_1<0.112772853,
 \qquad \widehat I^{113}_3<0.136772805,
 \qquad \widehat I^{113}_4<0.141729589.
\]
A single endpoint discrepancy valid uniformly for all three assertions is
\[
 \widehat\epsilon_{113}=600250.
\]
Moreover
\[
 \sum_{p>113}\frac1{p^2}<0.001386157=: \tau_{113}.
\]
\end{lemma}

\begin{proof}
Exact rational evaluation of the Bonferroni polynomials gives the displayed
densities; for each $r$ the maximum occurs when the coincident primes are the
$r$ smallest members of $\mathcal P_{113}$.  Each individual degree-$5$
polynomial has
\[
 \sum_{j=1}^5\binom{28}{j}=122437
\]
nonconstant counts.  The zero-coincidence joint polynomial, which maximizes
the endpoint count, has degree counts
\[
 56,\quad1512,\quad26208,\quad327600.
\]
Thus the uniform total is
\[
 2(122437)+56+1512+26208+327600=600250.
\]
Finally subtract the exact rational sum
$\sum_{47<p\leq113}p^{-2}$ from the proved bound
$\sum_{p>47}p^{-2}<0.003887$ to obtain $\tau_{113}$.
\end{proof}

Define $\mathfrak T^{D,113}_{H_1,H_2,k}$ by replacing $\tau_\sharp$ with
$\tau_{113}$ in $\mathfrak T^{D,\sharp}_{H_1,H_2,k}$.  For the last short
block put
\[
 h_\flat:=\sum_{\substack{1\leq a\leq5540\\(a,43890)=1}}\frac1a
 <2.277374,
 \qquad
 K_\flat(R)=2\eta_*Rh_\flat+6.744\sqrt R,
\]
and define $\Delta_q^\flat,\Delta_W^\flat,\Delta_8^\flat$ by replacing
$K_0$ with $K_\flat$ in the same manner as above.  On
$4.5654\cdot10^{10}\leq N\leq4.5715\cdot10^{10}$ the exact short-block proof
applies, and at the left endpoint
\[
 \begin{array}{lll}
 \Delta_{25}^\flat<0.002763777,&
 \Delta_{50}^\flat<0.002763200,&
 \Delta_{100}^\flat<0.002762911,\\
 \Delta_W^\flat<0.004572225,&
 \Delta_8^\flat<0.004572081.
 \end{array}
\]

\begin{theorem}[Extended-window descent]
\label{thm:45654e10-close}
The exact bound \eqref{eq:target}, and equivalently the Hall inequality \eqref{eq:hall-defect},
holds for every $N\geq4.5654\cdot10^{10}$.
\end{theorem}

\begin{proof}
Theorem~\ref{thm:45715e10-close} leaves
\[
 4.5654\cdot10^{10}\leq N<4.5715\cdot10^{10}.
\]
At the left endpoint,
\[
 \begin{array}{lll}
 \mathfrak T^{D,113}_{512,512,131}<0.002648917,&
 \mathfrak T^{D,113}_{512,1024,114}<0.002988705,&
 \mathfrak T^{D,113}_{1024,1024,104}<0.003258556,\\
 \mathfrak T^{D,113}_{2048,2048,82}<0.004023067,&
 \mathfrak T^{D,113}_{4096,4096,65}<0.004982955.
 \end{array}
\]
Use Lemma~\ref{lem:extended-bonferroni-window} in the low--low alternative.
The high-root alternatives remain smaller.  The first four exhaustive
allocations are below
\[
 0.039620,\qquad0.0382,\qquad0.0372,\qquad0.0349.
\]

Split at $e=0.0000000002N$.  Above that cutoff, two even pivots lie within
$5000000000$.  Five coincidences would force
\[
 2(3\cdot7\cdot11\cdot13\cdot17)^2=5212409202
\]
to divide their difference, so at most four coincide.  Hence
\[
 \frac{|A|}{N}<\frac{69}{100}C_{\rm quad}+\Delta_W^\flat
 +\frac{2\widehat I^{113}_4}{25}
 +\frac{2\widehat\epsilon_{113}}N
 +\mathfrak T^{D,113}_{4096,4096,65}<0.039790.
\]

Below the cutoff, the directly evaluated strip
$\mathcal R(N)>0.00000003669N$ leaves more than $0.00000003649N$ odd pivots
in one class modulo $4$.  Two lie within $27405000$, below
\[
 4(3\cdot7\cdot11\cdot13)^2=36072036.
\]
Thus at most three pairs coincide, and
\[
 \frac{|A|}{N}<0.0000000002+\frac{23}{100}C_{\rm quad}
 +\Delta_{100}^\flat+\frac2{100}
 +\frac{6\widehat I^{113}_3}{100}
 +\frac{2+6\widehat\epsilon_{113}}N
 +\mathfrak T^{D,113}_{512,512,131}<0.039987.
\]
Every one of the six bounds is strictly below $1/25-7/(25N)$, closing the
original Hall balance.
\end{proof}


For the longer descent use
\[
 \mathcal P_{131}=\{p:3\leq p\leq131,\ p\ne5\},
 \qquad |\mathcal P_{131}|=30.
\]

\begin{lemma}[Thirty-prime Bonferroni window]
\label{lem:thirty-prime-window}
The degree-$5$/degree-$4$ truncation over $\mathcal P_{131}$ gives
\[
 \widehat I^{131}_1<0.112782897,
 \qquad
 \widehat I^{131}_2<0.129880546,
 \qquad
 \widehat\epsilon_{131}=821632,
\]
and the remaining prime-square tail satisfies
\[
 \sum_{p>131}\frac1{p^2}<0.001265886=: \tau_{131}.
\]
\end{lemma}

\begin{proof}
Exact rational evaluation is maximal at coincidence sets $\{3\}$ and
$\{3,7\}$, respectively.  Each individual polynomial now has
$\sum_{j=1}^5\binom{30}{j}=174436$ nonconstant counts.  The zero-coincidence
joint degree counts are
\[
 60,\quad1740,\quad32480,\quad438480,
\]
so the uniform endpoint total is
\[
 2(174436)+60+1740+32480+438480=821632.
\]
Exact subtraction of $\sum_{47<p\leq131}p^{-2}$ from $0.003887$ proves the
tail bound.
\end{proof}

Define $\mathfrak T^{D,131}$ by using $\tau_{131}$ in the same tail formula.
The flat short-block proof remains valid on
$3.71\cdot10^{10}\leq N\leq4.5654\cdot10^{10}$.  At the new left endpoint,
\[
 \begin{array}{lll}
 \Delta_{25}^\flat<0.002953746,&
 \Delta_{50}^\flat<0.002953088,&
 \Delta_{100}^\flat<0.002952758,\\
 \Delta_W^\flat<0.004891857,&
 \Delta_8^\flat<0.004891694,
 \end{array}
\]
and
\[
 \begin{array}{lll}
 \mathfrak T^{D,131}_{512,512,123}<0.002817694,&
 \mathfrak T^{D,131}_{512,1024,107}<0.003183824,&
 \mathfrak T^{D,131}_{1024,1024,97}<0.003474609,\\
 \mathfrak T^{D,131}_{2048,2048,77}<0.004298780,&
 \mathfrak T^{D,131}_{4096,4096,60}<0.005333610.
 \end{array}
\]

\begin{theorem}[Half-root descent]
\label{thm:371e10-close}
The exact bound \eqref{eq:target}, and equivalently the Hall inequality \eqref{eq:hall-defect},
holds for every $N\geq3.71\cdot10^{10}$.
\end{theorem}

\begin{proof}
Theorem~\ref{thm:45654e10-close} leaves
\[
 3.71\cdot10^{10}\leq N<4.5654\cdot10^{10}.
\]
Use Proposition~\ref{prop:half-root-correlated-degree}, the flat diagonal
errors, and Lemma~\ref{lem:thirty-prime-window}.  The first four allocation
bounds are
\[
 \frac{|A|}{N}<0.039998,\qquad0.038327,\qquad0.037417,\qquad0.034754.
\]

If $e>10^{-5}N$, two even pivots lie within $100000$.  Three coincidences
would force
\[
 2(3\cdot7\cdot11)^2=106722
\]
to divide their difference, so at most two coincide.  Therefore
\[
 \frac{|A|}{N}<\frac{69}{100}C_{\rm quad}+\Delta_W^\flat
 +\frac{2\widehat I^{131}_2}{25}
 +\frac{2\widehat\epsilon_{131}}N
 +\mathfrak T^{D,131}_{4096,4096,60}<0.039530.
\]

Otherwise $\mathcal R(N)>0.001593N$ leaves more than $0.001583N$ odd pivots
in one class modulo $4$, so two lie within $632$.  Two coincidences would
force $4(3\cdot7)^2=1764$ to divide their difference.  Thus at most one
coincides, and
\[
 \frac{|A|}{N}<10^{-5}+\frac{23}{100}C_{\rm quad}
 +\Delta_{100}^\flat+\frac2{100}
 +\frac{6\widehat I^{131}_1}{100}
 +\frac{2+6\widehat\epsilon_{131}}N
 +\mathfrak T^{D,131}_{512,512,123}<0.038971.
\]
All six bounds decrease and lie strictly below $1/25-7/(25N)$, closing
directly in the original Hall balance.
\end{proof}


We next sharpen the first two of those terms.

\begin{lemma}[Three-term prime count and integrated square tail]
\label{lem:three-term-prime-tail}
For $y\geq355991$, Dusart's explicit estimate gives
\[
 \pi(y)\leq\frac{y}{\log y}
 \left(1+\frac1{\log y}+\frac{2.51}{\log^2y}\right).
\]
Moreover
\[
 \sum_{p>131}\frac1{p^2}<0.001257623=: \tau_{131}^+.
\]
Consequently, if $L=\log(N/(k+1))$, the intermediate-prime payment in the
two-pivot tail is at most
\[
 \frac4k\left(\frac1L+\frac1{L^2}+\frac{2.51}{L^3}\right).
\]
\end{lemma}

\begin{proof}
The displayed prime-count estimate is Theorem~6.9 of
\cite{DusartPrimes}.  Exact rational summation gives
\[
 \sum_{131<p\leq10^5}\frac1{p^2}<0.0012554421.
\]
Partial summation and the classical bound
$\pi(t)<1.25506t/\log t$ give
\[
 \sum_{p>10^5}\frac1{p^2}
 \leq2(1.25506)\int_{10^5}^{\infty}\frac{dt}{t^2\log t}
 <\frac{2(1.25506)}{10^5\log(10^5)}<0.0000021803.
\]
This proves the square-tail constant.  Finally
$N/(k+1)\leq\lfloor N/k\rfloor\leq N/k$; monotonicity of the displayed
prime-count majorant gives the normalized payment.
\end{proof}

Define $\mathfrak T^{+,131}_{H_1,H_2,k}$ by
\begin{align*}
 \mathfrak T^{+,131}_{H_1,H_2,k}(N)={}&\frac{4\tau_{131}^+}{25}
 +\frac4k\left(\frac1L+\frac1{L^2}+\frac{2.51}{L^3}\right)\\
 &+\frac{2(H_1+H_2)}N\left(
 \frac{2(N/25+2)}{N/(k+1)}
 +\frac{N^2+25N+1}{(N/(k+1))^2}\right),
\end{align*}
where $L=\log(N/(k+1))$.

The remaining diagonal divisor sum can be evaluated rather than enveloped.
Put
\[
 \mathcal S(R)=\sum_{\substack{s\leq R\\(s,43890)=1}}d(s),
 \qquad
 C(y)=\#\{a\leq y:(a,43890)=1\}.
\]

\begin{lemma}[Exact hyperbola certificate]
\label{lem:exact-hyperbola-certificate}
For $3.51\cdot10^{10}\leq N\leq3.71\cdot10^{10}$,
\[
 \frac{\mathcal S(r_N+1)}N<0.000623885.
\]
Hence the three aggregate diagonal estimates hold with
\begin{align*}
 \Delta_q^{\rm ex}(N)&=\frac{23}{N}+3(0.000623885)
 +\frac{46(\log(r_N+1)+2)}{q r_N}
 +\frac{N^2+1}{Nr_N^2}(1+\log_{2+\sqrt3}N),\\
 \Delta_W^{\rm ex}(N)&=\frac{69}{N}+6(0.000623885)
 +\frac{69(\log(r_N+1)+2)}{50r_N}
 +\frac{N^2+1}{Nr_N^2}(1+\log_{2+\sqrt3}N),\\
 \Delta_8^{\rm ex}(N)&=\frac{115}{N}+6(0.000623885)
 +\frac{23(\log(r_N+1)+2)}{20r_N}
 +\frac{N^2+1}{Nr_N^2}(1+\log_{2+\sqrt3}N).
\end{align*}
At the left endpoint these are below, respectively,
\[
 0.002903682,\quad0.002903000,\quad0.002902659,
 \quad0.004774997,\quad0.004774828.
\]
\end{lemma}

\begin{proof}
The exact hyperbola identity is
\[
 \mathcal S(R)=2\!\sum_{\substack{a\leq\sqrt R\\(a,43890)=1}}
 C(\lfloor R/a\rfloor)-C(\lfloor\sqrt R\rfloor)^2,
\]
where
\[
 C(y)=8640\left\lfloor\frac y{43890}\right\rfloor
 +C(y\bmod43890).
\]
On the stated interval, $\lfloor r_N+1\rfloor$ runs from $25728590$ to
$26696881$.  Divide this range into consecutive blocks of length $1000$.
On each block use the exact value of $\mathcal S$ at the upper endpoint and
the least integer $N$ mapping into the block.  The maximum occurs on the first
block and equals
\[
 \frac{21898353}{35100000000}<0.000623885.
\]
This finite calculation uses only the displayed identity and the $43890$
residue prefix.  Substitution proves the diagonal formulas; all remaining
terms decrease, and direct endpoint evaluation gives the five numbers.
\end{proof}

\begin{theorem}[Exact-diagonal descent]
\label{thm:351e10-close}
The exact bound \eqref{eq:target}, and equivalently the Hall inequality \eqref{eq:hall-defect},
holds for every $N\geq3.51\cdot10^{10}$.
\end{theorem}

\begin{proof}
Theorem~\ref{thm:371e10-close} leaves
\[
 3.51\cdot10^{10}\leq N<3.71\cdot10^{10}.
\]
At the left endpoint the optimized new tails are
\[
 \begin{array}{lll}
 \mathfrak T^{+,131}_{512,512,120}<0.002865809,&
 \mathfrak T^{+,131}_{512,1024,105}<0.003238848,&
 \mathfrak T^{+,131}_{1024,1024,95}<0.003535139,\\
 \mathfrak T^{+,131}_{2048,2048,75}<0.004374804,&
 \mathfrak T^{+,131}_{4096,4096,59}<0.005429329.
 \end{array}
\]
Use Lemmas~\ref{lem:thirty-prime-window},
\ref{lem:three-term-prime-tail}, and
\ref{lem:exact-hyperbola-certificate}.  The first four allocations give
\[
 \frac{|A|}{N}<0.039998,\qquad0.038289,\qquad0.037398,\qquad0.034775.
\]
If $e>10^{-5}N$, the even pair again has at most two coincidences and gives
\[
 \frac{|A|}{N}<\frac{69}{100}C_{\rm quad}+\Delta_W^{\rm ex}
 +\frac{2\widehat I^{131}_2}{25}
 +\frac{2\widehat\epsilon_{131}}N
 +\mathfrak T^{+,131}_{4096,4096,59}<0.039511.
\]
Otherwise $\mathcal R(N)>0.001424N$ leaves more than $0.001414N$ odd pivots
in one class modulo $4$, so two lie within $707<1764$ and at most one pair
coincides.  Thus
\[
 \frac{|A|}{N}<10^{-5}+\frac{23}{100}C_{\rm quad}
 +\Delta_{100}^{\rm ex}+\frac2{100}
 +\frac{6\widehat I^{131}_1}{100}
 +\frac{2+6\widehat\epsilon_{131}}N
 +\mathfrak T^{+,131}_{512,512,120}<0.038976.
\]
All six bounds are strictly below $1/25-7/(25N)$ and close the original Hall
balance directly.
\end{proof}


Let
\[
 \mathcal P_{137}=\{p:3\leq p\leq137,\ p\ne5\},
 \qquad |\mathcal P_{137}|=31.
\]
Define the individual degree-$5$ lower polynomial and the joint degree-$4$
upper polynomial as in Lemma~\ref{lem:bonferroni-one-coincidence}.

\begin{lemma}[Conditional endpoint certificate at $137$]
\label{lem:conditional-137-window}
If no residue pair coincides, the two-form intersection density is below
\[
 \widehat I^{137}_0<0.012349183
\]
and its endpoint discrepancy is at most $954056$.  If exactly one pair
coincides, the corresponding bounds are
\[
 \widehat I^{137}_1<0.112787353,
 \qquad \widehat\epsilon^{137}_1=919775.
\]
Consequently, whenever at most one pair coincides and
$N\geq3.4984\cdot10^{10}$, the contribution from the two base progressions is
at most
\[
 \frac{2\widehat I^{137}_1}{25}N+2\widehat\epsilon^{137}_1.
\]
Moreover
\[
 \sum_{p>137}\frac1{p^2}<0.001204343=: \tau_{137}^+.
\]
\end{lemma}

\begin{proof}
Exact rational evaluation of $1-2L_5+U_4$ gives the first two density
bounds; in the one-coincidence case the maximum is attained at the prime
$3$.  Each individual polynomial has
\[
 \sum_{j=1}^5\binom{31}{j}=206367
\]
nonconstant counts.  The joint degree counts are
\[
 62,\ 1860,\ 35960,\ 503440
\]
with no coincidence and
\[
 61,\ 1800,\ 34220,\ 470960
\]
with exactly one.  Hence the two endpoint totals are respectively
$954056$ and $919775$.  At the least allowed $N$, the two normalized
alternatives are below
\[
 \frac{2\widehat I^{137}_0}{25}+\frac{2(954056)}N
 <0.001042478,
 \qquad
 \frac{2\widehat I^{137}_1}{25}+\frac{2(919775)}N
 <0.009075571.
\]
Both endpoint terms decrease, so the second alternative dominates on the
whole range.  This proves the combined assertion while retaining the larger
endpoint count in the zero-coincidence calculation.

Finally, exact rational summation gives
\[
 \sum_{137<p\leq10^5}\frac1{p^2}<0.0012021627.
\]
The integrated tail in Lemma~\ref{lem:three-term-prime-tail} is below
$0.0000021803$; their strict sum is below $0.001204343$.
\end{proof}

\begin{lemma}[Exact bottom short-interval blocks]
\label{lem:35e10-exact-blocks}
For
\[
 3.4984\cdot10^{10}\leq N\leq3.51\cdot10^{10}
\]
one has
\[
 \frac{\mathcal S(\lfloor r_N+1\rfloor)}N<0.000624485942,
 \qquad
 \frac{4\pi(\lfloor N/120\rfloor)}N<0.001808618752.
\]
Thus the exact diagonal estimates hold with
\begin{align*}
 \Delta_q^{\rm ex,0}(N)&=\frac{23}{N}+3(0.000624485942)
 +\frac{46(\log(r_N+1)+2)}{q r_N}
 +\frac{N^2+1}{Nr_N^2}(1+\log_{2+\sqrt3}N),\\
 \Delta_W^{\rm ex,0}(N)&=\frac{69}{N}+6(0.000624485942)
 +\frac{69(\log(r_N+1)+2)}{50r_N}
 +\frac{N^2+1}{Nr_N^2}(1+\log_{2+\sqrt3}N),\\
 \Delta_8^{\rm ex,0}(N)&=\frac{115}{N}+6(0.000624485942)
 +\frac{23(\log(r_N+1)+2)}{20r_N}
 +\frac{N^2+1}{Nr_N^2}(1+\log_{2+\sqrt3}N).
\end{align*}
At the left endpoint the five relevant values are below
\[
 0.002906492,\quad0.002905809,\quad0.002905468,
 \quad0.004779610,\quad0.004779440.
\]
\end{lemma}

\begin{proof}
Here $\lfloor r_N+1\rfloor$ runs from $25671873$ to $25728590$.
Partition that integer interval into blocks of length $10$ and use the
exact hyperbola identity of Lemma~\ref{lem:exact-hyperbola-certificate} at
each upper endpoint.  The largest normalized value is
\[
 \frac{\mathcal S(25671892)}{34984018923}
 =\frac{21847028}{34984018923}<0.000624485942.
\]

For the prime block, $Y=\lfloor N/120\rfloor$ runs from $291533333$ to
$292500000$.  An Eratosthenes sieve gives
\[
 \pi(291533332)=15818171,
 \qquad
 \pi(292500000)-\pi(291533332)=49764.
\]
Partition the $Y$-interval into blocks of length $10$, using the prime
count at the upper endpoint and the least $N$ mapping into the block.  The
maximum is
\[
 \frac{4\pi(291534732)}{34984166760}
 =\frac{4(15818255)}{34984166760}<0.001808618752.
\]
Both scans are finite integer certificates.  Substitution in the aggregate
diagonal formulas gives the remaining assertions.
\end{proof}

On this short interval put
\begin{align*}
 \mathfrak T^{\rm pc,137}_{512,512,120}(N)={}&
 \frac{4\tau_{137}^+}{25}+0.001808618752\\
 &+\frac{2048}{N}\left(
 \frac{2(N/25+2)}{N/121}
 +\frac{N^2+25N+1}{(N/121)^2}\right).
\end{align*}
Lemma~\ref{lem:35e10-exact-blocks} and the original two-pivot tail proof
show that this is a valid upper bound; at the left endpoint it is below
$0.002858980$.

\begin{theorem}[Conditional-endpoint descent]
\label{thm:35e10-close}
The exact bound \eqref{eq:target}, and equivalently the Hall inequality \eqref{eq:hall-defect},
holds for every $N\geq3.4984\cdot10^{10}$.
\end{theorem}

\begin{proof}
Theorem~\ref{thm:351e10-close} leaves
\[
 3.4984\cdot10^{10}\leq N<3.51\cdot10^{10}.
\]
At the left endpoint Proposition~\ref{prop:half-root-correlated-degree}
gives a matching margin greater than $2.4\cdot10^7$, and
Corollary~\ref{cor:thin-nonbase-strip} gives
$\mathcal R(N)>0.001414N$.  The first four exhaustive allocations, using
Lemma~\ref{lem:35e10-exact-blocks}, are below
\[
 0.039999998,\qquad0.038299,\qquad0.037409,\qquad0.034782.
\]
For the dense $E_1$ branch, two pivots lie within $834$, hence at most one
small-prime residue pair coincides.  The low--low alternative is
\begin{align*}
 \frac{|A|}{N}<{}&\frac{23}{25}C_{\rm quad}
 +\Delta_{25}^{\rm ex,0}
 +\frac{2\widehat I^{137}_1}{25}
 +\frac{2\widehat\epsilon^{137}_1}{N}
 +\mathfrak T^{\rm pc,137}_{512,512,120}\\
 <{}&0.039999998.
\end{align*}
The high--low alternative from Lemma~\ref{lem:one-high-root-e1}, with its
original cutoff-$47$ tail with $k=105$, is below $0.031852$ and is therefore
smaller.

If $e>10^{-5}N$, the even pair has at most two coincidences.  Keeping the
cutoff-$131$ bounds gives
\[
 \frac{|A|}{N}<\frac{69}{100}C_{\rm quad}+\Delta_W^{\rm ex,0}
 +\frac{2\widehat I^{131}_2}{25}
 +\frac{2\widehat\epsilon_{131}}N
 +\mathfrak T^{+,131}_{4096,4096,59}<0.039522.
\]
Otherwise more than $0.001404N$ odd pivots remain in one class modulo $4$.
Two lie within $713<1764$, so at most one pair coincides, and
\[
 \frac{|A|}{N}<10^{-5}+\frac{23}{100}C_{\rm quad}
 +\Delta_{100}^{\rm ex,0}+\frac2{100}
 +\frac{6\widehat I^{131}_1}{100}
 +\frac{2+6\widehat\epsilon_{131}}N
 +\mathfrak T^{+,131}_{512,512,120}<0.038983.
\]
Every term not covered by a uniform finite block decreases on this interval.
All six bounds are therefore strictly below
$1/25-7/(25N)$ and directly close the original Hall balance.
\end{proof}


Let
\[
 \mathcal P_{509}=\{p:3\leq p\leq509,\ p\ne5\},
 \qquad |\mathcal P_{509}|=95.
\]

\begin{lemma}[Conditional degree-$3/2$ certificate at $509$]
\label{lem:conditional-509-window}
Use the degree-$3$ individual lower polynomial and degree-$2$ joint upper
polynomial on $\mathcal P_{509}$.  With no coincidence the intersection
density and endpoint discrepancy are bounded by
\[
 \widehat J^{509}_0<0.013594002,\qquad \zeta^{509}_0=304000.
\]
With exactly one coincidence they are bounded by
\[
 \widehat J^{509}_1<0.113400471,\qquad \zeta^{509}_1=303811.
\]
For $N\geq3.2532\cdot10^{10}$ and at most one coincidence, the contribution
from both base progressions is therefore at most
\[
 \frac{2\widehat J^{509}_1}{25}N+2\zeta^{509}_1.
\]
The remaining square tail satisfies
\[
 \sum_{p>509}\frac1{p^2}<0.000266795=:\tau_{509}^+.
\]
\end{lemma}

\begin{proof}
Exact rational evaluation of $1-2L_3+U_2$ gives the displayed densities;
the exactly-one maximum again occurs at $3$.  Each individual polynomial has
\[
 \sum_{j=1}^3\binom{95}{j}=142975
\]
nonconstant counts.  The zero-coincidence joint degree counts are
$190,17860$, while the exactly-one counts are $189,17672$.  Hence
\[
 2(142975)+190+17860=304000,
 \qquad
 2(142975)+189+17672=303811.
\]
At the least $N$, the two normalized alternatives are below
$0.001106210$ and $0.009090716$, respectively; the latter dominates
thereafter.  Finally exact rational summation gives
\[
 \sum_{509<p\leq10^5}\frac1{p^2}<0.0002646140.
\]
Adding the strict integrated tail $0.0000021803$ proves the last assertion.
\end{proof}

\begin{lemma}[Exact $509$-window interval blocks]
\label{lem:509-window-blocks}
For
\[
 3.2532\cdot10^{10}\leq N\leq3.4984\cdot10^{10}
\]
one has
\[
 \frac{\mathcal S(\lfloor r_N+1\rfloor)}N<0.000638391308,
 \qquad
 \frac{4\pi(\lfloor N/117\rfloor)}N<0.001859823072.
\]
Write $\Delta_q^{\rm ex,509},\Delta_W^{\rm ex,509},
\Delta_8^{\rm ex,509}$ for the exact diagonal formulas with
$0.000638392$ in place of the normalized divisor sum.  Their five
left-endpoint errors are below
\[
 0.002970568,\quad0.002969853,\quad0.002969495,
 \quad0.004885388,\quad0.004885211.
\]
\end{lemma}

\begin{proof}
The integer $\lfloor r_N+1\rfloor$ runs from $24457858$ to $25671873$.
Using length-$100$ blocks and the exact hyperbola identity, the maximum is
\[
 \frac{\mathcal S(24457957)}{32532000000}
 =\frac{20768146}{32532000000}<0.000638391308.
\]
For $Y=\lfloor N/117\rfloor$, the range is $278051282\leq Y\leq299008547$.
The exact sieve gives
\[
 \pi(278051281)=15125934,\qquad
 \pi(299008547)-\pi(278051281)=1075477.
\]
The same length-$100$ scan is maximal at the block ending at $278052281$:
\[
 \frac{4\pi(278052281)}{32532105294}
 =\frac{4(15125990)}{32532105294}<0.001859823072.
\]
Substitution proves the five diagonal bounds.
\end{proof}

On this interval define
\begin{align*}
 \mathfrak T^{\rm pc,509}_{512,512,117}(N)={}&
 \frac{4\tau_{509}^+}{25}+0.001859823072\\
 &+\frac{2048}{N}\left(
 \frac{2(N/25+2)}{N/118}
 +\frac{N^2+25N+1}{(N/118)^2}\right).
\end{align*}
This is below $0.002779668$ at the left endpoint and decreases apart from
the already uniform prime block.

\begin{theorem}[Order-optimized conditional descent]
\label{thm:3255e10-close}
The exact bound \eqref{eq:target}, and equivalently the Hall inequality \eqref{eq:hall-defect},
holds for every $N\geq3.2532\cdot10^{10}$.
\end{theorem}

\begin{proof}
Theorem~\ref{thm:35e10-close} leaves
\[
 3.2532\cdot10^{10}\leq N<3.4984\cdot10^{10}.
\]
At the left endpoint the matching margin exceeds $1.9\cdot10^7$ and
$\mathcal R(N)>0.001187N$.  Lemmas~\ref{lem:conditional-509-window} and
\ref{lem:509-window-blocks} give the first four allocation bounds
\[
 0.039999904,\qquad0.038521,\qquad0.037660,\qquad0.034950.
\]
In the dense $E_1$ branch, the low--low alternative is
\[
 \frac{|A|}{N}<\frac{23}{25}C_{\rm quad}+\Delta_{25}^{\rm ex,509}
 +\frac{2\widehat J^{509}_1}{25}
 +\frac{2\zeta^{509}_1}{N}
 +\mathfrak T^{\rm pc,509}_{512,512,117}<0.039999904,
\]
where $\Delta^{\rm ex,509}$ denotes the formulas of
Lemma~\ref{lem:509-window-blocks}.  The original high--low alternative,
optimized at $k=102$, is below $0.031996$.

If $e>10^{-5}N$, the cutoff-$131$ two-coincidence estimate gives
\[
 \frac{|A|}{N}<0.039772.
\]
Otherwise the single odd class contains more than $0.001177N$ pivots; two
lie within $850<1764$, so at most one residue pair coincides, and the
cutoff-$131$ estimate gives
\[
 \frac{|A|}{N}<0.039129.
\]
All six bounds are uniform on the interval and lie strictly below
$1/25-7/(25N)$.  Hence every allocation directly contradicts the original
Hall balance.
\end{proof}

\begin{lemma}[Terminal cutoff comparison]
\label{lem:terminal-cutoff-comparison}
Put $N_\dagger=32530000000$.  Consider the following family in the
dense-$E_1$ branch: an arbitrary prime cutoff, any valid odd individual
lower Bonferroni degree and even joint upper degree (including full
inclusion--exclusion), separate zero/one-coincidence endpoint counts, an
arbitrary integer split $k\geq2$ with $Y$ replaced by its lower bound
$N/(k+1)$ in the root term, and the exact diagonal sum.
At $N=N_\dagger$ its unique minimum uses cutoff $509$, Bonferroni degrees
$3/2$, and $k=117$, and the resulting envelope equals
\[
 0.040000010216\ldots
 >\frac1{25}-\frac7{25N_\dagger}.
\]
No larger cutoff or split parameter gives a smaller envelope.
\end{lemma}

\begin{proof}
At $N_\dagger$ the exact values are
\[
 \lfloor r_N+1\rfloor=24456856,\quad
 \mathcal S(24456856)=20767164,
\]
and
\[
 \lfloor N/117\rfloor=278034188,\qquad
 \pi(278034188)=15125057.
\]
For cutoffs through $509$, exact enumeration of all $77223$ cutoff/order
triples is minimized by cutoff $509$ and degrees $3/2$.  Exact enumeration
of $2\leq k\leq800$ is minimized by $k=117$; for $k\geq801$, the last
root-tail term alone exceeds
\[
 \frac{2048(802)^2}{N_\dagger}>0.0405.
\]

It remains to make the cutoff enumeration terminal.  For the minimizing
degree-$3/2$ pair, adjoining a new noncoincident prime $p$ when $n$ primes
are already present changes the endpoint-plus-density-plus-square-tail
payment by
\[
 \frac{2(n^2+5n+2)}{N_\dagger}
 -\frac{4}{25p^2}
 \left(1-\left(S_1+S_2-\frac19\right)\right),
\]
where $S_1=\sum q^{-2}$ and $S_2=\sum_{q<q'}q^{-2}q'^{-2}$ over the current
window.  At the first step, $n=95,p=521$ and
$S_1+S_2-1/9>0.05754$, so this increment is greater than
$2.87\cdot10^{-8}$.  Thereafter its positive endpoint term increases while
the magnitude of its negative term decreases.  The eight other potentially
competitive degree pairs are checked exactly through cutoff $4999$.  For
degree $1/0$ the finite-density increment exactly cancels the square-tail
saving, leaving a positive endpoint increment.  For the other seven, beyond
$4999$ the endpoint increment exceeds the maximal possible $4/(25p^2)$
density-tail saving.  Higher degrees are already worse at cutoff $509$, and
the same increment comparison only strengthens thereafter.

Finally
\[
 \frac1{25}-\frac7{25N_\dagger}
 =0.039999999991\ldots,
\]
which proves the stated comparison.
\end{proof}

\begin{lemma}[Exact split-point prime--root block]
\label{lem:exact-split-prime-root}
Put $Y=\lfloor N/118\rfloor$.  On
\[
 3.2275\cdot10^{10}\leq N\leq3.2532\cdot10^{10},
\]
the combined intermediate-prime and transformed-root payment satisfies
\begin{align*}
 &\frac{4\pi(Y)}N+\frac{2048}{N}\left(
 \frac{2(N/25+2)}Y+\frac{N^2+25N+1}{Y^2}\right)\\
 &\hspace{35mm}<0.002729824014.
\end{align*}
Also
\[
 \frac{\mathcal S(\lfloor r_N+1\rfloor)}N<0.000639920595.
\]
With $0.000639921$ in the exact diagonal formulas, the five left-endpoint
errors are below
\[
 0.002977624,\quad0.002976905,\quad0.002976546,
 \quad0.004897029,\quad0.004896851.
\]
\end{lemma}

\begin{proof}
For the prime--root term divide
\[
 273516949\leq Y\leq275694915
\]
into blocks $u\leq Y\leq v$ of length $10$.  On each block use
$\pi(v)$ and the least mapped $N$ in the prime term.  In the root term use
$Y\geq u$ and the largest mapped $N$: after expansion its only increasing
part is $2048N/u^2$, so this is a valid upper bound.  The maximum block is
\[
 u=273517049,\quad v=273517058,
\]
with
\[
 \pi(v)=14892354,\quad
 32275011782\leq N\leq32275012961,
\]
and gives $0.0027298240131\ldots$.

The diagonal range is $24328878\leq\lfloor r_N+1\rfloor\leq24457858$.
The length-$10$ hyperbola scan is maximal at
\[
 \frac{\mathcal S(24328897)}{32275018458}
 =\frac{20653449}{32275018458}<0.000639920595.
\]
Substitution gives the five displayed errors.
\end{proof}

Define
\[
 \mathfrak V^{509}_{118}
 :=\frac{4\tau_{509}^+}{25}+0.002729824014
 <0.002772511214.
\]

\begin{theorem}[Exact split-point descent]
\label{thm:32275e10-close}
The exact bound \eqref{eq:target}, and equivalently the Hall inequality \eqref{eq:hall-defect},
holds for every $N\geq3.2275\cdot10^{10}$.
\end{theorem}

\begin{proof}
Theorem~\ref{thm:3255e10-close} leaves
\[
 3.2275\cdot10^{10}\leq N<3.2532\cdot10^{10}.
\]
At the left endpoint the matching margin exceeds $1.8\cdot10^7$ and
$\mathcal R(N)>0.001162N$.  Lemmas~\ref{lem:conditional-509-window} and
\ref{lem:exact-split-prime-root} give the first four allocation bounds
\[
 0.039999953,\qquad0.038546,\qquad0.037688,\qquad0.034969.
\]
The dense $E_1$ low--low alternative is
\[
 \frac{|A|}{N}<\frac{23}{25}C_{\rm quad}+\Delta_{25}^{\rm ex,Y}
 +\frac{2\widehat J^{509}_1}{25}
 +\frac{2\zeta^{509}_1}{N}+\mathfrak V^{509}_{118}
 <0.039999953,
\]
where $\Delta^{\rm ex,Y}$ uses the diagonal constant from
Lemma~\ref{lem:exact-split-prime-root}; the high--low alternative remains
smaller.  If $e>10^{-5}N$, the even branch is below $0.039800$.  Otherwise
more than $0.001152N$ pivots remain in the single odd class, so a pair lies
within $869<1764$ and the last branch is below $0.039145$.  Every bound is
uniform and strictly below $1/25-7/(25N)$, closing the original Hall balance.
\end{proof}


\begin{lemma}[Boundary-corrected prime-square tail]
\label{lem:boundary-corrected-square-tail}
One has
\[
 \sum_{p>509}\frac1{p^2}<0.000265420=:\tau_{509}^{++}.
\]
\end{lemma}

\begin{proof}
Exact rational summation through $10^7$ gives
\[
 \sum_{509<p\leq10^7}\frac1{p^2}<0.000265410358.
\]
Partial summation retains the usually discarded lower boundary:
\[
 \sum_{p>X}\frac1{p^2}
 =-\frac{\pi(X)}{X^2}+2\int_X^\infty\frac{\pi(t)}{t^3}\,dt.
\]
With $X=10^7$, $\pi(X)=664579$, and
$\pi(t)<1.25506t/\log t$, this is below
\[
 -\frac{664579}{10^{14}}
 +\frac{2(1.25506)}{10^7\log(10^7)}
 <0.000000008928.
\]
The two strict bounds sum to less than $0.000265420$.
\end{proof}

\begin{lemma}[Final boundary-tail block]
\label{lem:boundary-tail-block}
On
\[
 3.22715\cdot10^{10}\leq N\leq3.2275\cdot10^{10},
\]
the exact-$Y$ combined prime--root term with $k=118$ is below
$0.002729928779$, and
\[
 \frac{\mathcal S(\lfloor r_N+1\rfloor)}N<0.000639941094.
\]
At the left endpoint the corresponding dense-$E_1$ envelope, using
$\tau_{509}^{++}$, is below $0.039999935$.
\end{lemma}

\begin{proof}
The same length-$10$ block calculation as
Lemma~\ref{lem:exact-split-prime-root} is maximal at
\[
 273487318\leq Y\leq273487327,
\quad \pi(273487327)=14890805,
\]
with
\[
 32271503524\leq N\leq32271504703.
\]
It gives $0.0027299287782\ldots$.  The diagonal scan is maximal at its
first block:
\[
 \frac{\mathcal S(24327128)}{32271500000}
 =\frac{20651859}{32271500000}<0.000639941094.
\]
Using $0.000639942$ in the diagonal formula, the degree-$3/2$ conditional
survivor of Lemma~\ref{lem:conditional-509-window}, and
Lemma~\ref{lem:boundary-corrected-square-tail} gives
$0.03999993450\ldots$.
\end{proof}

\begin{theorem}[Boundary-tail descent]
\label{thm:322715e10-close}
The exact bound \eqref{eq:target}, and equivalently the Hall inequality \eqref{eq:hall-defect},
holds for every $N\geq3.22715\cdot10^{10}$.
\end{theorem}

\begin{proof}
Theorem~\ref{thm:32275e10-close} leaves only the interval in
Lemma~\ref{lem:boundary-tail-block}.  The dense $E_1$ branch is below
$0.039999935$.  The other five bounds remain, respectively, below
\[
 0.038546,\quad0.037688,\quad0.034969,\quad0.039800,\quad0.039145.
\]
All are uniform and strictly below $1/25-7/(25N)$, so the original Hall
balance closes directly.
\end{proof}


\begin{lemma}[Root count versus inactive $509$-events]
\label{lem:root-inactive-509}
Let two pivots have at most $r$ coincident residue pairs in
$\mathcal P_{509}$, where $r=1,2,3,4$.  For two low-root pivots the
degree-$3/2$ density and endpoint pairs are bounded by
\[
\begin{array}{c|cccc}
r&1&2&3&4\\ \hline
I^{LL}_r&0.113400471&0.130297684&0.137108297&0.142005658\\
\epsilon^{LL}_r&303811&303623&303436&303250.
\end{array}
\]
If one pivot has the upper root count and the other is low-root, the bounds
improve to
\[
\begin{array}{c|ccc}
r&1&2&3\\ \hline
I^{HL}_r&0.111673122&0.114566348&0.115175656\\
\epsilon^{HL}_r&269183&269003&268824.
\end{array}
\]
For $r=4$, reactivating every event of a high-root pivot shows that the
last low--low pair is also a valid upper bound for every mixed root
alternative.  For two valuation-one upper-root pivots at distance at most $834$, the
$r=1$ bounds improve further to
\[
 I^{HH}_1<0.008402835,\qquad \epsilon^{HH}_1\leq234616.
\]
\end{lemma}

\begin{proof}
For a valuation-one pivot, root count $512$ requires at least nine odd
prime factors in the effective modulus.  At most one is $5$, so the pivot
has at least eight non-$5$ odd prime divisors.  Two of those cannot both
exceed $509$, since already
\[
 2(521)(523)(3)(7)(11)(13)(17)(19)>5.28\cdot10^{11}.
\]
Thus at least seven events in $\mathcal P_{509}$ are inactive.
The same eight-inactive-event relaxation is valid for a root-$4096$
high pivot in the higher $2$-adic classes.

At the upper endpoint $3.22715\cdot10^{10}$, enumerate the eight-subsets
$D\subseteq\mathcal P_{509}$ with $2\prod_{p\in D}p\leq N$ and the
seven-subsets with room for one further prime at least $521$.  There are
$9223$ resulting active-event patterns.  For each pattern evaluate the
individual $L_3$ polynomials and the joint $U_2$ polynomial, with the
$r$ smallest active common primes coincident.  The high--low maxima give
the second table; for $r=1$ the maximizing inactive set is
\[
 \{7,11,13,19,23,29,31,41\}.
\]
The first table is the exact full-active evaluation.  The fourth column is
obtained by making the four smallest active local counts equal to one; exact
rational evaluation gives the displayed density and endpoint count.  Any
forced inactive event can only decrease this full-active upper bound.

If both valuation-one pivots are high, every common odd divisor divides
their even difference, so the product of the common divisors is at most
$417$.  Exact evaluation of all gap-compatible pattern pairs gives the
last two bounds.  Without this small-gap condition, reactivating every
event of the second pivot can only enlarge both the intersection upper
bound and its endpoint count, so the high--low table is a valid upper bound
for the other high--high uses below.
\end{proof}

\begin{lemma}[Uniform root--inactive interval block]
\label{lem:root-inactive-interval-block}
On
\[
 2.631\cdot10^{10}\leq N\leq3.22715\cdot10^{10},
\]
the exact prime--root term with $k=121$ and root sum $768$ is below
$0.002678109464$, and
\[
 \frac{\mathcal S(\lfloor r_N+1\rfloor)}N<0.000680804714.
\]
With $0.000680805$ in the diagonal formulas, the five left-endpoint errors
are below
\[
 0.003165882,\quad0.003165064,\quad0.003164655,
 \quad0.005207890,\quad0.005207687.
\]
\end{lemma}

\begin{proof}
For $Y=\lfloor N/121\rfloor$, use length-$1000$ blocks exactly as in
Lemma~\ref{lem:exact-split-prime-root}.  The maximum is the first block:
\[
 217438016\leq Y\leq217439015,\quad
 \pi(217439015)=11989378,
\]
with $26310000000\leq N\leq26310120935$, and gives
$0.0026781094634\ldots$.  The length-$1000$ diagonal scan is also maximal
on its first block:
\[
 \frac{\mathcal S(21231414)}{26310000000}
 =\frac{17911972}{26310000000}<0.000680804714.
\]
Substitution gives the five errors.
\end{proof}

\begin{theorem}[Root--inactive-event descent]
\label{thm:2631e10-close}
The exact bound \eqref{eq:target}, and equivalently the Hall inequality \eqref{eq:hall-defect},
holds for every $N\geq2.631\cdot10^{10}$.
\end{theorem}

\begin{proof}
Theorem~\ref{thm:322715e10-close} leaves the interval of
Lemma~\ref{lem:root-inactive-interval-block}.  In the dense $E_1$ branch,
split according to whether each selected pivot has root count at most $256$
or exactly $512$.  The low--low, high--low, and high--high alternatives,
using Lemma~\ref{lem:root-inactive-509}, are respectively below
\[
 0.039826,\qquad0.039999725,\qquad0.032014.
\]
The high--low alternative uses the exact block in the preceding lemma and
is the largest.  The next three allocation bounds are
\[
 0.039232,\qquad0.038464,\qquad0.035494.
\]

It remains to treat the one-odd-class allocation.  Put
$e=|E_1|+|E_2|+|E_3|$.  If $e>10^{-5}N$, one valuation class contains more
than $e/3$, hence supplies a pair within $300000$.  Four coincidences would
force $2(3\cdot7\cdot11\cdot13)^2>1.8\cdot10^7$ to divide the difference,
so $r\leq3$.  Splitting each pivot at root count $2048$ and applying the
$r=3$ rows of Lemma~\ref{lem:root-inactive-509} gives three root
alternatives, all below $0.039743$.

If $e\leq10^{-5}N$, the mixed strip leaves more than $0.000476N$ pivots in
the single odd class.  Two lie within $2101<4(3\cdot7\cdot11)^2$, so
$r\leq2$.  Splitting at root count $256$ and using the $r=2$ rows gives
all three root alternatives below $0.039758$.  Thus all six exhaustive
allocations lie strictly below $1/25-7/(25N)$ and directly close the
original Hall balance.
\end{proof}


\begin{lemma}[Unit-block root--inactive endpoint]
\label{lem:unit-root-inactive-block}
On $2.6305\cdot10^{10}\leq N\leq2.631\cdot10^{10}$, unit $Y$-blocks give
the high--low prime--root term below $0.002678282125$, while unit diagonal
blocks give
\[
 \frac{\mathcal S(\lfloor r_N+1\rfloor)}N<0.000680809352.
\]
The dense $E_1$ branch is therefore below $0.039999978$.
\end{lemma}

\begin{proof}
The prime--root maximum occurs at
\[
 Y=217396723,\quad \pi(Y)=11987194,\quad
 26305003483\leq N\leq26305003603.
\]
The diagonal maximum is
\[
 \frac{\mathcal S(21227725)}{26305000000}
 =\frac{17908690}{26305000000}<0.000680809352.
\]
Substitution in the high--low formula of
Theorem~\ref{thm:2631e10-close} gives $0.03999997766\ldots$.
\end{proof}

\begin{theorem}[Terminal root--inactive descent]
\label{thm:26305e10-close}
The exact bound \eqref{eq:target}, and equivalently the Hall inequality \eqref{eq:hall-defect},
holds for every $N\geq2.6305\cdot10^{10}$.
\end{theorem}

\begin{proof}
Theorem~\ref{thm:2631e10-close} leaves only the interval in
Lemma~\ref{lem:unit-root-inactive-block}.  Its dense $E_1$ bound is below
$0.039999978$; the other five bounds in the preceding theorem have much
larger slack and decrease.  This closes every allocation in the original
Hall balance.
\end{proof}


\begin{lemma}[Dynamic cutoff--pattern endpoint]
\label{lem:dynamic-491-pattern-block}
On
\[
 2.6262\cdot10^{10}\leq N\leq2.6305\cdot10^{10},
\]
the high--low valuation-one branch admits the cutoff-$491$ bounds
\[
 I^{HL}_1<0.111634279,
 \qquad \epsilon^{HL}_1\leq244029,
 \qquad \sum_{p>491}\frac1{p^2}<0.000277248.
\]
For $Y=\lfloor N/121\rfloor$, the exact intermediate-prime and root payment
is below $0.002679811009$, while
\[
 \frac{\mathcal S(\lfloor r_N+1\rfloor)}N<0.000681147190.
\]
The resulting dense $E_1$ envelope is below $0.039999965$.
\end{lemma}

\begin{proof}
A root-$512$ valuation-one pivot has at least eight non-$5$ odd prime
divisors.  At most one can exceed $491$, since
\[
 2(499)(503)(3)(7)(11)(13)(17)(19)=486918618186
 >2.6305\cdot10^{10}.
\]
At the upper endpoint enumerate the eight-subsets of
$\mathcal P_{491}$ with $2\prod p\leq N$ and the seven-subsets with room for
one further prime at least $499$.  There are $6167$ distinct feasible
inactive-event patterns.  Exact evaluation of the individual degree-$3$
and joint degree-$2$ polynomials, including their endpoint counts, is
maximized by
\[
 D=\{7,11,13,17,23,29,31,37\},
 \qquad 2\prod_{p\in D}p=26037677666,
\]
and gives the first two displayed bounds.  Exact summation from $491$ through
$10^7$, followed by the boundary-corrected tail estimate of
Lemma~\ref{lem:boundary-corrected-square-tail}, gives the third.

Unit $Y$-blocks are maximal at $Y=217041322$; after subtracting the
$94$ primes already paid by the finite cutoff, the exact prime count is
$11968438$.  The corresponding normalized prime--root payment is below
$0.002679811009$.  Unit diagonal blocks are maximal at
\[
 \frac{\mathcal S(21204587)}{26262002235}
 =\frac{17888289}{26262002235}<0.000681147190.
\]
Using this ratio in the exact diagonal formula gives
$\Delta_{25}<0.003167512089$.  Substitution of all the preceding strict
bounds gives $0.03999996408\ldots$.
\end{proof}

\begin{theorem}[Dynamic-pattern descent]
\label{thm:26262e10-close}
The exact bound \eqref{eq:target}, and equivalently the Hall inequality \eqref{eq:hall-defect},
holds for every $N\geq2.6262\cdot10^{10}$.
\end{theorem}

\begin{proof}
Theorem~\ref{thm:26305e10-close} leaves only the interval of
Lemma~\ref{lem:dynamic-491-pattern-block}.  The dense $E_1$ branch is below
$0.039999965$.  The other five allocation bounds in
Theorem~\ref{thm:2631e10-close} have larger slack.  The only structural
change at the new endpoint is that the mixed strip is $>0.000486N$; after
removing an even part of size at most $10^{-5}N$, the final odd class still
contains more than $0.000476N$ pivots and supplies a pair within $2101$.
Thus the same at-most-two-coincidence argument applies, and all six
allocations close directly in the original Hall balance.
\end{proof}


\begin{lemma}[Quadratic-product endpoint block]
\label{lem:quadratic-product-endpoint-block}
On
\[
 2.6260\cdot10^{10}\leq N\leq2.6262\cdot10^{10},
\]
the cutoff-$491$ prime--root payment is below $0.002679890050$, and
\[
 \frac{\mathcal S(\lfloor r_N+1\rfloor)}N<0.000681163934.
\]
With $C_{\rm quad}<0.027346508$, the dense $E_1$ envelope is below
$0.039999953$.
\end{lemma}

\begin{proof}
The inactive-pattern family and its density/endpoint pair are those of
Lemma~\ref{lem:dynamic-491-pattern-block}.  Unit $Y$-blocks are maximal at
\[
 Y=217024793,\qquad \pi(Y)-\pi(491)=11967617,
\]
and give $0.0026798900498\ldots$.  Unit diagonal blocks are maximal at
\[
 \frac{\mathcal S(21203533)}{26260044182}
 =\frac{17887395}{26260044182}<0.000681163934.
\]
The exact diagonal formula is then below $0.003167587519$.  Substitution,
now using the sharpened quadratic Euler product, gives
$0.03999995253\ldots$.
\end{proof}

\begin{theorem}[Quadratic-product descent]
\label{thm:26260e10-close}
The exact bound \eqref{eq:target}, and equivalently the Hall inequality \eqref{eq:hall-defect},
holds for every $N\geq2.6260\cdot10^{10}$.
\end{theorem}

\begin{proof}
The preceding theorem leaves only the interval of
Lemma~\ref{lem:quadratic-product-endpoint-block}.  Its dense branch is below
$0.039999953$; the other five allocations retain larger slack.  The
mixed-strip and coincidence argument used there
is unchanged.  Hence every allocation directly contradicts the original
Hall balance.
\end{proof}


\begin{lemma}[Second root--inactive correlation]
\label{lem:second-root-inactive-correlation}
Let a high valuation-one pivot have root count $512$, and split the other
pivot according as its root count is at most $128$ or equals $256$.  In the
first case the cutoff-$491$ density/endpoint pair remains
\[
 (I,\epsilon)<(0.111634279,244029),
\]
but the root sum is only $640$.  In the second case one has
\[
 (I,\epsilon)<(0.111612165,231221),
\]
with root sum $768$.
\end{lemma}

\begin{proof}
Only the second case is new.  For a valuation-one pivot, root count greater
than $128$ requires at least eight distinct odd prime factors in the effective
modulus.  At most one is the modulus prime $5$, so at least seven non-$5$
odd primes divide the pivot.  At least three of them are at most $47$, since
otherwise the least possible product is
\[
 2(3)(7)(53)(59)(61)(67)(71)=38110106118
 >2.6260\cdot10^{10}.
\]
Thus at least three events in
\[
 \{3,7,11,13,17,19,23,29,31,37,41,43,47\}
\]
are inactive for the root-$256$ form.

At the upper endpoint there are $6143$ feasible seven/eight-inactive patterns
for the root-$512$ form.  Reactivate every event of the root-$256$ form except
for an arbitrary three-element subset of the displayed thirteen primes.
Exact degree-$3/2$ evaluation of all
\[
 6143\binom{13}{3}=1756898
\]
relaxed pattern pairs is maximized by the high inactive set
\[
 \{7,11,13,17,23,29,31,37\}
\]
and the low inactive set $\{41,43,47\}$, with the prime-$3$ residue pair
coincident.  This gives the second density/endpoint pair.  Since all other
forced low-form inactive events were reactivated, it is a uniform upper bound.
\end{proof}

\begin{lemma}[Second-correlation endpoint block]
\label{lem:second-correlation-endpoint-block}
On
\[
 2.6222\cdot10^{10}\leq N\leq2.6260\cdot10^{10},
\]
the root-sum-$640$ and root-sum-$768$ prime--root payments are respectively
below
\[
 0.002538247876,\qquad0.002681279433,
\]
and
\[
 \frac{\mathcal S(\lfloor r_N+1\rfloor)}N<0.000681458623.
\]
The two high--low alternatives are below $0.039859701$ and $0.039999987$.
\end{lemma}

\begin{proof}
Both unit $Y$-block maxima occur at
\[
 Y=216710779,\qquad \pi(Y)-\pi(491)=11951269.
\]
Substitution with root sums $640$ and $768$ gives the first two numbers.
The unit diagonal maximum is
\[
 \frac{\mathcal S(21183049)}{26222000000}
 =\frac{17869208}{26222000000}<0.000681458623,
\]
which gives $\Delta_{25}<0.003168950774$.  Apply
Lemma~\ref{lem:second-root-inactive-correlation}, the cutoff-$491$ tail, and
$C_{\rm quad}<0.027346508$ to obtain
$0.03985970054\ldots$ and $0.03999998609\ldots$.
\end{proof}

\begin{theorem}[Second root--inactive descent]
\label{thm:26222e10-close}
The exact bound \eqref{eq:target}, and equivalently the Hall inequality \eqref{eq:hall-defect},
holds for every $N\geq2.6222\cdot10^{10}$.
\end{theorem}

\begin{proof}
Theorem~\ref{thm:26260e10-close} leaves the interval in the preceding lemma.
The low-root split closes the largest high--low branch in both
alternatives.  Every other allocation retains larger slack.  At the new
endpoint the mixed strip is $>0.000481N$; after removing an even part of size
at most $10^{-5}N$, the final odd class has more than $0.000471N$ pivots and
supplies a pair within $2124<4(3\cdot7\cdot11)^2$.  Thus its
at-most-two-coincidence argument is unchanged, and all six allocations close
the original Hall balance.
\end{proof}


\begin{lemma}[Four-event low-root refinement]
\label{lem:four-event-low-root-refinement}
In the root-$256$ alternative of
Lemma~\ref{lem:second-root-inactive-correlation}, one may replace its
density/endpoint pair by
\[
 (I,\epsilon)<(0.111609892,227133).
\]
\end{lemma}

\begin{proof}
The seven non-$5$ odd prime divisors of the low pivot include at least four
primes at most $79$.  Indeed, if at most three did, their least possible
product would give
\[
 2(3)(7)(11)(83)(89)(97)(101)=33435142818
 >2.6222\cdot10^{10}.
\]
They still include at least three primes at most $47$.  Reactivate every
other low-form event and enumerate the $2717$ four-subsets of the primes
through $79$ having at least three members at most $47$.  At the upper
endpoint there are $6125$ feasible high-form patterns, so the enumeration has
$16641625$ pairs.  Its maximum uses high inactive set
\[
 \{7,11,13,17,23,29,31,37\}
\]
and low inactive set $\{41,43,47,79\}$, with the prime-$3$ pair coincident.
Exact degree-$3/2$ evaluation gives the displayed bound.
\end{proof}

\begin{lemma}[Four-event endpoint block]
\label{lem:four-event-endpoint-block}
On
\[
 2.6216\cdot10^{10}\leq N\leq2.6222\cdot10^{10},
\]
the root-sum-$768$ prime--root payment is below $0.002681504757$, and
\[
 \frac{\mathcal S(\lfloor r_N+1\rfloor)}N<0.000681507463.
\]
The root-$256$ high--low envelope is below $0.039999944$.
\end{lemma}

\begin{proof}
The unit prime--root maximum occurs at
\[
 Y=216661157,\qquad \pi(Y)-\pi(491)=11948721,
\]
and the diagonal maximum is
\[
 \frac{\mathcal S(21179821)}{26216006421}
 =\frac{17866404}{26216006421}<0.000681507463.
\]
The latter gives $\Delta_{25}<0.003169173036$.  Substitution of
Lemma~\ref{lem:four-event-low-root-refinement} gives
$0.03999994397\ldots$; the root-$128$ alternative remains far smaller.
\end{proof}

\begin{theorem}[Four-event descent]
\label{thm:26216e10-close}
The exact bound \eqref{eq:target}, and equivalently the Hall inequality \eqref{eq:hall-defect},
holds for every $N\geq2.6216\cdot10^{10}$.
\end{theorem}

\begin{proof}
The preceding theorem and lemma leave only the displayed interval and close
its controlling root-$256$ branch.  At the new endpoint the mixed strip is
$>0.000480N$; after the even part is removed, the last odd class supplies a
pair within $2128<4(3\cdot7\cdot11)^2$.  Every other allocation retains
larger slack, so the original Hall balance closes directly.
\end{proof}


\begin{lemma}[Five-event low-root refinement]
\label{lem:five-event-low-root-refinement}
In the root-$256$ alternative, one may use
\[
 (I,\epsilon)<(0.111609312,223134).
\]
\end{lemma}

\begin{proof}
At least five of the seven non-$5$ odd prime divisors are at most $157$:
if at most four were, the minimal product is already greater than
\[
 2(3)(7)(11)(13)(163)(167)(173)=28283653398
 >2.6216\cdot10^{10}.
\]
Combine this with the previously proved requirements
of at least four divisors at most $79$ and at least three at most $47$.

For fixed high-form data, write $x=p^{-2}$ for a low-form event that is
declared inactive.  Direct differentiation of the degree-$3$ individual
polynomial and degree-$2$ joint polynomial shows that the resulting
intersection upper bound is nonincreasing in $x$ for $0\leq x\leq1/9$;
the same check covers whether the high event at $p$ is active and whether the
least active common event changes.  Thus the componentwise largest admissible
relaxed five-set is the worst one.  It is
\[
 \{41,43,47,79,157\}.
\]
Evaluating this set against all $6121$ feasible high patterns is maximized by
the high inactive set
$\{7,11,13,17,23,29,31,37\}$ with the prime-$3$ pair coincident.  The exact
polynomials give the displayed density and endpoint count.
\end{proof}

\begin{lemma}[Five-event endpoint block]
\label{lem:five-event-endpoint-block}
On
\[
 2.6211\cdot10^{10}\leq N\leq2.6216\cdot10^{10},
\]
the root-sum-$768$ prime--root payment is below $0.002681688359$, and
\[
 \frac{\mathcal S(\lfloor r_N+1\rfloor)}N<0.000681546710.
\]
The controlling envelope is below $0.039999961$.
\end{lemma}

\begin{proof}
The unit prime--root maximum occurs at
\[
 Y=216620017,\qquad \pi(Y)-\pi(491)=11946587,
\]
and the unit diagonal maximum is
\[
 \frac{\mathcal S(21177127)}{26211004702}
 =\frac{17864024}{26211004702}<0.000681546710.
\]
Thus $\Delta_{25}<0.003169353914$.  With
Lemma~\ref{lem:five-event-low-root-refinement}, the envelope is
\[
 0.03999996020\ldots.
\]
\end{proof}

\begin{theorem}[Five-event descent]
\label{thm:26211e10-close}
The exact bound \eqref{eq:target}, and equivalently the Hall inequality \eqref{eq:hall-defect},
holds for every $N\geq2.6211\cdot10^{10}$.
\end{theorem}

\begin{proof}
The four-event theorem leaves the interval of the preceding lemma.  The
five-event bound closes its controlling root-$256$ alternative; every other
allocation retains larger slack.  The mixed strip remains $>0.000480N$, so
the final odd-class coincidence argument also remains unchanged.  This
closes the original Hall balance directly.
\end{proof}


\begin{lemma}[Nested six-event refinement]
\label{lem:nested-six-event-refinement}
At cutoff $503$, the root-$256$ high--low alternative admits
\[
 (I,\epsilon)<(0.111560667,234703),
\]
while the root-$128$ alternative admits
\[
 (I,\epsilon)<(0.111634809,260615)
\]
with root sum $640$.
\end{lemma}

\begin{proof}
The seven non-$5$ odd prime divisors of a root-$256$ pivot obey the nested
order constraints
\[
 d_1\leq17,\qquad d_2\leq29,\qquad d_3\leq47,\qquad
 d_4\leq79,\qquad d_5\leq157,\qquad d_6\leq503.
\]
For example, the six successive contrary minimal products are
\[
 51253692706, 30223218306, 38110106118,
 33435142818, 28283653398, 27076327278,
\]
each larger than $2.6211\cdot10^{10}$.  The last number is
$2(3)(7)(11)(13)(17)(509)(521)$.

The monotonicity calculation in
Lemma~\ref{lem:five-event-low-root-refinement} therefore reduces the relaxed
low inactive set to
\[
 \{17,29,47,79,157,503\}.
\]
At the upper endpoint, exact enumeration gives $6123$ feasible high patterns.
Evaluation against all of them is maximized by the high inactive set
$\{7,11,13,17,23,29,31,37\}$ with the prime-$3$ pair coincident, and gives
the first density/endpoint pair.  Reactivating every low event gives the
second pair.  A scan of every safe prime cutoff from $109$ through $1009$,
with its corresponding forced nested event count, is uniquely minimized at
$503$; beyond $1009$ the explicit endpoint increment exceeds the possible
tail decrement.
\end{proof}

\begin{lemma}[Nested six-event endpoint block]
\label{lem:nested-six-event-endpoint-block}
On
\[
 2.6153\cdot10^{10}\leq N\leq2.6211\cdot10^{10},
\]
the root-sum-$640$ and root-sum-$768$ prime--root payments are below
\[
 0.002540414524,\qquad0.002683823447,
\]
and
\[
 \frac{\mathcal S(\lfloor r_N+1\rfloor)}N<0.000682004015.
\]
The root-$128$/root-sum-$640$ and root-$256$/root-sum-$768$ high--low
envelopes are, respectively, below $0.039864461$ and $0.039999957$.
\end{lemma}

\begin{proof}
Exact summation and the negative boundary term give
\[
 \sum_{p>503}\frac1{p^2}<0.000269280.
\]
Both unit prime--root maxima occur at
\[
 Y=216140527,\qquad \pi(Y)-\pi(503)=11921650.
\]
The unit diagonal maximum is
\[
 \frac{\mathcal S(21145872)}{26153000000}
 =\frac{17836451}{26153000000}<0.000682004015,
\]
and gives $\Delta_{25}<0.003171459351$.
Substitution in the root-$128$ and root-$256$ cases of
Lemma~\ref{lem:nested-six-event-refinement} gives, respectively,
\[
 0.03986446071\ldots
 \qquad\hbox{and}\qquad
 0.03999995674\ldots.
\]
\end{proof}

\begin{theorem}[Nested six-event descent]
\label{thm:26153e10-close}
The exact bound \eqref{eq:target}, and equivalently the Hall inequality \eqref{eq:hall-defect},
holds for every $N\geq2.6153\cdot10^{10}$.
\end{theorem}

\begin{proof}
The five-event theorem leaves the interval of the preceding lemma, whose two
root alternatives close the controlling branch.  At the new endpoint the
mixed strip is $>0.000472N$; after removing the even part, the final odd class
still supplies a pair within $2165<4(3\cdot7\cdot11)^2$.  All other
allocations retain larger slack, so the original Hall balance closes directly.
\end{proof}


\begin{lemma}[Joint gap divisibility]
\label{lem:joint-gap-divisibility}
Let $b,c\in E_1$ be distinct, put $d=|b-c|\leq834$, and use the
cutoff-$503$ inactive sets $H,L$ for the two forms.  If
\[
 C=\prod_{p\in H\cap L}p
\]
and an active common prime $q$ has coincident residue classes, then
\[
 4Cq^2\mid d.
\]
In particular $4C\leq834$, and a coincident prime must be one of
$3,7,11,13$ and must satisfy $4Cq^2\leq834$.
\end{lemma}

\begin{proof}
Both valuation-one pivots are $2\pmod4$, so $4\mid d$.  Every prime in
$H\cap L$ divides both pivots and hence divides $d$.  A coincident active
residue pair at $q$ means $b\equiv c\pmod {q^2}$, so $q^2\mid d$.
The three factors $4,C,q^2$ are pairwise coprime: the cutoff primes are odd,
and an event cannot be simultaneously active and inactive.  Their product
therefore divides $d$.  The numerical conclusions follow from $d\leq834$.
\end{proof}

\begin{lemma}[Exact joint-compatible cutoff-$503$ patterns]
\label{lem:exact-joint-compatible-503}
In the valuation-one root-$512$/root-$256$ alternative on
\[
 2.3775\cdot10^{10}\leq N\leq2.6153\cdot10^{10},
\]
the degree-$3/2$ intersection density and endpoint discrepancy satisfy the
combined bound
\[
 \frac{2I^{HL}_{\rm joint}}{25}
 +\frac{2\epsilon^{HL}_{\rm joint}}N<0.001674991.
\]
The same maximizing pattern applies when the normalized endpoint term is
evaluated at either end of the interval.
\end{lemma}

\begin{proof}
At the upper endpoint there are $94$ cutoff primes.  Enumerate the root-$512$
patterns as the eight-subsets $H$ with $2\prod_{p\in H}p\leq N$, together
with the seven-subsets having room for one prime at least $509$.  There are
$6093$ such patterns.  For the root-$256$ pivot, enumerate the analogous
seven- and six-subsets; there are $560701$.

For fixed high data, Lemma~\ref{lem:five-event-low-root-refinement} and the
nested six-event argument give the relaxed low set
\[
 \{17,29,47,79,157,503\}.
\]
Evaluating this relaxation first excludes $5782$ high patterns whose upper
bound is already below the incumbent.  The remaining $311$ patterns require
\[
 311(560701)=174378011
\]
exact high--low evaluations.  Lemma~\ref{lem:joint-gap-divisibility} leaves
$31260793$ joint-compatible pairs.  For each one, the individual degree-$3$
and joint degree-$2$ polynomials are evaluated together with their actual
nonconstant CRT count.  The maximum is
\[
 \begin{split}
 H&=\{3,13,17,19,23,29,37,41\},\\
 L&=\{3,11,43,47,53,59,61\},
 \end{split}
\]
with $C=3$ and coincident prime $q=7$; indeed $4Cq^2=588\leq834$.
Exact rational evaluation gives
\[
 I=0.020694788706384609\ldots,
 \qquad \epsilon=230703.
\]
Repeating the combined density-plus-endpoint ordering at the lower endpoint
gives the same maximum.  The runner-up normalized payment is smaller by more
than $3.17\cdot10^{-7}$, while the first screened upper bound is smaller by
more than $8.2\cdot10^{-4}$, so the strict rounded bounds follow.
\end{proof}

\begin{lemma}[Joint-compatible endpoint block]
\label{lem:joint-compatible-endpoint-block}
On
\[
 2.3976\cdot10^{10}\leq N\leq2.6153\cdot10^{10},
\]
the six exhaustive Hall allocations are bounded, respectively, by
\[
 0.039985,\qquad0.0396,\qquad0.0389,\qquad0.0359,\qquad
 0.039999946,\qquad0.039921.
\]
The controlling fifth allocation is the one-odd-class case with
$e>10^{-5}N$, three allowed coincidences, and two low-root even pivots.
\end{lemma}

\begin{proof}
At the lower endpoint the extended half-root degree satisfies
\[
 \mathcal M_{1/2}(N)>\frac N{50}+2.07\cdot10^6,
 \qquad \mathcal R(N)>0.00017314N.
\]
The two root--Euler comparisons used in
Proposition~\ref{prop:half-root-correlated-degree} are, respectively, below
$0.001364$ and $0.005245$, still strictly below
$\kappa_5/200$ and $\kappa_5/75$.  Thus the matching and the six-way
allocation remain valid on the whole interval.

For the dense $E_1$ allocation, the low--low, exact joint high--low, and
high--high alternatives are below
\[
 0.039984900,\qquad0.032902569,\qquad0.032184709.
\]
The three intervening allocation envelopes increase by less than
$0.000322$ from their values in Theorem~\ref{thm:2631e10-close}; this gives
the second, third, and fourth displayed bounds.

If the even part in the final one-odd-class allocation exceeds $10^{-5}N$,
one valuation class supplies a pair within $300000$, so at most three
residue pairs coincide.  The low--low, high--low, and high--high root
alternatives are below
\[
 0.039999945687,\qquad0.038901921,\qquad0.039425069.
\]
For the controlling low--low case use cutoff $509$,
$(I^{LL}_3,\epsilon^{LL}_3)<(0.137108297,303436)$, root sum $4096$, and
$k=67$.  Unit blocks give
\[
 \begin{split}
 \text{prime--root payment}&<0.004739167218,\\
 \frac{\mathcal S(19955490)}{23976000000}
 &=\frac{16787505}{23976000000}<0.000700179555,\\
 \Delta_W&<0.005355245344.
 \end{split}
\]
Substitution gives $0.0399999456866949\ldots$.

Finally suppose $e\leq10^{-5}N$.  The single odd class contains more than
$0.00016314N$ pivots and hence supplies a pair within $6130$; since
$6130<4(3\cdot7\cdot11)^2$, at most two residue pairs coincide.  Its three
root alternatives are below
\[
 0.039920971,\qquad0.039306063,\qquad0.039573819.
\]
Every displayed bound is uniform and strictly below
$1/25-7/(25N)$.
\end{proof}

\begin{theorem}[Joint-compatible descent]
\label{thm:23976e10-close}
The exact bound \eqref{eq:target}, and equivalently the Hall inequality \eqref{eq:hall-defect},
holds for every $N\geq2.3976\cdot10^{10}$.
\end{theorem}

\begin{proof}
Theorem~\ref{thm:26153e10-close} leaves precisely the interval of
Lemma~\ref{lem:joint-compatible-endpoint-block}.  That lemma closes all six
mass allocations directly in the original Hall inequality.  No residual cut
is introduced.
\end{proof}


\begin{lemma}[Valuation--coincidence correlation]
\label{lem:valuation-coincidence-correlation}
In the final one-odd-class allocation with $e>10^{-5}N$, select two pivots
from one common $2$-adic valuation class at distance at most $300000$.
If that class is $E_1$, the three-coincidence low--low, high--low, and
high--high alternatives are below
\[
 0.037712106,\qquad0.036293211,\qquad0.036561762.
\]
If it is $E_2$ or $E_3$, at most two residue pairs coincide, and the three
corresponding alternatives are below
\[
 0.039482093,\qquad0.038881874,\qquad0.039406733.
\]
These bounds hold on
$2.3784\cdot10^{10}\leq N\leq2.3976\cdot10^{10}$.
\end{lemma}

\begin{proof}
For $E_1$, three coincidences are arithmetically possible, but the root split
is the valuation-one split at $256$: the three root sums are $512,768,1024$,
not $4096,6144,8192$.  Substitute those sums in the $r=3$ rows of
Lemma~\ref{lem:root-inactive-509}.

For $E_2$ and $E_3$, the two selected pivots are congruent modulo $8$.
Three coincident odd primes would therefore force
\[
 8(3\cdot7\cdot11)^2=426888
 \mid |b-c|,
\]
contrary to the gap bound.  Hence the $r=2$ rows apply.  Using the conservative
root sums $4096,6144,8192$ gives the second triple.  Direct endpoint
substitution gives the displayed values, and unit-block monotonicity makes
them uniform.  The value
\[
 0.039999945686\ldots
\]
combined mutually incompatible data: three coincidences from $E_1$ and root
sum $4096$ from a higher valuation class.
\end{proof}

\begin{lemma}[Valuation-correlated endpoint block]
\label{lem:valuation-correlated-endpoint-block}
On
\[
 2.3784\cdot10^{10}\leq N\leq2.3976\cdot10^{10},
\]
the six exhaustive Hall allocations are below, respectively,
\[
 0.039999979,\qquad0.0397,\qquad0.0389,
 \qquad0.0359,\qquad0.039483,\qquad0.039937.
\]
\end{lemma}

\begin{proof}
At the lower endpoint,
\[
 \mathcal M_{1/2}(N)>\frac N{50}+1.72\cdot10^6,
 \qquad \mathcal R(N)>0.00014499N,
\]
and the two root--Euler differences are below $0.001367$ and $0.005259$.
Thus the matching and allocation structure remain valid.

After Lemma~\ref{lem:valuation-coincidence-correlation}, the dense $E_1$
low--low branch is the only close alternative.  Use cutoff $509$,
$(I^{LL}_1,\epsilon^{LL}_1)<(0.113400471,303811)$, root sum $512$, and
$k=135$.  Exact unit blocks give
\[
 \begin{split}
 \text{prime--root payment}&<0.002438141663,\\
 \frac{\mathcal S(19848823)}{23784018997}
 &=\frac{16693717}{23784018997}<0.000701887978,\\
 \Delta_{25}&<0.003262997405.
 \end{split}
\]
The maximum prime--root block occurs at
$N=23784015555$, $Y=176177893$, and $\pi(Y)=9828932$.
Substitution gives
\[
 0.0399999788173336\ldots
 <\frac1{25}-\frac7{25N}.
\]
The exact joint-compatible high--low maximum of
Lemma~\ref{lem:exact-joint-compatible-503} is unchanged at the new lower
endpoint.  The three middle allocations retain their stated slack,
Lemma~\ref{lem:valuation-coincidence-correlation} gives the fifth, and the
three $r=2$ final-odd root alternatives are below
\[
 0.039936452,\qquad0.039322465,\qquad0.039591016.
\]
Here $e\leq10^{-5}N$ leaves more than $0.00013499N$ pivots in the final odd
class, hence a pair within $7408<4(3\cdot7\cdot11)^2$.  This completes all
six allocations.
\end{proof}

\begin{theorem}[Valuation-correlated descent]
\label{thm:23784e10-close}
The exact bound \eqref{eq:target}, and equivalently the Hall inequality \eqref{eq:hall-defect},
holds for every $N\geq2.3784\cdot10^{10}$.
\end{theorem}

\begin{proof}
Theorem~\ref{thm:23976e10-close} leaves precisely the interval of
Lemma~\ref{lem:valuation-correlated-endpoint-block}, whose six alternatives
all close in the original Hall balance.
\end{proof}


\begin{lemma}[Re-optimized low--low window]
\label{lem:reoptimized-low-low-window}
On
\[
 2.3775\cdot10^{10}\leq N\leq2.3784\cdot10^{10},
\]
the dense $E_1$ low--low branch is below
\[
 0.039999941.
\]
The optimal finite certificate uses cutoff $467$, individual degree $3$,
joint degree $2$, and split point $k=135$.
\end{lemma}

\begin{proof}
For an arbitrary prime cutoff $Q$, evaluate the odd individual lower
Bonferroni degrees $1,3,5$ and even joint upper degrees $0,2,4$, keeping the
zero- and one-coincidence endpoint counts separate.  Subtract the exact
finite prime-square sum from the boundary-corrected cutoff-$509$ tail and
retain $Y=\lfloor N/k\rfloor$ in the prime--root term.

At the lower endpoint, exact enumeration of every prime cutoff through
$4999$, every one of the nine degree pairs, and every $2\leq k\leq1000$ is
uniquely minimized at
\[
 Q=467,\qquad (d_{\rm ind},d_{\rm joint})=(3,2),\qquad k=135.
\]
The zero/one-coincidence densities and endpoint counts there are
\[
 \begin{array}{c|cc}
 &0&1\\ \hline
 I&0.013585721441&0.113397645522\\
 \epsilon&250980&250803.
 \end{array}
\]
The one-coincidence alternative controls.  Exact subtraction gives
\[
 \sum_{p>467}\frac1{p^2}<0.000289972.
\]
Higher Bonferroni degrees have larger endpoint counts and were already worse
at the larger endpoint in Lemma~\ref{lem:terminal-cutoff-comparison}; the
comparison only strengthens as $N$ decreases.  Beyond cutoff $4999$, the same
one-prime increment comparison used in that lemma has an increasing positive
endpoint term and a decreasing possible density-tail saving, so no omitted
cutoff can improve the bound.

Unit blocks on the stated interval give
\[
 \begin{split}
 \text{prime--root payment}&<0.002438481401,\\
 \frac{\mathcal S(19843811)}{23775011058}
 &=\frac{16689295}{23775011058}<0.000701967918,\\
 \Delta_{25}&<0.003263366199.
 \end{split}
\]
The first maximum occurs at $N=23775001605$, $Y=176111123$, and
$\pi(Y)=9825451$.  Substitution gives
\[
 0.0399999402040402\ldots
 <\frac1{25}-\frac7{25N}.
\]
All other Hall allocations retain the strict margins of
Lemma~\ref{lem:valuation-correlated-endpoint-block}.
\end{proof}

\begin{theorem}[Cutoff-$467$ descent]
\label{thm:23775e10-close}
The exact bound \eqref{eq:target}, and equivalently the Hall inequality \eqref{eq:hall-defect},
holds for every $N\geq2.3775\cdot10^{10}$.
\end{theorem}

\begin{proof}
Theorem~\ref{thm:23784e10-close} leaves only the interval of
Lemma~\ref{lem:reoptimized-low-low-window}, whose controlling branch and all
five complementary allocations close directly.
\end{proof}


\begin{lemma}[Integer negative-Pell count]
\label{lem:integer-negative-pell-count}
In every exact aggregate diagonal formula, the last factor
\[
 1+\log_{2+\sqrt3}N
\]
may be replaced by
\[
 P(N):=\left\lfloor1+\log_{2+\sqrt3}N\right\rfloor.
\]
In particular $P(N)=19$ throughout
\[
 2.2827\cdot10^{10}\leq N\leq2.3775\cdot10^{10}.
\]
\end{lemma}

\begin{proof}
The proof of Lemma~\ref{lem:aggregate-diagonal} bounds, for each fixed $d$,
the number of positive solutions of the negative Pell equation by
$1+\log_{2+\sqrt3}N$.  That number is an integer, so it is at most the floor
of the displayed real bound.  Finally direct squaring in
$\mathbb Q(\sqrt3)$ gives
\[
 (2+\sqrt3)^{18}<2.2827\cdot10^{10}
 <2.3775\cdot10^{10}<(2+\sqrt3)^{19},
\]
which makes the floor equal to $19$.
\end{proof}

\begin{lemma}[Integer-Pell endpoint block]
\label{lem:integer-pell-endpoint-block}
On
\[
 2.3669\cdot10^{10}\leq N\leq2.3775\cdot10^{10},
\]
the dense $E_1$ low--low branch is below $0.039999960$ and all five
complementary Hall allocations retain strict slack.
\end{lemma}

\begin{proof}
Use Lemma~\ref{lem:integer-negative-pell-count} in the exact diagonal term and
repeat the cutoff/order/split enumeration.  At the lower endpoint the unique optimum
is cutoff $463$, degrees $3/2$, and $k=135$.  Its one-coincidence data are
\[
 I=0.113397118126\ldots,
 \qquad \epsilon=242617,
 \qquad \sum_{p>463}\frac1{p^2}<0.000294557.
\]
With this fixed certificate, unit blocks give
\[
 \begin{split}
 \text{prime--root payment}&<0.002442402990,\\
 \frac{\mathcal S(19784791)}{23669021538}
 &=\frac{16637351}{23669021538}<0.000702916721,\\
 \Delta_{25}^{\rm Pell}&<0.003259369807.
 \end{split}
\]
The prime--root maximum occurs at $N=23669006085$, $Y=175325971$, and
$\pi(Y)=9784038$.  Thus the controlling envelope is
\[
 0.0399999596197116\ldots
 <\frac1{25}-\frac7{25N}.
\]

At the same endpoint the half-root matching margin exceeds
$1.51\cdot10^6$, the two root--Euler differences are below $0.001369$ and
$0.005268$, and $\mathcal R(N)>0.00012798N$.  If the final even part is at
most $10^{-5}N$, the remaining odd class therefore supplies a pair within
$8476<4(3\cdot7\cdot11)^2$.  Every structural and complementary allocation
used above remains valid; the integer Pell replacement only lowers their
diagonal payments.
\end{proof}

\begin{theorem}[Integer-Pell descent]
\label{thm:23669e10-close}
The exact bound \eqref{eq:target}, and equivalently the Hall inequality \eqref{eq:hall-defect},
holds for every $N\geq2.3669\cdot10^{10}$.
\end{theorem}

\begin{proof}
Theorem~\ref{thm:23775e10-close} leaves the interval of
Lemma~\ref{lem:integer-pell-endpoint-block}, which closes every allocation in
the original Hall inequality.
\end{proof}


\begin{lemma}[Restricted small-modulus diagonal sum]
\label{lem:restricted-small-modulus-sum}
Put
\[
 \mathcal F(R)=
 \sum_{\substack{s\leq R\ s\ {
m squarefree}\\
       p\mid s\Rightarrow p\equiv1\pmod4,\ p\geq13}}
       2^{\omega(s)}.
\]
In every exact aggregate diagonal formula, the occurrence of
$\mathcal S(R)$ in the small-modulus rounding term may be replaced by
$\mathcal F(R)$.  On
\[
 2.2827\cdot10^{10}\leq N\leq2.3669\cdot10^{10}
\]
unit blocks give
\[
 \frac{\mathcal F(19312757)}{22827035690}
 =\frac{3991879}{22827035690}<0.000174875050.
\]
With the integer Pell factor $19$, the corresponding maxima satisfy
\[
 \Delta_{25}^{\mathcal F}<0.001689240822,
 \qquad
 \Delta_W^{\mathcal F}<0.002213420548,
 \qquad
 \Delta_{100}^{\mathcal F}<0.001687899155.
\]
On the extended lower interval
\[
 1.382\cdot10^{10}\leq N\leq2.2827\cdot10^{10},
\]
the corresponding unit-block maxima are
\[
 \frac{\mathcal F(13821334)}{13820000000}
 =\frac{2856719}{13820000000}<0.000206709045
\]
and
\[
 \Delta_{25}^{\mathcal F}<0.001997137693,
 \qquad
 \Delta_W^{\mathcal F}<0.002616654327,
 \qquad
 \Delta_{100}^{\mathcal F}<0.001995296367.
\]
\end{lemma}

\begin{proof}
In the proof of Lemma~\ref{lem:aggregate-diagonal}, a modulus $s$ contributes
only when it is squarefree and every prime divisor is $1\pmod4$ and at least
$13$.  Its exact root count is $2^{\omega(s)}$.  Thus the small-$s$ rounding
loss is at most $3\mathcal F(R)$ in the $q$-formulas and
$6\mathcal F(R)$ in the $W$-formula.  Replacing this by
$\mathcal S(R)$ would enlarge both the support and the weight.

For completeness, the finite certificate starts with weight one at $s=1$.
From a permitted squarefree product $s$, append each permitted prime larger
than its last prime factor, provided $sp\leq19784778$, and assign weight
$2\cdot2^{\omega(s)}$ to $sp$.  Unique increasing prime factorization makes
this an exact, duplication-free enumeration.  There are $1224624$ products
at the upper end.  Prefixing their integer weights and scanning every change
of $\lfloor r_N+1\rfloor$ gives the displayed maximum.  Substitution in the
three diagonal formulas, with Lemma~\ref{lem:integer-negative-pell-count},
gives the last line.  Repeating the same prefix and unit-block scan on the
extended range gives the second set of maxima.  The conservative Pell factor
$19$ remains valid there (and drops to $18$ below
$(2+\sqrt3)^{18}$).
\end{proof}

\begin{lemma}[Restricted-diagonal endpoint block]
\label{lem:restricted-diagonal-endpoint-block}
On
\[
 2.2827\cdot10^{10}\leq N\leq2.3669\cdot10^{10},
\]
the six exhaustive Hall allocations are below, respectively,
\[
 0.038464,\qquad0.0389,\qquad0.0382,\qquad0.0363,
 \qquad0.038539,\qquad0.039711.
\]
\end{lemma}

\begin{proof}
At the lower endpoint, Proposition~\ref{prop:half-root-correlated-degree}
gives
\[
 \mathcal M_{1/2}(N)>\frac N{50}+411,
 \qquad \mathcal R(N)>0.000000036N.
\]
The root--Euler differences are below $0.001386$ and $0.005332$, so the
matching and the six-way allocation remain valid.

For the dense $E_1$ low--low branch retain the cutoff-$463$, degree-$3/2$,
$k=135$ certificate of Lemma~\ref{lem:integer-pell-endpoint-block}, but use
Lemma~\ref{lem:restricted-small-modulus-sum}.  Unit $Y$- and $R$-blocks give
\[
 \begin{split}
 \text{prime--root payment}&<0.002474825011,\\
 \Delta_{25}^{\mathcal F}&<0.001689240822,
 \end{split}
\]
and the full envelope is
\[
 0.038463008852\ldots.
\]
Direct substitution of the same restricted diagonal terms in the next three
allocation formulas gives the second, third, and fourth displayed bounds.

It remains to check the one-odd-class allocation at the almost-vanishing
strip.  Split at $e=1.5\cdot10^{-8}N$.  If $e$ is larger, one of the three
even valuation classes contains at least $115$ pivots and supplies a pair at
distance below $2.01\cdot10^8$.  If $e$ is smaller, the final odd class
contains at least $481$ pivots and supplies a pair at distance below
$4.8\cdot10^7$.  In both cases the difference is divisible by $4$, and five
coincident odd primes would force
\[
 4(3\cdot7\cdot11\cdot13\cdot17)^2
 =10424818404
\]
to divide it.  Hence at most four residue pairs coincide.
Lemma~\ref{lem:root-inactive-509} gives
\[
 I^{LL}_4<0.142005658,\qquad\epsilon^{LL}_4=303250.
\]
Reactivating every finite event is a valid common upper bound for all root
alternatives.  For the even pair, use the worst root sum $8192$ and fixed
$k=52$; for the odd pair use root sum $1024$ and $k=105$.  Unit blocks give
\[
 \begin{array}{c|cc}
 &\text{prime--root}&\text{full envelope}\\ \hline
 \text{even}&0.006026552133&0.038538552454\\
 \text{odd}&0.003090437113&0.039710563116.
 \end{array}
\]
The latter row includes the raw $1.5\cdot10^{-8}$ even charge.  All six
bounds are therefore strictly below $1/25-7/(25N)$ throughout the interval.
\end{proof}

\begin{theorem}[Restricted-diagonal descent]
\label{thm:22827e10-close}
The exact bound \eqref{eq:target}, and equivalently the Hall inequality \eqref{eq:hall-defect},
holds for every $N\geq2.2827\cdot10^{10}$.
\end{theorem}

\begin{proof}
Theorem~\ref{thm:23669e10-close} leaves precisely the interval of
Lemma~\ref{lem:restricted-diagonal-endpoint-block}, which closes every
allocation directly in the original Hall inequality.
\end{proof}


\begin{lemma}[Eight-prime endpoint block]
\label{lem:eight-prime-endpoint-block}
On
\[
 1.382\cdot10^{10}\leq N\leq2.2827\cdot10^{10},
\]
all six Hall allocations close directly.  The dense $E_1$ branch is below
$0.039366$, the worst higher-valuation even branch is below $0.039753$, and
the final odd branch is below $0.039577$.
\end{lemma}

\begin{proof}
At the lower endpoint, Proposition~\ref{prop:eight-prime-correlated-degree}
and Corollary~\ref{cor:thin-nonbase-strip} give
\[
 \mathcal M_{1/4}(N)>\frac N{50}+16356,
 \qquad \mathcal R(N)>0.00000236N.
\]
Use Lemma~\ref{lem:restricted-small-modulus-sum} in all diagonal terms.  The
preceding computation already covers $N\geq1.383\cdot10^{10}$; scan the new unit
interval down to $1.382\cdot10^{10}$.  For the cutoff-$463$, $k=135$ dense
certificate, the maxima are
\[
 \text{prime--root}<0.003055994166,\qquad
 \Delta_{25}^{\mathcal F}<0.001997137693,
\]
and the full envelope is $0.039365928857\ldots$.  Direct substitution in the
three middle allocation formulas leaves larger slack.

For the last two allocations split at $e=10^{-7}N$.  If $e$ is larger, one
even valuation class contains more than $e/3$ pivots and supplies a pair at
distance below $3.01\cdot10^7$.  Four coincidences would force
$4(3\cdot7\cdot11\cdot13)^2=36072036$ to divide the gap, so $r\leq3$.
For $E_1$ use the worst root sum $1024$; for $E_2,E_3$ use the conservative
root sum $8192$.  Unit blocks give
\[
 0.036240855052\ldots,\qquad0.039752885052\ldots.
\]

If $e\leq10^{-7}N$, the final odd class contains more than
$0.00000226N$ pivots and supplies a pair at distance below $4.43\cdot10^5$.
Four coincidences would force $4(3\cdot7\cdot11\cdot13)^2=36072036$ to
divide the gap, so $r\leq3$.  A low--low pair uses
$(I^{LL}_3,\epsilon^{LL}_3)<(0.137108297,303436)$ and root sum $512$.
If one pivot is high-root, the forced inactive events give
$(I^{HL}_3,\epsilon^{HL}_3)<(0.115175656,268824)$; use root sums $768$ and
$1024$ for the high--low and high--high alternatives.  On this new unit
interval the exact Pell factor is $18$.  Optimizing at fixed
$k=113,99,89$ gives, respectively,
\[
 0.039576187839\ldots,\qquad
 0.038656450868\ldots,\qquad
 0.038982482756\ldots.
\]
The higher-valuation even branch is controlling, and it too is strictly below
$1/25-7/(25N)$.  All six therefore close on the stated interval.
\end{proof}

\begin{theorem}[Eight-prime root--Euler descent]
\label{thm:1382e10-close}
The exact bound \eqref{eq:target}, and equivalently the Hall inequality \eqref{eq:hall-defect},
holds for every $N\geq1.382\cdot10^{10}$.
\end{theorem}

\begin{proof}
Theorem~\ref{thm:22827e10-close} leaves the interval of
Lemma~\ref{lem:eight-prime-endpoint-block}, whose six allocations all close
inside the original Hall inequality.
\end{proof}


\begin{lemma}[Second eight-prime endpoint block]
\label{lem:second-eight-prime-endpoint-block}
On
\[
 1.29\cdot10^{10}\leq N\leq1.382\cdot10^{10},
\]
all six Hall allocations close directly.  The dense branch is below
$0.039519$, the worst even branch is below $0.039926$, and the final odd
branch is below $0.039706$.
\end{lemma}

\begin{proof}
At the lower endpoint, Proposition~\ref{prop:second-eight-prime-degree} and
Corollary~\ref{cor:thin-nonbase-strip} give
\[
 \mathcal M_{1/8}(N)>\frac N{50}+6728,
 \qquad \mathcal R(N)>0.00000104N.
\]
The exact restricted-diagonal scan gives
\[
 \frac{\mathcal F(13200941)}{12900021174}
 =\frac{2728601}{12900021174}<0.000211519111,
\]
\[
 \Delta_{25}^{\mathcal F}<0.002043605517,
 \qquad \Delta_W^{\mathcal F,18}<0.002603499910.
\]
The cutoff-$463$ dense envelope is $0.039518076041\ldots$.

Split at $e=9\cdot10^{-8}N$.  In the dense-even case a valuation class
supplies a pair below $3.34\cdot10^7$, so $r\leq3$.  The $E_1$ alternative
is below $0.036401$; for $E_2,E_3$, use root sum $8192$, exact Pell factor
$18$, and $(I^{LL}_3,\epsilon^{LL}_3)$ to obtain
\[
 0.039925196317\ldots.
\]
If $e\leq9\cdot10^{-8}N$, the final odd class has more than
$0.00000095N$ pivots, hence a pair below $1.06\cdot10^6$ and again $r\leq3$.
The low--low, high--low, and high--high root-stratified envelopes are
\[
 0.039705845760\ldots,\qquad
 0.038796221886\ldots,\qquad
 0.039129213946\ldots.
\]
The first four allocation formulas retain larger slack.  Thus every branch is
strictly below $1/25-7/(25N)$.
\end{proof}

\begin{theorem}[Second eight-prime root--Euler descent]
\label{thm:129e10-close}
The exact bound \eqref{eq:target}, and equivalently the Hall inequality \eqref{eq:hall-defect},
holds for every $N\geq1.29\cdot10^{10}$.
\end{theorem}

\begin{proof}
Theorem~\ref{thm:1382e10-close} leaves the interval of
Lemma~\ref{lem:second-eight-prime-endpoint-block}, which closes every
allocation in the original Hall inequality.
\end{proof}


\begin{lemma}[Third eight-prime endpoint block]
\label{lem:third-eight-prime-endpoint-block}
On
\[
 1.279\cdot10^{10}\leq N\leq1.29\cdot10^{10},
\]
all six Hall allocations close directly.  The dense, worst even, and final
odd envelopes are below $0.039538$, $0.039957$, and $0.039723$.
\end{lemma}

\begin{proof}
At the lower endpoint,
\[
 \mathcal M_{1/16}(N)>\frac N{50}+5090,
 \qquad \mathcal R(N)>0.00000079N.
\]
Unit blocks give
\[
 \frac{\mathcal F(13125781)}{12790008075}
 =\frac{2713031}{12790008075}<0.000212121133,
\]
and the cutoff-$463$ dense envelope is
$0.039537470760\ldots$.  Split at $e=9\cdot10^{-8}N$.
The dense-even valuation class supplies a pair below $3.34\cdot10^7$, so
$r\leq3$.  With exact Pell factor $18$, the worst $E_2,E_3$ envelope is
\[
 0.039956442803\ldots.
\]
In the sparse-even case the final odd class contains more than
$0.00000070N$ pivots, hence a pair below $1.43\cdot10^6$ and again $r\leq3$.
The low--low, high--low, and high--high envelopes are
\[
 0.039722473493\ldots,\qquad
 0.038814112743\ldots,\qquad
 0.039147949978\ldots.
\]
All other allocations retain larger slack, proving the claim.
\end{proof}

\begin{theorem}[Third eight-prime root--Euler descent]
\label{thm:1279e10-close}
The exact bound \eqref{eq:target}, and equivalently the Hall inequality \eqref{eq:hall-defect},
holds for every $N\geq1.279\cdot10^{10}$.
\end{theorem}

\begin{proof}
Theorem~\ref{thm:129e10-close} leaves the interval of
Lemma~\ref{lem:third-eight-prime-endpoint-block}, which closes every Hall
allocation directly.
\end{proof}


\begin{lemma}[Retained-modulus endpoint block]
\label{lem:retained-modulus-endpoint-block}
On
\[
 1.2643\cdot10^{10}\leq N\leq1.279\cdot10^{10},
\]
all six Hall allocations close directly.  The dense, worst even, and final
odd envelopes are below $0.039564$, $0.039999135$, and $0.039745$.
\end{lemma}

\begin{proof}
At the lower endpoint, Proposition~
\ref{prop:modulus-correlated-eight-prime-degree} and Corollary~
\ref{cor:thin-nonbase-strip} give
\[
 \mathcal M_{1/32}(N)>\frac N{50}+1442,
 \qquad \mathcal R(N)>0.000000228N.
\]
The exact restricted-diagonal unit scan gives
\[
 \frac{\mathcal F(13025077)}{12643098938}
 =\frac{2692239}{12643098938}<0.000212941386,
\]
and the cutoff-$463$ dense envelope is
$0.039563962205\ldots$.  Split at $e=10^{-7}N$.  In the dense-even case one
valuation class supplies a pair below $3\cdot10^7$, so $r\leq3$.  With exact
Pell factor $18$, the worst $E_2,E_3$ envelope is
\[
 0.039999134344\ldots<\frac1{25}-\frac7{25N}.
\]
If $e\leq10^{-7}N$, the final odd class contains more than
$0.000000128N$ pivots and supplies a pair below $7.82\cdot10^6$, again with
$r\leq3$.  The low--low, high--low, and high--high envelopes are
\[
 0.039744956074\ldots,\qquad
 0.038838362547\ldots,\qquad
 0.039173393787\ldots.
\]
The remaining allocations have larger slack.  Thus every branch is strictly
below the exact target.
\end{proof}

\begin{theorem}[Retained-modulus root--Euler descent]
\label{thm:12643e10-close}
The exact bound \eqref{eq:target}, and equivalently the Hall inequality \eqref{eq:hall-defect},
holds for every $N\geq1.2643\cdot10^{10}$.
\end{theorem}

\begin{proof}
Theorem~\ref{thm:1279e10-close} leaves precisely the interval of
Lemma~\ref{lem:retained-modulus-endpoint-block}, which closes every
allocation in the original Hall inequality.
\end{proof}


\begin{lemma}[Quadratic-residue endpoint block]
\label{lem:qr-endpoint-block}
On
\[
 1.15\cdot10^{10}\leq N\leq1.2643\cdot10^{10},
\]
all six Hall allocations close directly.  The dense, worst even, and final
odd envelopes are below $0.039791$, $0.039896$, and $0.039534$.
\end{lemma}

\begin{proof}
At the lower endpoint, Proposition~\ref{prop:qr-screened-degree} and
Corollary~\ref{cor:thin-nonbase-strip} give
\[
 \mathcal M_{\rm QR}(N)>\frac N{50}+154290,
 \qquad \mathcal R(N)>0.0000268N.
\]
Twelve exact blocks of width at most $10^8$ cover the interval.  The
restricted-diagonal scan is controlled by the first block and gives
\[
 \frac{\mathcal F(12227654)}{11500000000}
 =\frac{2527449}{11500000000}<0.000219778174,
\]
with cutoff-$463$ dense envelope
$0.039790171170\ldots$.

Split at $e=1.5\cdot10^{-5}N$.  In the dense-even case one valuation class
has more than $5\cdot10^{-6}N$ pivots and supplies a pair below $200001$.
Three coincident active primes would make
\[
 4(3\cdot7\cdot11)^2=213444
\]
divide their difference, so $r\leq2$.  Exact Pell factor $18$, root sum
$8192$, and $(I^{LL}_2,\epsilon^{LL}_2)$ give the worst $E_2,E_3$ envelope
\[
 0.039895678510\ldots.
\]
If $e\leq1.5\cdot10^{-5}N$, the final odd class contains more than
$0.0000118N$ pivots, so a pair lies below $84746$ and again $r\leq2$.
The low--low, high--low, and high--high envelopes, using the exact
root sums $512,768,1024$, are
\[
 0.039533065858\ldots,\qquad
 0.039010120997\ldots,\qquad
 0.039358389351\ldots.
\]
All block maxima occur at the lower end, and the other allocations retain
larger slack.  Hence every branch is strictly below the exact target.
\end{proof}

\begin{theorem}[Quadratic-residue root-tail descent]
\label{thm:115e10-close}
The exact bound \eqref{eq:target}, and equivalently the Hall inequality \eqref{eq:hall-defect},
holds for every $N\geq1.15\cdot10^{10}$.
\end{theorem}

\begin{proof}
Theorem~\ref{thm:12643e10-close} leaves the interval of
Lemma~\ref{lem:qr-endpoint-block}, which closes every allocation directly in
the original Hall inequality.
\end{proof}


\begin{lemma}[QR-hierarchy endpoint block]
\label{lem:qr-hierarchy-endpoint-block}
On
\[
 9.38\cdot10^9\leq N\leq1.15\cdot10^{10},
\]
all six Hall allocations close directly.  The dense, worst even, and final
odd envelopes are below $0.039998505$, $0.038930$, and $0.039954$.
\end{lemma}

\begin{proof}
At the lower endpoint, Proposition~\ref{prop:qr-root-hierarchy} and
Corollary~\ref{cor:thin-nonbase-strip} give
\[
 \mathcal M_{\rm QR2}(N)>\frac N{50}+7633040,
 \qquad \mathcal R(N)>0.001627N.
\]
For the dense branch optimize the complete finite-window family while
retaining the exact restricted diagonal sum.  The lower block is uniquely
minimized by cutoff $367$, Bonferroni degrees $3/2$, and $k=100$, with
\[
 \frac{\mathcal F(10674474)}{9380000000}
 =\frac{2206435}{9380000000}<0.000235227612,
\]
and the full dense envelope is
\[
 0.039998504675\ldots<\frac1{25}-\frac7{25N}.
\]

Split at $e=1.5\cdot10^{-5}N$.  A dense even valuation class supplies a pair
below $200001$, so $r\leq2$ as before.  On the present interval a pivot cannot
have more than $2048$ transformed roots; consequently the two-pivot root sum
in the $E_2,E_3$ branch is at most $4096$.
With exact Pell factor $18$ its envelope is
\[
 0.038929045621\ldots.
\]
In the sparse-even case the final odd class contains more than $0.001612N$
pivots and supplies a pair below $621$, again with $r\leq2$.  The low--low,
high--low, and high--high envelopes are
\[
 0.039953388923\ldots,\qquad
 0.039460160705\ldots,\qquad
 0.039834609222\ldots.
\]
All later unit blocks decrease and the remaining allocations have larger
slack.  Hence every branch closes in the original Hall balance.
\end{proof}

\begin{theorem}[Quadratic-residue hierarchy descent]
\label{thm:938e9-close}
The exact bound \eqref{eq:target}, and equivalently the Hall inequality \eqref{eq:hall-defect},
holds for every $N\geq9.38\cdot10^9$.
\end{theorem}

\begin{proof}
Theorem~\ref{thm:115e10-close} leaves the interval of
Lemma~\ref{lem:qr-hierarchy-endpoint-block}, which closes every allocation
directly.
\end{proof}

\begin{lemma}[Cutoff-$383$ root--inactive endpoint block]
\label{lem:383-root-inactive-endpoint-block}
On
\[
 8.78\cdot10^9\leq N\leq9.38\cdot10^9,
\]
all six Hall allocations close directly.  The dense, worst even, and final
odd envelopes are below $0.039997644$, $0.039192926$, and $0.039954208$.
\end{lemma}

\begin{proof}
At the lower endpoint, Proposition~\ref{prop:qr-root-hierarchy} and
Corollary~\ref{cor:thin-nonbase-strip} give
\[
 \mathcal M_{\rm QR2}(N)>\frac N{50}+6261235,
 \qquad \mathcal R(N)>0.001426N.
\]
The higher-root lower bound in the proof of the proposition is still more
than $9.49\cdot10^6$ above $\mathcal M_{\rm QR2}$ there, so no degree stratum
has been lost by extending the endpoint.

It remains to strengthen the controlling dense $E_1$ branch.  Its selected
pair has even gap at most $834$ and at most one coincident finite-prime
residue pair.  A valuation-one pivot has root count at most $512$ on the
present interval.  First split each root count as at most $128$, exactly
$256$, or exactly $512$.

A root-$256$ pivot has at least seven non-$5$ odd prime divisors.  At most one
can exceed $383$, since the least product with two such primes already
exceeds the upper endpoint.  Thus at cutoff $383$ it forces either seven
inactive events with
\[
 2\prod_{p\in D}p\leq9380000000,
\]
or six inactive events with room for one further prime at least $389$.
There are respectively $129618$ and $713$ such forced patterns.  Exact
degree-$3/2$ evaluation gives the following coordinatewise density/endpoint
bounds for a full-active pair and a pair with one root-$256$ pivot:
\[
 (I_{128,128},\epsilon_{128,128})
 <(0.113388315,146003),
\]
\[
 (I_{256,128},\epsilon_{256,128})
 <(0.112531936,130146).
\]

For two root-$256$ pivots, let $C$ be the product of their common forced
inactive primes and let $q$ be the sole coincident active prime.  By
Lemma~\ref{lem:joint-gap-divisibility}, the exact compatibility condition is
$4Cq^2\leq834$.  Reactivating all events of either pivot gives an
upper bound for every pair whose one-form bound is too small to control.
Only $6955$ of the $130331$ patterns remain.  Exhausting their
gap-compatible Cartesian square gives
\[
 I_{256,256}<0.111792015,\qquad \epsilon_{256,256}\leq109217.
\]
The density maximum has common product $19$, coincidence prime
$3$, and inactive sets
\[
 \{7,17,19,31,37,41,43\},\qquad
 \{11,13,19,23,29,47,53\}.
\]
The complete enumeration is reproduced by
\path{scripts/dense_root_inactive_359_scan.py} with \texttt{--q 383}.

Now
\[
 \sum_{p>383}\frac1{p^2}<0.000371164574418.
\]
The exact restricted diagonal and prime blocks on the whole interval, with
root sums $256,384,512$ and split points $124,108,98$, give
\[
 0.039443871,\qquad0.039759223,\qquad0.039997643212.
\]
The controlling block has
\[
 \frac{\mathcal F(10214297)}{8780025882}
 =\frac{2111247}{8780025882}<0.000240460226.
\]
A root-$512$ pivot forces eight non-$5$ odd prime divisors.  At cutoff $359$
there are only $658$ feasible forced patterns; even after reactivating every
event of the other pivot one has
\[
 (I,\epsilon)<(0.021511201,105714).
\]
The worst root-sum-$1024$ envelope is therefore below $0.033652$, so these
higher-root alternatives close with large slack.

For the dense higher-valuation even branch, the pair gap remains below
$200001$, hence $r\leq2$, and the actual two-pivot root sum is $4096$.
The exact Pell-$18$ block is below $0.039192926$.  If the even mass is at
most $1.5\cdot10^{-5}N$, the final odd class contains more than
$0.001411N$ pivots and supplies a pair below $710$, still with $r\leq2$.
At cutoff $359$, splitting its two root counts at $128/256$ gives the three
exact block bounds
\[
 0.039301217,\qquad0.039647372,\qquad0.039954208.
\]
The remaining allocations retain larger slack.  Hence every branch closes
inside the original Hall balance on the stated interval.
\end{proof}

\begin{theorem}[Dense root--inactive descent]
\label{thm:878e9-close}
The exact bound \eqref{eq:target}, and equivalently the Hall inequality \eqref{eq:hall-defect},
holds for every $N\geq8.78\cdot10^9$.
\end{theorem}

\begin{proof}
Theorem~\ref{thm:938e9-close} leaves precisely the interval of
Lemma~\ref{lem:383-root-inactive-endpoint-block}, which closes all six
allocations directly.
\end{proof}


\begin{lemma}[Free integer split endpoint block]
\label{lem:free-integer-split-endpoint-block}
On
\[
 8.77\cdot10^9\leq N\leq8.78\cdot10^9,
\]
all six Hall allocations close directly.  The dense, worst even, and final
odd envelopes are below $0.039999903$, $0.039197329$, and $0.039956604$.
\end{lemma}

\begin{proof}
The intermediate-prime/root decomposition is valid for every integer split
point $Y$; the earlier substitution $Y=\lfloor N/k\rfloor$ was a convenient
one-parameter subfamily, not an arithmetic restriction.  Use cutoff $379$.
At the upper pattern endpoint there are $129618$ forced root-$256$ inactive
patterns.  Applying $4Cq^2\leq834$ and exact combined
density-plus-endpoint ordering leaves $6954$ candidates and gives
\[
 I_{256,256}<0.111791625,
 \qquad \epsilon_{256,256}\leq104529.
\]
The maximizing inactive sets are the same two sets displayed in
Lemma~\ref{lem:383-root-inactive-endpoint-block}; again $C=19$ and $q=3$.
Moreover
\[
 \sum_{p>379}\frac1{p^2}<0.000377981718171.
\]

Choose the fixed integer threshold
\[
 Y=89823400,
 \qquad \pi(Y)=5207279.
\]
The prime--root payment is convex as a function of $N$ on this block, so its
maximum occurs at an endpoint; exact evaluation gives
\[
 \mathfrak P_{379,512,Y}<0.003488984652.
\]
The restricted-diagonal block is maximal at
\[
 N=8770002913,\qquad R=10206521,\qquad \mathcal F(R)=2109663,
\]
and is below $0.002324486044$.  Substitution gives the dense envelope
\[
 0.039999902986\ldots
 <\frac1{25}-\frac7{25N}.
\]
At the new lower endpoint the QR degree remains more than $6238718$ above
half a base class and the mixed strip exceeds $0.001422N$.  The exact
higher-valuation even root-$4096$ block is below $0.039197329$; the three
cutoff-$359$ final-odd root alternatives are below
\[
 0.039303347,\qquad0.039649632,\qquad0.039956604.
\]
All remaining allocations retain larger slack, proving the claim inside the
original Hall balance.
\end{proof}

\begin{theorem}[Free-split dense descent]
\label{thm:877e9-close}
The exact bound \eqref{eq:target}, and equivalently the Hall inequality \eqref{eq:hall-defect},
holds for every $N\geq8.77\cdot10^9$.
\end{theorem}

\begin{proof}
Theorem~\ref{thm:878e9-close} leaves precisely the short block of
Lemma~\ref{lem:free-integer-split-endpoint-block}, which closes every
allocation directly.
\end{proof}

\begin{lemma}[Joint-multiplier and one-coincidence block]
\label{lem:joint-multiplier-one-coincidence-block}
On
\[
 7.8\cdot10^9\leq N\leq8.77\cdot10^9,
\]
all six Hall allocations close directly.  The controlling dense high--low,
worst even, and final odd envelopes are below $0.039989608$,
$0.039667100$, and $0.039176885$.
\end{lemma}

\begin{proof}
Retain cutoff $379$ and recompute the root-$256$ patterns at the upper
endpoint.  A seven-event pattern has actual divisor
\[
 M_D=2\prod_{p\in D}p.
\]
For every such pattern enumerate all multiples $mM_D\leq8770000000$.
Reactivating the other form first leaves $32775$ one-form candidates.  Of
these, $70$ are six-event relaxations with one omitted large prime and are
compared against every candidate.  The exact seven-event patterns give
$53150$ actual multiple entries; sorting them leaves only $865$ support
pairs represented by two values within $834$.  Applying
$4Cq^2\leq834$ to every retained pair gives the safe coordinatewise bound
\[
 I_{256,256}<0.020715563,\qquad
 \epsilon_{256,256}\leq109217.
\]
Thus the dense high--high bound is strictly below the target.
The full-active and one-root-$256$ bounds at the same cutoff are
\[
 (I_{128,128},\epsilon_{128,128})<(0.113387532,140307),
\]
\[
 (I_{256,128},\epsilon_{256,128})<(0.112523483,124883).
\]

Exact restricted-diagonal blocks give the root-$256$ low--low envelope below
$0.039658794$.  For the high--low root sum $384$, use the five split blocks
\[
\begin{array}{c|c|c}
\text{interval for }N&k&\text{envelope}\\ \hline
{[8.50,8.77]\cdot10^9}&108&0.039817900\\
{[8.30,8.50]\cdot10^9}&107&0.039864973\\
{[8.10,8.30]\cdot10^9}&106&0.039913645\\
{[7.95,8.10]\cdot10^9}&105&0.039951123\\
{[7.80,7.95]\cdot10^9}&104&0.039989608.
\end{array}
\]
The joint-multiplier high--high envelope is below $0.034$ throughout.

At the lower endpoint the QR degree is more than $4113070$ above half a base
class and the mixed strip exceeds $0.001054N$.  In the sparse-even
allocation the final odd class therefore contains more than $0.001039N$
pivots and supplies a pair below $963$.  Two coincident active primes would
force
\[
 4(3\cdot7)^2=1764
\]
to divide its difference, which is impossible.  Hence this branch has
$r\leq1$.  The cutoff-$359$ root-$128$/root-$256$
alternatives are respectively below
\[
 0.038513847,\qquad0.038856177,\qquad0.039176885.
\]
The actual higher-valuation even root-$4096$ block is below $0.039667100$,
and all remaining allocations have larger slack.  Every branch therefore
closes in the original Hall balance.
\end{proof}

\begin{theorem}[Joint-multiplier descent]
\label{thm:78e9-close}
The exact bound \eqref{eq:target}, and equivalently the Hall inequality \eqref{eq:hall-defect},
holds for every $N\geq7.8\cdot10^9$.
\end{theorem}

\begin{proof}
Theorem~\ref{thm:877e9-close} leaves precisely the interval of
Lemma~\ref{lem:joint-multiplier-one-coincidence-block}, which closes every
allocation directly.
\end{proof}


\begin{lemma}[Exact finite high--low block]
\label{lem:exact-finite-high-low-block}
On
\[
 7.74855\cdot10^9\leq N\leq7.8\cdot10^9,
\]
all six Hall allocations close directly.  The controlling dense high--low,
worst even, and final-odd envelopes are below $0.039999999884$,
$0.039694959$, and $0.039191886$.
\end{lemma}

\begin{proof}
Use cutoff $353$ and retain the actual floor terms in the degree-$3/2$
Bonferroni count.  For $t\in\{7,18\}$ write
\[
 X_t(N)=\left\lfloor\frac{N-t}{25}\right\rfloor+1.
\]
If $p$ is active for a pivot $b$, the event $p^2\mid (t+25n)b+1$ is the
single class
\[
 n\equiv r_{t,b,p}:=-(tb+1)(25b)^{-1}\pmod {p^2}.
\]
Let $C_X(r,m)$ denote the number of integers $0\leq n<X$ in the class
$r\pmod m$.  Every finite intersection of events is therefore counted
exactly by CRT and the floor
\[
 C_X(r,m)=\begin{cases}
 0,&r\geq X,\\
 1+\lfloor(X-1-r)/m\rfloor,&r<X.
 \end{cases}
\]

For one pivot let $L_b=X-S_1(b)+S_2(b)-S_3(b)$ be the odd Bonferroni lower
bound for its survivor set, with every $S_j$ evaluated by the preceding
CRT floor.  For the union of the two event systems, merging coincident
classes at the same prime, let $U_{b,c}=X-\Sigma_1(b,c)+\Sigma_2(b,c)$ be
the even Bonferroni upper bound for the joint survivor set.  Hence the number
of common bad values in this progression is at most
\begin{equation}
 X-L_b-L_c+U_{b,c}.                                      \label{eq:exact-common-bad}
\end{equation}
This is an integer identity in signed CRT floors, so its maximum on a finite
$N$-interval occurs at the left endpoint or at one of the explicitly
enumerated floor jumps.

At the upper support endpoint there are $98208$ feasible six/seven-event
supports.  Ranking the uniform one-coincidence bounds gives first
\[
 D_1=\{13,17,19,23,29,31,43\},\qquad
 b_1=2\prod_{p\in D_1}p=7466754178,
\]
then
\[
 D_2=\{13,17,19,23,29,31,41\},\qquad
 b_2=2\prod_{p\in D_2}p=7119463286.
\]
Every other support is bounded by the third uniform value
\[
 (I,\epsilon)<(0.112497017,103036),
\]
and hence contributes at most
\begin{equation}
 \frac{2I}{25}+\frac{2\epsilon}{7.74855\cdot10^9}
 <0.009026356261.                                      \label{eq:third-support-pair}
\end{equation}
Both displayed bases exceed half the upper endpoint, so they are the actual
high pivots.  Since there is at most one active coincidence, its prime is one
of $3,7,11,13$ and $4q^2\mid b_i-c$ with $0<|b_i-c|\leq834$.
There are $56$ admissible offsets for each base.  Expanding
\eqref{eq:exact-common-bad} and scanning
every jump on the whole interval gives
\[
 \max\frac1N\sum_{t=7,18}(X_t-L_{b_1}-L_c+U_{b_1,c})
 <0.009000586,
\]
\[
 \max\frac1N\sum_{t=7,18}(X_t-L_{b_2}-L_c+U_{b_2,c})
 <0.009000350.
\]
The witnesses occur at $N=7773532793$, $c=b_1+108$, and at
$N=7781651968$, $c=b_2+288$, respectively.  Thus
\eqref{eq:third-support-pair} is the valid uniform pair payment.  The complete support
ranking, CRT expansion, and jump scan are reproduced by
\path{scripts/exact_dense_high_low_bonferroni.py}.

Now
\[
 \sum_{p>353}p^{-2}<0.000407314700.
\]
On $7.75\cdot10^9\leq N\leq7.8\cdot10^9$, choose the fixed large-prime split
$Y=74680000$, for which $\pi(Y)=4376612$.  The exact restricted-diagonal
scan with root sum $384$ gives
\[
 0.039999616073\ldots
 <\frac1{25}-\frac7{25N}.
\]
On the bottom block choose instead $Y=74751552$, for which
$\pi(Y)=4380561$.  On every interval with fixed restricted-diagonal parameter
$R$, the full envelope minus the Hall target is a sum of an affine function
and positive multiples of $1/N$, and is therefore convex.  Scanning both
endpoints of every such interval on
$7.74855\cdot10^9\leq N\leq7.75\cdot10^9$ gives its maximum at the lower
endpoint:
\[
 0.039999999883517563
 <0.039999999963864207.
\]
The full-active low--low block is below $0.039676122$; the retained actual
multiplier enumeration keeps the high--high block below $0.034$, and every
root-$512$ alternative remains below $0.033652$.

At the lower endpoint
\[
 \mathcal M_{\rm QR2}(N)>\frac N{50}+4003743,
 \qquad \mathcal R(N)>0.001033N.
\]
The higher-root direct progression lower bound is in fact greater than
$N/50+10258173$, so the hierarchy itself retains strict slack.  The exact
higher-valuation even block with actual two-pivot root sum $4096$ is below
$0.039694959$.  If the even mass is at most $1.5\cdot10^{-5}N$, the final
odd class has more than $0.001018N$ pivots and supplies a pair at distance
below $982$.  Two coincidences would force
$4(3\cdot7)^2=1764$ to divide that nonzero difference, so again $r\leq1$.
The three cutoff-$359$ root alternatives are below
\[
 0.038527089,\qquad0.038870308,\qquad0.039191886.
\]
All six allocations therefore contradict the Hall inequality.
\end{proof}

\begin{theorem}[Exact finite-pair descent]
\label{thm:774855e9-close}
The exact bound \eqref{eq:target}, and equivalently the Hall inequality \eqref{eq:hall-defect},
holds for every $N\geq7.74855\cdot10^9$.
\end{theorem}

\begin{proof}
Theorem~\ref{thm:78e9-close} leaves precisely the interval of
Lemma~\ref{lem:exact-finite-high-low-block}, which closes every allocation
inside the original Hall cut.
\end{proof}


\begin{lemma}[Twenty-exception exact finite block]
\label{lem:twenty-exception-finite-block}
On
\[
 7.69835\cdot10^9\leq N\leq7.74855\cdot10^9,
\]
all six Hall allocations close directly.  The controlling dense high--low,
worst even, and final-odd envelopes are below $0.039999992524$,
$0.039722475$, and $0.039206719$.
\end{lemma}

\begin{proof}
Keep cutoff $353$ and the exact CRT/floor expression
\eqref{eq:exact-common-bad}.  At the upper support endpoint there are
$97283$ feasible six/seven-event supports.  Exhaustive ranking shows that
only the following twenty actual bases exceed the twenty-first uniform
record:
\[
\begin{array}{rrrr}
7466754178&7119463286&6424881502&7685998826\\
7540865882&7190127934&6905745814&6318022766\\
6024161242&5436438194&7410796822&7466049734\\
7705963694&7597780762&7601061182&6830641246\\
7325070038&6737654638&6984369106&6954162358.
\end{array}
\]
Every base exceeds half the support endpoint and hence has no nontrivial
multiplier.  For each base enumerate its at most $56$ offsets satisfying
$0<|b-c|\leq834$ and the exact one-coincidence divisibility.  Scanning every
signed-floor jump on the full interval (starting the base $7705963694$ only
at its first feasible value) gives
\[
 \max_{b\text{ in the displayed list}}
 \frac1N\sum_{t=7,18}(X_t-L_b-L_c+U_{b,c})
 <0.009000587.
\]
The maximum is attained for $b=7466754178$, $c=b+108$, at
$N=7703499518$.

Every remaining support is bounded by the twenty-first record
\[
 D_{21}=\{11,13,19,23,29,37,53\},\qquad
 (I_{21},\epsilon_{21})<(0.112331560,103036).
\]
Consequently the true pair payment on this interval is
\[
 \frac{2I_{21}}{25}+\frac{2\epsilon_{21}}{7698350000}
 <0.009013292578,
\]
strictly larger than every exact exception and therefore valid uniformly.
These support rankings and interval jump scans are reproduced by
\path{scripts/exact_dense_high_low_bonferroni.py}.

Choose the integer split
\[
 Y=74567262,\qquad \pi(Y)=4370376.
\]
As in the preceding lemma, total envelope minus target is convex on every
fixed-$R$ interval.  A joint endpoint scan of the whole block, rather than a
sum of unrelated component maxima, has its maximum at the lower endpoint and
gives
\[
 0.039999992523402584
 <0.039999999963628570.
\]
The full-active low--low block is below $0.039688991$; actual multiplier
compatibility keeps the high--high block below $0.034$, and every root-$512$
alternative remains below $0.033652$.

At the lower endpoint
\[
 \mathcal M_{\rm QR2}(N)>\frac N{50}+3897425,
 \qquad \mathcal R(N)>0.001012N,
\]
while the higher-root direct bound exceeds $N/50+9990592$.  The exact
higher-valuation even root-$4096$ block is below $0.039722475$.  In the
sparse-even case the final odd class has more than $0.000997N$ pivots and
therefore supplies a pair at distance below $1003$.  Since $1003<1764$, two
coincidences are again impossible.  The three one-coincidence root
alternatives are below
\[
 0.038540147,\qquad0.038884298,\qquad0.039206719.
\]
Thus every allocation closes in the unchanged Hall balance.
\end{proof}

\begin{theorem}[Twenty-exception finite descent]
\label{thm:769835e9-close}
The exact bound \eqref{eq:target}, and equivalently the Hall inequality \eqref{eq:hall-defect},
holds for every $N\geq7.69835\cdot10^9$.
\end{theorem}

\begin{proof}
Theorem~\ref{thm:774855e9-close} leaves precisely the interval of
Lemma~\ref{lem:twenty-exception-finite-block}, which closes all six
allocations.
\end{proof}


\begin{lemma}[Fifty-exception exact finite block]
\label{lem:fifty-exception-finite-block}
On
\[
 7.68768\cdot10^9\leq N\leq7.69835\cdot10^9,
\]
all six Hall allocations close directly.  The controlling dense high--low,
worst even, and final-odd envelopes are below $0.039999997172$,
$0.039728326$, and $0.039209886$.
\end{lemma}

\begin{proof}
Keep cutoff $353$ and the exact CRT/floor expression
\eqref{eq:exact-common-bad}.  At the upper support endpoint there are
$96431$ feasible supports.  The leading fifty bases are obtained from the
twenty-base list in Lemma~\ref{lem:twenty-exception-finite-block} by deleting
the now-infeasible base $7705963694$ and adjoining
\[
\begin{array}{rrrr}
7107601358&6630712946&6164237222&6389954714\\
5877528514&6302967242&7528040806&5766544498\\
6853887898&7294521806&6629170262&5498333126\\
5955017354&7532934838&7006010726&5280864446\\
7567600898&7390692166&6680149762&6798014366\\
7046939042&4831429174&6111626378&6554010034\\
7319347178&4606711538&7079368582&6028427834\\
6249172358&6359432794&5515370146.&
\end{array}
\]
Denote the resulting fifty-element set by $\mathcal E_{50}$.
Every one exceeds half the support endpoint.  Exhausting all at most $56$
one-coincidence offsets and every signed-floor jump for each of the fifty
actual bases on the whole interval gives
\[
 \max_{b\in\mathcal E_{50}}
 \frac1N\sum_{t=7,18}(X_t-L_b-L_c+U_{b,c})
 <0.009000587.
\]
The maximum again uses $b=7466754178$, $c=b+108$, now at
$N=7694521343$.

Every remaining support is bounded by the fifty-first record
\[
 D_{51}=\{11,13,17,23,29,41,43\},\qquad
 (I_{51},\epsilon_{51})<(0.112296072,103036).
\]
Thus the valid uniform pair payment is
\[
 \frac{2I_{51}}{25}+\frac{2\epsilon_{51}}{7687680000}
 <0.009010491293,
\]
which strictly dominates every exact exception.  The exhaustive ranking and
jump scans are reproduced by
\path{scripts/exact_dense_high_low_bonferroni.py}.

The integer split remains
\[
 Y=74567262,\qquad\pi(Y)=4370376.
\]
The joint convex endpoint scan on the whole block gives
\[
 0.039999997171827493
 <0.039999999963578089.
\]
The full-active low--low block is below $0.039691740$; actual multiplier
compatibility keeps the high--high block below $0.034$, and root-$512$
alternatives remain below $0.033652$.

At the lower endpoint
\[
 \mathcal M_{\rm QR2}(N)>\frac N{50}+3874872,
 \qquad \mathcal R(N)>0.001008N,
\]
and the higher-root direct bound exceeds $N/50+9933718$.  The exact
higher-valuation even root-$4096$ block is below $0.039728326$.  In the
sparse-even case the final odd class has more than $0.000993N$ pivots and
supplies a pair below $1007<1764$, so two coincidences remain impossible.
The three one-coincidence alternatives are below
\[
 0.038542937,\qquad0.038887268,\qquad0.039209886.
\]
Every allocation therefore closes in the unchanged Hall balance.
\end{proof}

\begin{theorem}[Fifty-exception finite descent]
\label{thm:768768e9-close}
The exact bound \eqref{eq:target}, and equivalently the Hall inequality \eqref{eq:hall-defect},
holds for every $N\geq7.68768\cdot10^9$.
\end{theorem}

\begin{proof}
Theorem~\ref{thm:769835e9-close} leaves precisely the interval of
Lemma~\ref{lem:fifty-exception-finite-block}, which closes all six
allocations.
\end{proof}


\begin{lemma}[Ninety-five-exception exact finite block]
\label{lem:ninetyfive-exception-finite-block}
On
\[
 7.67769\cdot10^9\leq N\leq7.68768\cdot10^9,
\]
all six Hall allocations close directly.  The controlling dense high--low,
worst even, and final-odd envelopes are below $0.039999999924$,
$0.039733922$, and $0.039212887$.
\end{lemma}

\begin{proof}
Use cutoff $353$ and the exact signed-floor
expansion \eqref{eq:exact-common-bad}.  At the upper support endpoint there
are $96249$ feasible six/seven-event supports.  Adjoin to the fifty bases in
Lemma~\ref{lem:fifty-exception-finite-block} the following next forty-five:
\[
\begin{array}{rrrr}
5717327902&4157276266&5258841302&5639497006\\
7338804902&7654533458&7137741754&6735615458\\
5159539814&6132426014&7427283578&5931362866\\
4919561218&6907749706&5328173422&6586459022\\
7133521538&7511532314&4724983978&7265252566\\
6025909318&7162158718&7680584054&6927333842\\
7428761626&6526413322&4322857682&7522987186\\
6222859214&7099156922&6673294342&7026917326\\
4121794534&7359108614&6796526594&5787574078\\
6463411526&6700083962&6251496394&6480409078\\
6105354398&5518384586&7235876362&7037633174\\
5615750998.&&&
\end{array}
\]
Call the resulting set $\mathcal E_{95}$.  Every displayed base, and every
base already in $\mathcal E_{50}$, exceeds half the upper support endpoint.
Thus all ninety-five are actual pivots with no multiplier ambiguity.  For a
base or its paired offset that first appears inside the interval, start its
floor scan at $\max\{7.67769\cdot10^9,b,c\}$; this is the exact first
feasible endpoint, not an extrapolation to the left.

Exhausting every admissible one-coincidence offset and every signed-floor
jump on the whole interval gives
\[
 \max_{b\in\mathcal E_{95}}
 \frac1N\sum_{t=7,18}(X_t-L_b-L_c+U_{b,c})
 <0.009000589.
\]
The maximum is attained for $b=7466754178$, $c=b+108$, at
$N=7683469568$, where the exact normalized count is
$0.009000588001027402$.

The ninety-sixth record is
\[
 D_{96}=\{11,13,17,23,29,31,37\},\qquad
 (I_{96},\epsilon_{96})<(0.112262740,103036).
\]
Its minimal base is $3719668238$, below half the support endpoint, so it does
not satisfy the actual-base condition and remains the uniform controller.
On the present interval its payment is
\[
 \frac{2I_{96}}{25}+\frac{2\epsilon_{96}}{7677690000}
 <0.009007859565,
\]
strictly above every exact exception.  The exhaustive ranking and the
first-feasible signed-floor scans are reproduced by
\path{scripts/exact_dense_high_low_bonferroni.py}.

Choose the fixed integer split
\[
 Y=74315802,\qquad \pi(Y)=4356538.
\]
Using the rounded upper bound $I_{96}<0.112262740$, the joint convex endpoint
scan on the entire block has its maximum at the lower endpoint and gives
\[
 0.039999999923493333
 <0.039999999963530697.
\]
The full-active low--low block is below $0.039694322$; actual multiplier
compatibility keeps high--high below $0.034$, and every root-$512$
alternative remains below $0.033652$.

At the lower endpoint,
\[
 \mathcal M_{\rm QR2}(N)>\frac N{50}+3853771,
 \qquad \mathcal R(N)>0.0010038N,
\]
while the higher-root direct progression bound exceeds $N/50+9880469$.
The exact higher-valuation even root-$4096$ block is below $0.039733922$.
After charging an even mass of at most $1.5\cdot10^{-5}N$, the final odd
class retains more than $0.0009888N$ pivots and hence has a pair at distance
below $1012<1764$.  Thus two finite-prime coincidences are impossible.  The
three one-coincidence root alternatives are below
\[
 0.038545557,\qquad0.038890115,\qquad0.039212887.
\]
Every allocation therefore closes directly in the unchanged Hall balance.
\end{proof}

\begin{theorem}[Ninety-five-exception finite descent]
\label{thm:767769e9-close}
The exact bound \eqref{eq:target}, and equivalently the Hall inequality \eqref{eq:hall-defect},
holds for every $N\geq7.67769\cdot10^9$.
\end{theorem}

\begin{proof}
Theorem~\ref{thm:768768e9-close} leaves precisely the interval of
Lemma~\ref{lem:ninetyfive-exception-finite-block}, which closes all six
allocations.
\end{proof}




\begin{theorem}[Erd\H{o}s Problem 848: terminal assembly]
\label{thm:erdos-848-complete}
For every integer $N\geq1$, the maximum cardinality of a set
$A\subseteq[1,N]$ for which $ab+1$ is nonsquarefree for every $a,b\in A$ is
\[
 \left|\{n\leq N:n\equiv7\pmod {25}\}\right|.
\]
Equivalently, the original Hall inequality in the Hall inequality \eqref{eq:hall-defect} holds
for every $N$.
\end{theorem}

\begin{proof}
Theorem~\ref{thm:finite-prefix-colouring-close} gives the upper bound for
$N\leq1.5\cdot10^6$.
Theorems~\ref{thm:first-low-parity-close} and
\ref{thm:second-low-parity-close} give it on
$[1.5\cdot10^6,3\cdot10^6)$ and
$[3\cdot10^6,5\cdot10^6)$, respectively.
Theorem~\ref{thm:mixed-cell-five-million-close} gives it directly for every
$N\geq5\cdot10^6$.  These paper ranges cover all positive integers.
Proposition~\ref{prop:sharpness} supplies the matching lower bound, and
Theorem~\ref{thm:hall-equivalence} gives the equivalent Hall formulation.
\end{proof}

\begin{thebibliography}{9}

\bibitem{Sawhney848}
M.~Sawhney,
\emph{On $A\subseteq[N]$ such that $ab+1$ is never squarefree for
$a,b\in A$}, 2025.
\url{https://www.math.columbia.edu/~msawhney/Problem_848.pdf}

\bibitem{Sothanaphan848}
N.~Sothanaphan,
\emph{An Explicit Threshold in Erd\H{o}s Problem 848}, March 24, 2026.
\url{https://drive.google.com/file/d/1ujhm4_WYpgRV_rd1rJXIfHyvx16COEKe/view}

\bibitem{RosserSchoenfeld}
J.~B. Rosser and L.~Schoenfeld,
\emph{Approximate formulas for some functions of prime numbers},
Illinois J. Math. 6 (1962), 64--94.
\url{https://doi.org/10.1215/ijm/1255631807}

\bibitem{DusartPrimes}
P.~Dusart,
\emph{Estimates of some functions over primes without R.H.},
arXiv:1002.0442, 2010.
\url{https://arxiv.org/abs/1002.0442}

\end{thebibliography}

\end{document}
